\documentclass[11pt,a4paper]{article}
\usepackage[margin=2.4cm]{geometry}
\usepackage{amsmath,amssymb,amsthm,mathtools}
\usepackage[protrusion=true,expansion=false]{microtype}
\usepackage{booktabs,enumitem}
\usepackage{placeins}
\usepackage{tikz}
\usetikzlibrary{arrows.meta,positioning,calc,fit,backgrounds}
\definecolor{tierspend}{RGB}{200,60,60}
\definecolor{tierfull}{RGB}{210,140,40}
\definecolor{tierin}{RGB}{60,130,180}
\definecolor{tierrecv}{RGB}{70,150,90}
\usepackage[colorlinks=true,linkcolor=blue!50!black,citecolor=blue!50!black,urlcolor=blue!50!black]{hyperref}

\emergencystretch=3em
\hbadness=10000
\hfuzz=2pt

\newtheorem{theorem}{Theorem}[section]
\newtheorem{lemma}[theorem]{Lemma}
\newtheorem{proposition}[theorem]{Proposition}
\newtheorem{corollary}[theorem]{Corollary}
\theoremstyle{definition}
\newtheorem{definition}[theorem]{Definition}
\newtheorem{construction}[theorem]{Construction}
\newtheorem{assumption}[theorem]{Assumption}
\newtheorem{remark}[theorem]{Remark}
\newtheorem{example}[theorem]{Example}

\newcommand{\F}{\mathbb{F}}
\newcommand{\Z}{\mathbb{Z}}
\newcommand{\N}{\mathbb{N}}
\newcommand{\G}{\mathbb{G}}
\newcommand{\E}{\mathbb{E}}
\newcommand{\PP}{\mathbb{P}}
\renewcommand{\Pr}{\mathrm{Pr}}
\newcommand{\Fp}{\F_p}
\newcommand{\Fq}{\F_q}
\newcommand{\ip}[2]{\langle #1,\, #2 \rangle}
\newcommand{\Comm}{\mathsf{Commit}}
\newcommand{\Open}{\mathsf{Open}}
\newcommand{\Setup}{\mathsf{Setup}}
\newcommand{\Verify}{\mathsf{Verify}}
\newcommand{\negl}{\mathsf{negl}}
\newcommand{\poly}{\mathsf{poly}}
\newcommand{\PPT}{\mathsf{PPT}}
\newcommand{\Adv}{\mathcal{A}}
\newcommand{\Prov}{\mathcal{P}}
\newcommand{\Ver}{\mathcal{V}}
\newcommand{\Sim}{\mathcal{S}}
\newcommand{\Ext}{\mathcal{E}}
\newcommand{\Rel}{\mathcal{R}}
\newcommand{\Lang}{\mathcal{L}}
\newcommand{\nfsym}{\mathsf{nf}}
\newcommand{\cmsym}{\mathsf{cm}}
\newcommand{\cmx}{\mathsf{cmx}}
\newcommand{\cvsym}{\mathsf{cv}}
\newcommand{\rk}{\mathsf{rk}}
\newcommand{\Pallas}{\ensuremath{\mathsf{Pallas}}}
\newcommand{\Vesta}{\ensuremath{\mathsf{Vesta}}}
\newcommand{\Ep}{\ensuremath{\mathcal{E}}}
\newcommand{\NoteCommit}{\ensuremath{\mathsf{NoteCommit}}}
\newcommand{\Sinsemilla}{\ensuremath{\mathsf{Sinsemilla}}}
\newcommand{\ShortCommit}{\ensuremath{\mathsf{ShortCommit}}}
\newcommand{\CommitIvk}{\ensuremath{\mathsf{Commit}^{\mathsf{ivk}}}}
\newcommand{\Extract}{\ensuremath{\mathsf{Extract}}}
\newcommand{\Acc}{\ensuremath{\mathsf{Acc}}}
\newcommand{\GroupHash}{\ensuremath{\mathsf{GroupHash}}}
\newcommand{\ask}{\ensuremath{\mathsf{ask}}}
\newcommand{\aksym}{\ensuremath{\mathsf{ak}}}
\newcommand{\nksym}{\ensuremath{\mathsf{nk}}}
\newcommand{\ivk}{\ensuremath{\mathsf{ivk}}}
\newcommand{\fvk}{\ensuremath{\mathsf{fvk}}}
\newcommand{\ovk}{\ensuremath{\mathsf{ovk}}}
\newcommand{\dk}{\ensuremath{\mathsf{dk}}}
\newcommand{\pkd}{\ensuremath{\mathsf{pk_d}}}
\newcommand{\gd}{\ensuremath{\mathsf{g_d}}}
\newcommand{\rcv}{\ensuremath{\mathsf{rcv}}}
\newcommand{\rcm}{\ensuremath{\mathsf{rcm}}}
\newcommand{\rivk}{\ensuremath{\mathsf{r_{ivk}}}}
\newcommand{\rsk}{\ensuremath{\mathsf{rsk}}}

\renewenvironment{abstract}{%
  \small\begin{center}{\bfseries\abstractname}\end{center}\medskip\noindent\ignorespaces
}{\par\medskip}

\title{\textbf{\Huge ZSA Guide}\\[6pt]\large Zcash Shielded Assets: issuance, transfer, and the counterfeiting boundary}
\author{m@rek.onl}
\date{}
\begin{document}
\maketitle
\thispagestyle{empty}
\begin{abstract}
Assuming the Orchard protocol of the \emph{Orchard Guide}, this volume of
\emph{The Zcash Arboretum} documents Zcash Shielded Assets: per-asset value
bases and asset identity, transparent issuance and finalization (ZIP-227),
shielded transfer and burn (ZIP-226), and the split-note construction that
guards the counterfeiting boundary. The upstream ZIPs are in flux --- their
transaction carrier was withdrawn and their fee terms removed upstream --- so
every claim is classified as specified, designed-but-unspecified, or open,
and where ZIP text, fork text, and the implementation disagree (as they do on
the split-note nullifier derivation), all three readings are stated with the
divergence flagged rather than silently resolved.
\end{abstract}
\clearpage
\tableofcontents
\clearpage
\section{Scope, status, and sources}\label{sec:zsa-scope}

Zcash Shielded Assets (ZSA) is the proposed extension of Orchard that carries
user-issued assets in the shielded pool. Two consensus ZIPs define it:
ZIP~226, ``Transfer and Burn of Zcash Shielded Assets'', and ZIP~227,
``Issuance of Zcash Shielded Assets'' --- both Status Draft, Category
Consensus, created 2022-05-01, owned by Pablo Kogan and Vivek Arte (QED-it),
Daira-Emma Hopwood, and Jack Grigg, with credits to Daniel Benarroch,
Aurelien Nicolas, Deirdre Connolly, and Teor, and discussed in the single
thread zcash/zips\#618 (\texttt{zip-0226.rst} and \texttt{zip-0227.rst}
lines~1--18 @ \texttt{69610984}). The combined protocol, OrchardZSA, enables
issuance (ZIP~227) and transfer and burn (ZIP~226) of Custom Assets on the
Zcash chain; ZIP~227 must only be implemented in conjunction with ZIP~226, and
issuance is deliberately transparent --- issuance outputs are unencrypted ---
so that every asset's circulating supply can be tracked on chain
(\texttt{zip-0226.rst} lines~42--43; \texttt{zip-0227.rst} lines~47, 54).

Terminology is ZIP~227's: an \emph{Asset} is a type of note transferable on
the Zcash blockchain; ZEC is the default --- and currently the only defined ---
Asset on mainnet, TAZ on testnet; a \emph{Custom Asset} is any Asset other
than ZEC and TAZ (\texttt{zip-0227.rst} lines~35--39).

The design reduces to one alteration of Orchard's value algebra. The fixed
value base $V^{\mathsf{Orchard}}$ of the homomorphic Pedersen value commitment
(\emph{Orchard Guide}, \S``Conservation of value: value commitments and the
binding signature'') is replaced, per Asset, by an \emph{Asset Base} witnessed
inside the circuit; the same Asset Base must appear in the old and the new
note commitments and as the base of the Action's value commitment, so balance
holds separately per Asset Identifier without revealing which Assets --- or
how many distinct ones --- a transaction transfers (\texttt{zip-0226.rst}
lines~59--63). The transparent protocol is unchanged: unshielding is
impossible for Custom Assets, burn is the only supply-reduction mechanism and
sends to no address, and the Native Asset cannot be burnt
(\texttt{zip-0226.rst} lines~206--230).

\subsection{Scope and deferrals}\label{sec:zsa-deferrals}

This volume develops only the delta over Orchard: the per-asset note type and
note commitment; the Asset Base and its derivation from an Asset Identifier;
the per-asset value commitment, the burn mechanism, and the burn-aware balance
equation; split notes, the mechanism that replaces dummy spends for Custom
Assets; the changes to the Action circuit; the issuance layer of ZIP~227 ---
issuance keys, issue notes, reference notes, and the node-side issuance
state --- and the note-plaintext extension.

Everything the delta rests on is constructed in the companion volumes and is
cited, never re-derived: notes and note commitments (\emph{Orchard Guide},
\S``Representation of a note: notes and note commitments''); value
commitments, the bases $V^{\mathsf{Orchard}}$ and $R^{\mathsf{Orchard}}$, and
the binding signature (\emph{Orchard Guide}, \S``Conservation of value: value
commitments and the binding signature''); the Action statement and dummy
spends (\emph{Orchard Guide}, \S``The Action statement, formally'');
nullifiers (\emph{Orchard Guide}, \S``Prevention of double-spending:
nullifiers''); Sinsemilla (\emph{Orchard Guide}, \S``The commitment primitive:
Sinsemilla''); and, at the wallet layer, transaction construction, PCZT, and
fees (\emph{Wallet Guide}, \S``Transaction construction'', \S``Partially
created transactions (PCZT)'', \S``Fees (ZIP-317)'').

Deployment status is stated once, in \S\ref{sec:zsa-status}. Subsequent
sections construct their objects without requalifying each of them with it.

\subsection{Status}\label{sec:zsa-status}

Both ZIPs are Status Draft. ZIP~226's Deployment section reads: ``The Zcash
Shielded Assets protocol is scheduled to be deployed in Network Upgrade~7
(NU7). This ZIP currently assumes that this will be the case''
(\texttt{zip-0226.rst} lines~425--426). ZIP~227's Deployment section is
``TBD'', and its Test Vectors and Reference Implementation entries read ``LINK
TBD'' (\texttt{zip-0227.rst} lines~661--675).

The network-upgrade reality is thinner than the assumption. ZIP~254, the NU7
deployment ZIP, is Withdrawn and delegates to \texttt{draft-arya-deploy-nu7}
(Status Draft), which fixes the consensus branch ID $\mathtt{0x77190AD8}$ and
the minimum network protocol versions 170150 (testnet) and 170160 (mainnet),
leaves the activation height TBD on both networks, and lists no feature ZIPs
at all (\texttt{zip-0254.md}, \texttt{draft-arya-deploy-nu7.md} @
\texttt{69610984}). The repository README's ``NU7 Candidate ZIPs'' list
(218, 230, 231, 233, 234, 235, 2002, 2003, plus a possible ZIP~317 update)
does not include ZIP~226 or ZIP~227, still names the withdrawn ZIP~230, and is
explicitly non-decisional: ``no decision has been made on whether to include
each of these ZIPs''; the deployment draft ``will define which ZIPs are
included in NU7'' (\texttt{README.rst} lines~70--92 @ \texttt{69610984}). In
the Zcash Foundation's February 2026 NU7 sentiment polls, the proposal that
drew ``universal or near-universal support'' was Tachyon, not ZSA
(\emph{Tachyon Guide}, \S``Standing''). ZSA's own standing in those polls is
split: ZCAP supported ZSAs at 70.4\%, the coinholder poll ran 98.6\%
against, and the Foundation did not recommend ZSAs for NU7, writing that
``before any decision is made, we need to understand why coinholders oppose
ZSAs so strongly'' (zfnd.org, ``NU7 Polling Results: What We Heard and Where
We Go From Here'', 2026-02-23,
\url{https://zfnd.org/nu7-polling-results-what-we-heard-and-where-we-go-from-here/}).

The specification's baseline moved beneath it in 2026. ZIP~226 is now defined
relative to the protocol with ZIP~2005 (Ironwood Quantum Recoverability;
Status Proposed, deploying with NU6.3 ``Ironwood'') applied, and states that
ZIP~2005 ``is expected to deploy before any ZSA activation''
(\texttt{zip-0226.rst} line~45, which still cites the ZIP under its older
name ``Orchard Quantum Recoverability''; \texttt{zip-2005.md} header). The dependency points upward: if the
NU6.3 content changes, the ZSA specification inherits the change.

The carrier format has no live specification. ZIP~226 defines its transaction
structure and sighash by deferring to ZIP~230 (the v6 transaction format) and
ZIP~246 (the v6 digests) (\texttt{zip-0226.rst} lines~373--405, 456--462);
both are now Withdrawn. ZIP~230 is retitled ``Withdrawn Version 6 Transaction
Format'' and warns: ``This ZIP has been obsoleted by ZIP 229, and will not be
deployed. Transaction version number 6 is now defined by ZIP 229''
(\texttt{zip-0230.rst} lines~5, 14, 37--40; withdrawal commit
\texttt{04db8968}, 2026-05-16). ZIP~246 is retitled ``Digests for the
Withdrawn Version 6 Transaction Format'' (\texttt{zip-0246.rst} lines~4--11,
36--44; commit \texttt{00f37219}, merged 2026-07-09). The ZIP~229 that now
owns transaction version~6 (Status Draft, created 2026-06-13, for NU6.3)
excludes ZSA by declared non-requirement: ``The v6 transaction format need not
support ZSAs, the zip233Amount field, the explicit fee field, or
per-transparent-input sighash information that appeared in the withdrawn ZIP
230, or those features and other extensibility affordances planned for ZIP
248'', and ``The value 6 for the transaction version number need not avoid
collisions with the withdrawn ZIP 230 or with the current draft of ZIP 248''
(\texttt{zip-0229.md} lines~154--162). ZIP~248, the intended extensible-format
home for the ZIP~230/231/246 material, exists only as the unmerged pull
request zcash/zips\#1156. The rot reaches into the ZSA texts themselves:
ZIP~227's SigHash consensus rule cites ``modifications described in ZIP 226
[\#zip-0226-txiddigest]'', an anchor ZIP~226 no longer contains
(\texttt{zip-0227.rst} lines~474, 693). This volume therefore bins the entire
transaction-format layer --- wire format, digest tree, sighash --- as designed
but unspecified (\S\ref{sec:zsa-method}), even though the cryptographic core
above it is specified.

The v6 note-plaintext lead byte is formally unassigned. Occurrences of
ZIP~230's constant (originally $\mathtt{0x03}$) were changed to the
placeholder \texttt{\{\{LEADBYTE\}\}} ``to clarify that this ZIP does not
reserve that lead byte value, and to avoid confusion with the different
specification of lead byte 0x03 in ZIP 2005''; ZIP~226 correspondingly
requires the placeholder \texttt{\{\{ZSALEADBYTE\}\}} (\texttt{zip-0230.rst}
lines~41--44; \texttt{zip-0226.rst} line~162; commit \texttt{8418662},
2026-05-16). The pinned implementation ships $\mathtt{0x03}$
(\S\ref{sec:zsa-sources}).

The ZSA fee terms lost their upstream home on 2026-06-26: ZIP~317's change
history for that date records ``Removed the ZSA issuance and asset-creation
fee contributions (ZIP 227)'', so ZIP~227's ``Changes to ZIP 317'' section now
points at a fee calculation that contains no ZSA terms (\texttt{zip-0317.rst}
lines~11--15, 668--676; \texttt{zip-0227.rst} lines~604--610). The issuance
and asset-creation contributions survive only in the QED-it zips fork and in
\texttt{librustzcash-zsa} (\S\ref{sec:zsa-sources}); they are binned designed
but unspecified.

Two claims commonly attached to ZSA are stated here at exactly the rigor the
sources support. ZIP~226 lists reference implementations --- the
QED-it forks of zcashd, \texttt{orchard}, \texttt{librustzcash}, and
\texttt{halo2} --- and test vectors at QED-it/zcash-test-vectors
(\texttt{zip-0226.rst} lines~428--439). No artifact in the local corpus
records an audit of ZSA (a repository-wide search over the ZIP texts and all
three QED-it forks finds none), and none documents a testnet deployment; the
only local traces of node integration are the \texttt{temporary-zebra} cargo
feature and the fork ecosystem itself. Both claims therefore rest on web
sources, pinned here by URL and date, and both are scoped. Audit: Least
Authority TFA GmbH, \emph{OrchardZSA Protocol Security Audit Report}
(commissioned by ZCG as Zcash Ecosystem Security Lead), Updated Final Audit
Report 30 January 2025. Scope: the QED-it \texttt{orchard}
\texttt{zsa1}-vs-\texttt{main} diff (PR QED-it/orchard\#7, including the ZSA
circuit extensions) and the QED-it \texttt{halo2} gadget changes (PR
QED-it/halo2\#18), at revisions \texttt{a7c02d22} (initial) and
\texttt{01e85a5f} (verification); findings: no issues, four suggestions (all
but one resolved).
\url{https://leastauthority.com/wp-content/uploads/2025/01/Least-Authority-ZCG-OrchardZSA-Protocol-Updated-Final-Audit-Report.pdf}.
The audit does not cover \texttt{librustzcash-zsa} or the v6 transaction
format layer, and the audited revisions predate the pinned
\texttt{cf801a5d}. Testnet: ZecHub, \emph{Zcash Shielded News Vol.\,61}
(2025-12-17): ``A feature-complete version of ZSAs is now available on
testnet'' (\url{https://zechub.substack.com/p/zcash-shielded-news-vol61});
the testnet is QED-it's public Zebra-based single-node ZSA testnet (see
forum.zcashcommunity.com thread \#53293, ``ZSA Testnet Question''). This
volume accordingly asserts ``audited'' only of the orchard/halo2 circuit
layer at the audited revisions, and ``running on QED-it's public testnet''
of the stack as a whole --- and nothing stronger.

\subsection{Sources}\label{sec:zsa-sources}

Table~\ref{tab:zsa-sources} lists the primary sources, each pinned to a
commit. All facts in this volume were verified firsthand at the stated
commits; the only web sources cited are pinned by URL and date --- the Least
Authority OrchardZSA audit report (2025-01-30), ZecHub \emph{Shielded News}
Vol.\,61 (2025-12-17), and the Zcash Foundation's NU7 Polling Results post
(2026-02-23), all three in \S\ref{sec:zsa-status}.

\begin{table}[htbp]
\centering
\small
\begin{tabular}{@{}lll@{}}
\toprule
Artifact & Pin / date & Role \\
\midrule
\texttt{zcash/zips} (upstream)
  & \texttt{69610984}, 2026-07-09 & primary ZIP ground truth \\
\texttt{QED-it/orchard}, branch \texttt{zsa1}
  & \texttt{cf801a5d}, 2026-06-29 & crate \texttt{orchard} 0.14.0 \\
\texttt{QED-it/librustzcash}, branch \texttt{zsa1}
  & \texttt{3f9b53fc}, 2026-04-17 & v6 carrier, fee rule \\
\texttt{QED-it/zips}, branch \texttt{zsa1}
  & \texttt{fd71419a}, 2026-02-11 & fork baseline for diffing \\
\bottomrule
\end{tabular}
\caption{Primary ZSA sources, pinned.}
\label{tab:zsa-sources}
\end{table}

\paragraph{Upstream versus fork.} Upstream ZIP~226 diverged from the fork text
at \texttt{fd71419a} in three substantive ways: it added the ZIP~2005
dependency paragraph; it changed the split-note $\psi^{\nfsym}$ from
``sampled uniformly at random on $\F_{p_{\Pallas}}$'' to a
$\mathsf{PRFexpand}$ derivation whose domain-separator byte $\mathtt{0x0A}$ is
reserved by ZIP~2005; and it replaced the lead byte $\mathtt{0x03}$ with the
\texttt{\{\{ZSALEADBYTE\}\}} placeholder. ZIP~227 upstream differs from the
fork only in punctuation and one reference title.

\paragraph{Code versus ZIP.} Where the pinned code disagrees with the pinned
ZIP text, this volume presents the ZIP formula and flags the divergence
inline. Two content divergences head the list. First, the split-note nullifier
randomizer: ZIP~226 derives $\psi^{\nfsym}$ with domain byte $\mathtt{0x0A}$,
while \texttt{orchard-zsa} @ \texttt{cf801a5d} computes it with
$\mathsf{PRFexpand}$ domain $\mathtt{0x09}$ over a random per-note seed
(\texttt{src/circuit.rs} lines~316--319) --- the older fork-era derivation.
Second, the note-plaintext lead byte: a placeholder upstream, but
$\mathtt{0x03}$ in code
(\texttt{src/primitives/zcash\_note\_encryption\_domain.rs} lines~46--49) ---
the very value ZIP~2005 now specifies differently. A third divergence is one
of status rather than content: the digest personalisations in
\texttt{orchard-zsa} match the withdrawn ZIP~246 exactly, and
\texttt{librustzcash-zsa} implements \texttt{TxVersion::V6} with the withdrawn
ZIP~230 wire format --- the entire carrier layer the code implements is
specified only by withdrawn ZIPs (\S\ref{sec:zsa-status}). Further
code-versus-text divergences confirmed at the same pins --- the omitted
explicit-fee digest field, the missing memo branch and its non-spec
substitute digest, the empty-action issuance rule, and the
$\mathsf{valueBurn}$ bound wording --- are inventoried in
\S\ref{sec:zsa-impl}, \S\ref{sec:zsa-issuance-state}, and
\S\ref{sec:zsa-burn}.

\paragraph{Pin consistency.} The two implementation forks are not mutually
consistent snapshots: \texttt{librustzcash-zsa} @ \texttt{3f9b53fc} patches
\texttt{orchard} to QED-it rev \texttt{77d3274c}, older than the
\texttt{orchard-zsa} checkout head \texttt{cf801a5d}; it also pins
\texttt{sapling-crypto} @ \texttt{59535fb5} and \texttt{zcash\_spec} @
\texttt{d5e84264}, while \texttt{orchard-zsa} pins
\texttt{halo2\_proofs}/\texttt{halo2\_gadgets} @ \texttt{ef3d0ba2} and the
same \texttt{zcash\_spec} (\texttt{Cargo.toml} lines~234--239 and 118--123
respectively). Each fork is cited at its own pin; where the gap between them
is load-bearing, the text says so.

\paragraph{Gating.} In \texttt{librustzcash-zsa}, all ZSA and v6 code sits
behind \texttt{cfg(zcash\_unstable = "nu7")}; in \texttt{orchard-zsa},
issuance sits behind the cargo feature \texttt{zsa-issuance} (which pulls in
\texttt{secp256k1}) and node integration behind \texttt{temporary-zebra}.

\paragraph{Non-sources.} The branch \texttt{ltf/zsa} of upstream
\texttt{zcash/orchard} points at \texttt{a799082e} (2025-11-25), an ancestor
of \texttt{main}, and contains no ZSA code; it is not usable ground truth.
Node-side enforcement is a second gap: ZIP~227 specifies chained per-node
issuance state (the \texttt{issued\_assets} map), but no node implementation
is checked out locally, so the supply and burn logic verified for this volume
is library-side only (\texttt{orchard-zsa}
\texttt{src/bundle/burn\_validation.rs}, \texttt{src/issuance.rs}).
%% TODO: pin a QED-it node fork (zebra or zcashd) before asserting any
%% node-side enforcement behaviour of the issued_assets state machine.

\subsection{Method: the three bins for ZSA}\label{sec:zsa-method}

This volume adopts the series' classification of design-stage claims
(\emph{Tachyon Guide}, \S``The three-bin method''): every claim is
\emph{specified}, \emph{designed but unspecified}, or an \emph{open problem},
and the bin stays visible in the prose. The distribution differs sharply from
the Tachyon volume's. Bin~(i) holds the entire cryptographic core --- the
per-asset note and its commitment, the Asset Base derivation, value
commitments and the balance equation, split notes, the circuit statement,
issuance keys, issue notes and reference notes, the issuance state machine,
and burn --- because the ZIP texts state it at consensus rigor and the pinned
implementation cross-checks it. Bin~(ii) holds the carrier layer: the v6 wire
format, the digest tree and sighash, the note-plaintext lead byte, and the
ZSA fee terms, each specified today only by withdrawn ZIPs, the QED-it fork
texts, or code (\S\ref{sec:zsa-status}). Bin~(iii) holds NU7 inclusion and
activation, the eventual lead-byte assignment, and the carrier format's
future home. Per the series conventions, every formula in this volume is
transcribed verbatim from the pinned ZIP text or the pinned code, and the
citation names which.

\section{Asset identity}\label{sec:zsa-asset-identity}

%% Sources, verified firsthand:
%%   zcash/zips @ 69610984 (upstream, text of record): zip-0226.rst,
%%     zip-0227.rst, zip-0230.rst, zip-0032.rst, protocol/protocol.tex.
%%   Implementation (QEDit forks, branch zsa1): orchard-zsa @ cf801a5d
%%     (crate `orchard` 0.14.0), librustzcash-zsa @ 3f9b53fc.
%%   NOTE: the zcash/orchard branch origin/ltf/zsa (a799082) is NOT a ZSA
%%     implementation snapshot (upstream main plus a halo2 compatibility
%%     merge; it contains no asset_base.rs). All code citations in this
%%     section deliberately target the QEDit forks.
%% ZIP citations below are file:line at zcash/zips @ 69610984; code
%% citations are file:line at the fork commits above.

An OrchardZSA note denominates value in one of many assets sharing a single
shielded pool, so the protocol must attach to every note an asset identity
satisfying two demands. First, uniqueness: ``The Asset identification should
be unique (among all shielded pools) and different issuer public keys should
not be able to generate the same Asset Identifier''
(\texttt{zip-0227.rst}:76). Second, usability inside the circuit: the
identity must ultimately be a Pallas point, because it serves as the
asset-specific base of the value commitment
(\S\ref{sec:zsa-native-asset}). ZIP~227 meets both demands with one
derivation chain,
\[
(\mathsf{issuer}, \texttt{asset\_desc})
\;\longrightarrow\; \mathsf{assetDescHash}
\;\longrightarrow\; \mathsf{AssetId}
\;\longrightarrow\; \mathsf{AssetDigest}
\;\longrightarrow\; \mathsf{AssetBase},
\]
which the ZIP itself presents as a relation diagram
(\texttt{zip-0227.rst}:210--224): an issuer key and an issuer-chosen
description string are hashed into an identifier, the identifier into a
digest, and the digest onto the Pallas curve. Every Asset issued under
ZIP~227 is scoped to the issuer that issued it, and its Asset Identifier is
globally unique, ``used across all Zcash protocols that support ZSAs''
(\texttt{zip-0227.rst}:227--240).

This section constructs the chain end to end: the issuance key pair at its
root (\S\ref{sec:zsa-issuance-keys}); the hash steps from description to
identifier (\S\ref{sec:zsa-asset-id}); the group hash from identifier to
base point (\S\ref{sec:zsa-asset-base}); the native asset's special position
outside the chain (\S\ref{sec:zsa-native-asset}); and the parts of the chain
consensus sees on the wire versus recomputes (\S\ref{sec:zsa-identity-wire}).
Deployment status for the whole protocol is fixed once, in
\S\ref{sec:zsa-status}, and not restated per object.

\subsection{The issuance key pair}\label{sec:zsa-issuance-keys}

The OrchardZSA protocol adds exactly two keys to the key material of
Orchard: the \emph{issuance authorizing key} $\mathsf{isk}$, which
authorizes the issuance of Asset Identifiers by a given issuer and is used
only by that issuer, and the \emph{issuance validating key} $\mathsf{ik}$,
which validates issuance transactions (\texttt{zip-0227.rst}:84--91).
Neither key touches the spend path. The ZIP's rationale: the issuance key
structure is independent of the original key tree but derived analogously
via ZIP~32, which keeps issuance details and Asset Identifiers consistent
across multiple shielded pools, and separates issuance authority from spend
authority --- allowing a potential transfer of issuance authority without
compromising spend authority (\texttt{zip-0227.rst}:524).

\begin{definition}[Issuance authorization signature scheme]
\label{def:zsa-issueauthsig}
The scheme $\mathsf{IssueAuthSig}$ is instantiated as a BIP-340 Schnorr
signature over the secp256k1 curve, with constants set per the secp256k1
standard parameters as described in BIP~340. Its associated types are:
messages are arbitrary byte sequences; signatures are $64$-byte sequences
or $\bot$; public keys are $32$-byte sequences or $\bot$; private keys are
$32$-byte sequences (\texttt{zip-0227.rst}:122--135). Signing fixes the
BIP-340 auxiliary data to $a = [\mathtt{0x00}]^{32}$: the signature is
$\sigma := \mathsf{Sign}(\mathsf{isk}, M)$ with auxiliary data $a$, or
$\bot$ if $\mathsf{Sign}$ fails; on the wire it is encoded as
$\mathtt{issueAuthSig} = [\mathtt{0x00}] \,\|\, \sigma$
(\texttt{zip-0227.rst}:190--207).
\end{definition}

In the implementation, type \texttt{ZSASchnorr} realizes the scheme with
\texttt{ALGORITHM\_BYTE} $= \mathtt{0x00}$, secret keys as
\texttt{secp256k1::SecretKey}, public keys as \texttt{XOnlyPublicKey}, and
signing via \texttt{sign\_schnorr\_with\_aux\_rand} with the all-zero
$32$-byte auxiliary array (\texttt{orchard-zsa}
\texttt{src/issuance/auth.rs}:163--273 @ \texttt{cf801a5d}).

\begin{construction}[Issuance key derivation]\label{const:zsa-isk-derivation}
The issuance authorizing key is generated by ZIP~32's hardened-only key
derivation --- the maps $\mathsf{MKGh}$ and $\mathsf{CKDh}$, constructed in
the \emph{Wallet Guide}, \S``Hardened-only derivation'', and not re-derived
here --- instantiated in a dedicated $\mathsf{Issuance}$ context
(\texttt{zip-0227.rst}:139--166). Let $S$ be a seed of at least $32$ and at
most $252$ bytes.
\begin{enumerate}[label=(\arabic*)]
\item \emph{Context.} The issuance context sets
  $\mathsf{Issuance.MKGDomain} = \texttt{ZcashSA\_Issue\_V1}$ and
  $\mathsf{Issuance.CKDDomain} = \mathtt{0x81}$
  (\texttt{zip-0227.rst}:148--149).
\item \emph{Master key.} Set
  $m_{\mathsf{Issuance}} := \mathsf{MKGh}^{\mathsf{Issuance}}(S)$.
\item \emph{Child derivation.} Set
  $\mathsf{CKDsk}((\mathsf{sk}_{\mathrm{par}}, c_{\mathrm{par}}), i) :=
  \mathsf{CKDh}^{\mathsf{Issuance}}((\mathsf{sk}_{\mathrm{par}},
  c_{\mathrm{par}}), i)$. The two-argument call fixes the general form's
  $\mathit{lead} = 0$ and $\mathit{tag} = [\,]$, exactly the case whose
  suffix the ZIP-32 encoding omits (\texttt{zip-0032.rst}:449--462; the
  \emph{Wallet Guide}'s remark ``The general form: triples, and sibling
  instantiations'').
\item \emph{Path.} Derive along
  $m_{\mathsf{Issuance}} \,/\, \mathrm{purpose}' \,/\,
  \mathrm{coin\_type}' \,/\, \mathrm{account}'$, with purpose fixed to
  $227$ ($\mathtt{0xe3}$; hardened $\mathtt{0x800000e3}$ per BIP~43),
  coin\_type as in the ZIP~32 key path levels, and account fixed to index
  $0$ (\texttt{zip-0227.rst}:160--164).
\item \emph{Key extraction.} From the generated pair $(\mathsf{sk}, c)$,
  set $\mathsf{isk} := \mathsf{sk}$ (\texttt{zip-0227.rst}:164--166).
\end{enumerate}
\end{construction}

The implementation realizes the context as a struct \texttt{Issuance}
implementing the \texttt{zip32} crate's hardened-only \texttt{Context}
trait, with \texttt{MKG\_DOMAIN} set to \texttt{ZcashSA\_Issue\_V1} and
with \texttt{CKD\_DOMAIN} reusing the Orchard byte, constant
\texttt{PrfExpand::ORCHARD\_ZIP32\_CHILD}; the seed-level constructor
\texttt{from\_zip32\_seed} derives along
$[227', \mathrm{coin\_type}', 0']$ and rejects any non-zero account with
error \texttt{NonZeroAccount} (\texttt{orchard-zsa}
\texttt{src/issuance/auth.rs}:38--99, 219--236 @ \texttt{cf801a5d}). The
\texttt{zip32} crate 0.2.0 implements $\mathsf{MKGh}$ and $\mathsf{CKDh}$
exactly, and its \texttt{derive\_child(index)} calls
\texttt{ckdh\_internal(index, 0, [])} --- $\mathit{lead} = 0$,
$\mathit{tag} = [\,]$, omitted from the PRF message
(\texttt{src/hardened\_only.rs}:95--130).

\begin{remark}[Domain separation between the Orchard and Issuance trees]
\label{rem:zsa-ckd-domain}
The byte $\mathsf{Issuance.CKDDomain} = \mathtt{0x81}$ equals
$\mathsf{Orchard.CKDDomain}$ (\texttt{zip-0032.rst}:476), and the
implementation reuses \texttt{PrfExpand::ORCHARD\_ZIP32\_CHILD}
(\texttt{zcash\_spec} \texttt{src/prf\_expand.rs}:119 @
\texttt{d5e84264}). Separation between the two ZIP-32 trees therefore
rests entirely on the differing $\mathsf{MKGDomain}$: distinct master
secret keys and chain codes feed every subsequent
$\mathsf{PRFexpand}$ call. Neither ZIP states whether the byte reuse is
deliberate; ZIP~32 asks only that $\mathsf{CKDDomain}$ values be tracked as
part of the global set of domains defined for $\mathsf{PRFexpand}$. The
derivation is \emph{specified}; we flag the shared byte because the
separation argument is easy to misattribute to it.
\end{remark}

\begin{definition}[Issuance validating key and issuer identifier]
\label{def:zsa-ik}
The issuance validating key is
$\mathsf{ik} := \mathsf{PubKey}(\mathsf{isk})$, where $\mathsf{PubKey}$ is
the BIP-340 algorithm and the output is in big-endian order as defined in
BIP~340; the derivation returns $\bot$ if $\mathsf{PubKey}$ fails, which
happens with very low probability, in which case the keys are discarded
and derivation repeated with a different $\mathsf{isk}$
(\texttt{zip-0227.rst}:168--183). Its encoding prepends a signature-scheme
byte, which MUST be $\mathtt{0x00}$, indicating BIP~340:
\[
\mathsf{ik\_encoding} = [\mathtt{0x00}] \,\|\, \mathsf{ik}
\qquad (33 \text{ bytes}),
\]
and the issuer identifier is defined by
$\mathsf{issuer} = \mathsf{ik\_encoding}$
(\texttt{zip-0227.rst}:103--110, 178--179). The encoding currently appears
on chain only in the \texttt{issuer} field of an issuance bundle
(\texttt{zip-0227.rst}:178--180).
\end{definition}

\begin{remark}[Permanence of the issuer binding]\label{rem:zsa-rotation}
ZIP~227 notes that ``the equality of \texttt{issuer} and
\texttt{ik\_encoding} might not hold in future when key rotation is
specified for issuance key pairs'' (\texttt{zip-0227.rst}:103--110). No
rotation mechanism exists today: an asset's identity is permanently bound
to the one $\mathsf{ik}$ under which it was issued, and the ZIP's remedy
for key compromise is to finalize the asset and reissue under a new Asset
Identifier: ``If an issuer identifier is compromised, the
$\mathsf{finalize}$ boolean for all the Assets issued with that key should
be set to $1$; then the issuer should change to a new issuance authorizing
key (hence a new issuer identifier), and issue new Assets, each with a new
$\mathsf{AssetId}$'' (\texttt{zip-0227.rst}:640--645, Security and Privacy
Considerations, ``Issuance Key Compromise'').
This is a \emph{stated limitation} of the specified design; the rotation
mechanism itself is an open design area.
\end{remark}

\subsection{From description to identifier}\label{sec:zsa-asset-id}

An issuer names each of its assets with an \emph{asset description}
$\texttt{asset\_desc}$: a non-empty byte sequence, which SHOULD be a
well-formed UTF-8 code unit sequence according to Unicode 15.0.0 or later
(\texttt{zip-0227.rst}:237--241). Within the scope of an issuer, Asset
Identifier uniqueness is obtained by way of the asset description
(\texttt{zip-0227.rst}:237--240). The current ZIP imposes no maximum
length --- its rationale notes that a description ``can be longer than 512
bytes without incurring chain costs'', the 512-byte cap being an earlier
design (\texttt{zip-0227.rst}:528--539). The description is also global:
ZIP~226 states that ``Every Asset is defined by an Asset description,
\texttt{asset\_desc}, which is a global byte string (scoped across all
future versions of Zcash)'', so that in a future upgrade the same string
can carry over to a different shielded pool, with nodes transforming the
identifier chain between pools without balance violations
(\texttt{zip-0226.rst}:93--103).

\begin{definition}[Asset Identifier]\label{def:zsa-asset-id}
The description is hashed,
\[
\mathsf{assetDescHash} :=
\mathrm{BLAKE2b\text{-}256}(\texttt{ZSA-AssetDescCRH},\
\texttt{asset\_desc})
\]
(\texttt{zip-0227.rst}:243--247), and the Asset Identifier is the pair
\[
\mathsf{AssetId} := (\mathsf{issuer}, \mathsf{assetDescHash})
\]
(\texttt{zip-0227.rst}:249--253), with the canonical encoding
\[
\mathsf{EncodeAssetId}(\mathsf{AssetId}) :=
[\mathtt{0x00}] \,\|\, \mathsf{issuer} \,\|\, \mathsf{assetDescHash}
\qquad (66 = 1 + 33 + 32 \text{ bytes}),
\]
where the initial $\mathtt{0x00}$ is a version byte enabling future
issuance protocols (\texttt{zip-0227.rst}:255--261). The version byte is
not to be confused with the byte that follows it --- the first byte of
$\mathsf{issuer}$, which indicates the signature scheme
(Definition~\ref{def:zsa-ik}); every current encoding happens to begin
with two $\mathtt{0x00}$ bytes of distinct meaning.
\end{definition}

The pair structure carries the uniqueness argument. Because the issuer
identifier is a component of $\mathsf{AssetId}$, two different issuers
cannot produce the same identifier; in the ZIP's words, the Custom Asset
``is described via a combination of the issuer identifier and an asset
description string, to preclude the possibility of two different issuers
creating colliding Custom Assets'' (\texttt{zip-0227.rst}:525). Within one
issuer, distinct descriptions separate assets via
$\mathsf{assetDescHash}$; for every new Asset there MUST be a new and
unique Asset Identifier (\texttt{zip-0226.rst}:93--96).

\begin{remark}[Why a hash, not the string]\label{rem:zsa-desc-hash}
ZIP~227 replaces the direct description string of an earlier design with a
collision-resistant hash for five stated reasons
(\texttt{zip-0227.rst}:528--559): a fixed $32$ bytes per Issue Action costs
less bandwidth than up to $512$ bytes; descriptions can exceed $512$ bytes
without chain cost; carrying the byte string on chain would not make the
user-visible description consensus-visible anyway (the string could itself
be a hash of an off-chain description); absent key rotation, marking an
issuer as trusted is insufficient --- each Asset Identifier needs
independent out-of-band verification that conveys the description; and an
on-chain description is an ``attractive nuisance'' --- hashing forces
wallets to obtain the description from a trusted party, a registry, or
direct user approval. Consistent with the last point, wallets MUST NOT
display just the $\texttt{asset\_desc}$ string as the name of the Asset;
the ZIP suggests a petname system mapping names to Asset Identifiers via
client configuration, or the Asset Digest as a compact unique byte
sequence, e.g.\ in QR codes (\texttt{zip-0227.rst}:263--266).
\end{remark}

\paragraph{Implementation.} The identifier is the enum variant
\texttt{AssetId::V0}, holding the validating key and the $32$-byte
description hash; \texttt{encode\_asset\_id} emits the version byte
$\mathtt{0x00}$, the $33$-byte key encoding, and the hash --- $66$ bytes,
matching $\mathsf{EncodeAssetId}$ (\texttt{orchard-zsa}
\texttt{src/note/asset\_base.rs}:20--60 @ \texttt{cf801a5d}). The function
\texttt{compute\_asset\_desc\_hash} enforces non-emptiness through its
argument type \texttt{NonEmpty<u8>} and \emph{panics} if the description
is not well-formed UTF-8 (\texttt{orchard-zsa}
\texttt{src/issuance.rs}:135--154 @ \texttt{cf801a5d}). The specification
says SHOULD; the implementation enforces a hard MUST at its API boundary.
Since only the $32$-byte hash is consensus-visible, neither variant is
consensus-enforceable: we classify the UTF-8 constraint as an issuer-side
convention --- \emph{specified} as a SHOULD, with the stricter behaviour
an implementation choice.

\subsection{From identifier to base point}\label{sec:zsa-asset-base}

Two further steps land the identifier on the Pallas curve. The full chain,
collected:
\begin{align*}
\mathsf{assetDescHash} &:=
  \mathrm{BLAKE2b\text{-}256}(\texttt{ZSA-AssetDescCRH},\
  \texttt{asset\_desc}), \\
\mathsf{AssetId} &:= (\mathsf{issuer}, \mathsf{assetDescHash}), \\
\mathsf{AssetDigest}_{\mathsf{AssetId}} &:=
  \mathrm{BLAKE2b\text{-}512}(\texttt{ZSA-Asset-Digest},\
  \mathsf{EncodeAssetId}(\mathsf{AssetId})), \\
\mathsf{AssetBase}_{\mathsf{AssetId}} &:=
  \mathsf{ZSAValueBase}(\mathsf{AssetDigest}_{\mathsf{AssetId}}),
\end{align*}
where the digest is defined at \texttt{zip-0227.rst}:268--278, the base at
\texttt{zip-0227.rst}:280--290, and the OrchardZSA instantiation of the
final step is
\[
\mathsf{ZSAValueBase}(\mathsf{AssetDigest}) :=
\GroupHash^{\Pallas}(\mathtt{z.cash:OrchardZSA},\ \mathsf{AssetDigest})
\]
(\texttt{zip-0227.rst}:292--301). The subscript $\mathsf{AssetId}$ may be
omitted when clear from context (\texttt{zip-0227.rst}:303--308). The
group hash is the protocol specification's
$\GroupHash^{\PP}$ (\S5.4.9.8, ``Group Hash into Pallas and Vesta''),
whose construction --- \texttt{expand\_message\_xmd} over BLAKE2b, two
field elements through the simplified SWU map, summed and carried through
the isogeny --- the \emph{Orchard Guide} develops in ``Preliminaries:
hashing and encoding on the curve''; we reuse it as a deterministic,
public, nothing-up-my-sleeve map into the Pallas group and do not
re-derive it here.

ZIP~226 types the result as a \emph{non-trivial} point:
$\mathsf{AssetBase}$ is ``the unique element of the Pallas group that
identifies each Asset in the Orchard protocol \dots\ a valid group element
that is not the identity and is not $\bot$''
(\texttt{zip-0226.rst}:111--117), a non-identity point of
$\Ep(\F_{p_{\Pallas}})$ in our notation. Its byte form is
\[
\texttt{asset\_base} :=
\mathrm{LEBS2OSP}_{\ell_{\PP}}(\mathrm{repr}_{\PP}(\mathsf{AssetBase})),
\]
the $32$-byte representation stored in note plaintexts and burn fields
(\texttt{zip-0226.rst}:114; \texttt{zip-0230.rst}:540--545). ZIP~226 notes
the encoding assumes canonicity, ``true for the Pallas group, but may not
hold for future shielded protocols'' (\texttt{zip-0226.rst}:111--117).
Table~\ref{tab:zsa-identity-sizes} collects the sizes along the chain.

\begin{table}[ht]
\centering
\small
\begin{tabular}{@{}llll@{}}
\toprule
Object & Bytes & Form & Source \\
\midrule
$\mathsf{ik}$ & 32 & x-only BIP-340 key, big-endian &
  \texttt{zip-0227.rst}:130--133 \\
$\mathsf{ik\_encoding} = \mathsf{issuer}$ & 33 &
  $[\mathtt{0x00}] \,\|\, \mathsf{ik}$ &
  \texttt{zip-0227.rst}:178--179 \\
$\mathsf{assetDescHash}$ & 32 & BLAKE2b-256 output &
  \texttt{zip-0227.rst}:243--247 \\
$\mathsf{EncodeAssetId}(\mathsf{AssetId})$ & 66 &
  $[\mathtt{0x00}] \,\|\, \mathsf{issuer} \,\|\, \mathsf{assetDescHash}$ &
  \texttt{zip-0227.rst}:255--259 \\
$\mathsf{AssetDigest}$ & 64 & BLAKE2b-512 output &
  \texttt{zip-0227.rst}:268--278 \\
$\texttt{asset\_base}$ & 32 &
  $\mathrm{LEBS2OSP}_{\ell_{\PP}}(\mathrm{repr}_{\PP}(\cdot))$ &
  \texttt{zip-0230.rst}:540--545 \\
\bottomrule
\end{tabular}
\caption{Byte sizes along the asset-identity derivation chain.}
\label{tab:zsa-identity-sizes}
\end{table}

\paragraph{Why the base must be a genuine group-hash output.} The chain's
final step is not a formality. In a transfer, the output note of a split
Action cannot be allowed to carry an arbitrary Pallas point as its base:
``We must enforce it to be an actual output of a GroupHash computation
\dots\ Without this enforcement the prover could input a multiple (or
linear combination) of an existing Asset Base, and thereby attack the
network by overflowing the ZEC value balance and hence counterfeiting ZEC
funds'' (\texttt{zip-0226.rst}:299--307). The circuit therefore ensures
the value base points of all output notes are actual GroupHash outputs
``as they originate in the Issuance protocol which is publicly verified'',
and ZIP~226 adds: ``We prevent a potential malleability attack on the
Asset Identifier by ensuring the output notes receive an Asset Base that
exists on the global state'' (\texttt{zip-0226.rst}:104--106, 299--307).
The enforcement mechanism belongs to the transfer sections of this volume
(\S\ref{sec:zsa-split}).

\begin{remark}[The identity-point corner case]\label{rem:zsa-identity-point}
The protocol specification's $\GroupHash^{G}(D, M)$ has the full group as
codomain, so its output can in principle be the zero point
(\texttt{protocol.tex} \S5.4.9.8; the \emph{Crypto Guide} constructs the
hash-to-curve pipeline). ZIP~227's $\mathsf{ZSAValueBase}$ never
states what happens in that event; the non-identity requirement comes from
ZIP~226's typing of the note field
(\texttt{zip-0226.rst}:111--117) and from the implementation, which panics
inside \texttt{AssetBase::custom} (``this will happen with negligible
probability''), rejects the identity in \texttt{AssetBase::from\_bytes},
and carries an \texttt{Error::AssetBaseCannotBeIdentityPoint} check in
\texttt{IssueAction::verify} that the panic renders unreachable
(\texttt{orchard-zsa} \texttt{src/note/\allowbreak asset\_base.rs}:98--138,
\texttt{src/issuance.rs}:216--244 @ \texttt{cf801a5d}). Whether the
intended consensus behaviour is ``reject the transaction'' or ``the issuer
must pick a different description'' is unspecified in the ZIP text. The
event has negligible probability; the specification gap is nevertheless
real, and we classify the corner case as \emph{open}.
\end{remark}

\subsection{The native asset}\label{sec:zsa-native-asset}

ZEC is the default --- and currently the only defined --- Asset for
mainnet, TAZ for testnet; a \emph{Custom Asset} is any Asset other than
ZEC and TAZ (\texttt{zip-0227.rst}:35--39). The native asset does not pass
through the derivation chain at all. Its base is fixed by definition:
``the Asset Base of the ZEC Asset will be kept as the original value base
point, $V^{\mathsf{Orchard}}$'' (\texttt{zip-0226.rst}:98); for ZEC,
\[
\mathsf{AssetBase}_{\mathsf{AssetId}} := V^{\mathsf{Orchard}},
\]
``so that the value commitment for ZEC notes is computed identically to
the Orchard protocol deployed in NU5. As such
$\mathsf{ValueCommit}^{\mathsf{Orchard}}_{\rcv}(v)$ as defined in ZIP~224
is used as
$\mathsf{ValueCommit}^{\mathsf{OrchardZSA}}_{\rcv}(V^{\mathsf{Orchard}},
v)$'' (\texttt{zip-0226.rst}:182--194). The point
$V^{\mathsf{Orchard}} :=
\GroupHash^{\Pallas}(\mathtt{z.cash:Orchard\text{-}cv}, \texttt{"v"})$ is
already constructed, together with its companion $R^{\mathsf{Orchard}}$,
in the \emph{Orchard Guide}, ``Conservation of value: value commitments
and the binding signature'', Definition ``Net value commitment''; we cite
it and do not re-derive it.

Two consequences follow. First, indistinguishability: ZIP~230 states that
``An Orchard note is essentially equivalent to an OrchardZSA note with
$\mathsf{AssetBase} = V^{\mathsf{Orchard}}$; in particular, they have the
same note commitment \dots\ and need not be distinguished in note
plaintexts'' (\texttt{zip-0230.rst}:531--538), and ZIP~226's privacy
discussion states that OrchardZSA note commitments for the native asset
are indistinguishable from those for non-native Assets
(\texttt{zip-0226.rst}:83--84). Second, exclusion from burning: the first
consensus rule for $\mathtt{assetBurn}$ requires
$\mathsf{AssetBase} \neq V^{\mathsf{Orchard}}$ for every
$(\mathsf{AssetBase}, v) \in \mathtt{assetBurn}$
(\texttt{zip-0226.rst}:220--223) --- the native asset cannot be burnt.

\paragraph{Implementation.} The native base is
\texttt{AssetBase::zatoshi()}, computed as the Pallas hash-to-curve over
personalization \texttt{"z.cash:Orchard-cv"} of the message \texttt{"v"}
--- exactly $V^{\mathsf{Orchard}}$; \texttt{AssetBase::is\_zatoshi()} is a
constant-time equality with that point; and
\texttt{ValueCommitment::derive} computes $[v]\,V + [\rcv]\,R$ with
$V = \texttt{asset.cv\_base()}$ and $R$ the fixed base
$R^{\mathsf{Orchard}}$, so a ZEC note's value commitment is bit-identical
to NU5 Orchard's (\texttt{orchard-zsa} @ \texttt{cf801a5d}:
\texttt{note/asset\_base.rs}:140--155,
\texttt{constants/fixed\_bases.rs}:35--45,
\texttt{value.rs}:380--400).

\subsection{Identity on the wire}\label{sec:zsa-identity-wire}

Of the whole chain, consensus sees exactly two links on the issuance side:
the issuer and the description hash. The v6 transaction's issuance bundle
carries \texttt{issuerLength} (compactSize), \texttt{issuer} (the issuer
identifier per ZIP~227), \texttt{nIssueActions} (compactSize),
\texttt{vIssueActions}, and \texttt{issueAuthSig} --- the signature bytes
preceded by a $\mathtt{0x00}$ byte indicating BIP~340. These fields are
always present; if $\texttt{nIssueActions} = 0$ then \texttt{issuerLength}
MUST be $0$ (\texttt{zip-0230.rst}:166--179, 215--217, 430--443). Each
Issue Action carries \texttt{assetDescHash} as \texttt{byte[32]},
\texttt{nNotes} (compactSize, possibly zero, to finalize without
issuing), \texttt{vNotes} ($115$ bytes per note description: $43$-byte
recipient, $8$-byte value, $32$-byte $\rho$, $32$-byte
$\mathsf{rseed}$), and \texttt{flagsIssuance} with finalize in the least
significant bit; the burn encoding carries \texttt{assetBase} as
\texttt{byte[32]} (\texttt{zip-0230.rst}:367--383, 446--494). On the
transfer side, the v6 Orchard note plaintext inserts the asset between
$\mathsf{rseed}$ and the memo key,
\[
[\mathtt{leadByte}] \,\|\, \mathrm{LEBS2OSP}_{88}(d) \,\|\,
\mathrm{I2LEOSP}_{64}(v) \,\|\, \mathsf{rseed} \,\|\,
\texttt{asset\_base} \,\|\, \mathsf{K}^{\mathsf{memo}}
\]
(\texttt{zip-0230.rst}:523--553).

The consensus rules bind identity to validation: the
$\mathsf{ik\_encoding}$ used to validate \texttt{issueAuthSig} MUST be
taken from the \texttt{issuer} field and MUST start with byte
$\mathtt{0x00}$, and for each Issue Action, the Issue Note's
$\mathsf{AssetBase}$ is derived from the \texttt{assetDescHash} field of
the action and the \texttt{issuer} field of the bundle
(\texttt{zip-0227.rst}:483--491). The implementation makes the
recomputation structural. Function \texttt{read\_bundle} decodes the
issuer via \texttt{IssueValidatingKey::decode}, checking the algorithm
byte; function \texttt{read\_action} reads the $32$-byte hash and
recomputes the base as \texttt{AssetBase::custom} of
\texttt{AssetId::new\_v0(ik, asset\_desc\_hash)}. The Asset Base is never
carried on the wire for issuance, always re-derived
(\texttt{librustzcash-zsa} @ \texttt{3f9b53fc}, crate
\texttt{zcash\_primitives}, file
\texttt{transaction/components/issuance.rs}, lines 22--90). Redundantly
with construction, \texttt{IssueAction::verify} checks that every note's
base equals the one derived from the validating key and the description
hash, failing with error \texttt{IssueBundleIkMismatchAssetBase}
otherwise (\texttt{orchard-zsa} \texttt{src/issuance.rs}:216--244 @
\texttt{cf801a5d}). An asset's
identity is therefore not an on-chain object at all: it is a recomputable
function of two consensus-visible byte strings, anchored to the issuer key
that signs the bundle.

\paragraph{Classification.} The entire derivation chain of this section is
\emph{specified}: every formula above appears in the ZIP texts at rigor
and is realized by the pinned implementation forks at the pinned commits
(\texttt{orchard-zsa} @ \texttt{cf801a5d}, \texttt{librustzcash-zsa} @
\texttt{3f9b53fc}). Both ZIPs are Status Draft, Category Consensus, owned
by Pablo Kogan, Vivek Arte, Daira-Emma Hopwood, and Jack Grigg, created
2022-05-01 and discussed in zcash/zips issue \#618 and PR \#680
(\texttt{zip-0226.rst}:1--18, \texttt{zip-0227.rst}:1--18); deployment
status is stated in the scope section. Three reference gaps qualify the
bin. First, ZIP~227's Test Vectors and Reference Implementation sections
read ``LINK TBD'' and its Deployment section ``TBD''
(\texttt{zip-0227.rst}:661--675); the de-facto assertion-checked vectors
live in the fork: file \texttt{test\_vectors/asset\_base.rs}, generated
from the zcash-hackworks test-vectors script
\texttt{orchard\_zsa\_asset\_base.py}, is replayed by a test that
recomputes the full chain from $(\mathsf{ik\_encoding},
\texttt{asset\_desc})$ to \texttt{asset\_base}
(\texttt{src/note/asset\_base.rs}:236--261 @ \texttt{cf801a5d}), with
companion vectors for the signature scheme and the ZIP-32 derivation in
\texttt{issuance\_auth\_sig.rs} and \texttt{zip32.rs}.
The vectors' fixed \texttt{[u8; 512]} description arrays are a relic of
the dropped 512-byte cap and imply no bound. Second, the corner cases
flagged above remain: the identity-point behaviour is \emph{open}
(Remark~\ref{rem:zsa-identity-point}), the UTF-8 constraint is a SHOULD
with a stricter implementation (\S\ref{sec:zsa-asset-id}), and key
rotation is an open design area (Remark~\ref{rem:zsa-rotation}). Third,
upstream \texttt{zcash/zips} @ \texttt{69610984} is the text of record;
the QEDit zips fork (\texttt{zips-zsa} @ \texttt{fd71419a}) is essentially
identical for the asset-identity material, and the differences ---
upstream's later additions of a ZIP-2005 dependency, a nullifier
$\psi^{\nfsym}$ derivation, and a lead-byte rule --- do not touch this
chain.

\section{Issuance and finalization (ZIP-227)}\label{sec:zsa-issuance}

The preceding section moved Custom Assets between shielded notes; this one
brings them into existence. ZIP-227 (Status Draft, Category Consensus;
owners Pablo Kogan and Vivek Arte of QEDit, Daira-Emma Hopwood, and Jack
Grigg) specifies the issuance mechanism for Custom Assets on the Zcash
chain: a per-transaction \emph{issuance bundle} in which a publicly named
issuer creates new notes, a per-Asset \emph{finalize} switch that ends
issuance forever, and a node-maintained global supply map. The draft's Test
Vectors, Reference Implementation, and Deployment sections all read TBD
(ZIP-227, header and closing sections @ zips \texttt{6961098}). Deployment
standing --- Draft ZIPs, externally audited at circuit scope, running on
QED-it's public testnet, contested for NU7 --- is fixed once in
\S\ref{sec:zsa-scope} and not relitigated here.

Issuance is deliberately \emph{not} shielded. The ZIP's Motivation states:
``This ZIP only enables \emph{transparent} issuance. As a first step,
transparency will allow for proper testing of the applications that will be
most used in the Zcash ecosystem, and will enable the supply of Assets to
be tracked'' (ZIP-227, \emph{Motivation}). The issuer, the Asset, the
issued amounts, and the finalization status are all consensus-visible;
only the recipients of the issued notes are shielded, because the notes
themselves enter the Orchard note commitment tree
(\S\ref{sec:zsa-issuance-notes}).

\paragraph{Encoding ground truth.} One sourcing decision governs every
wire format and digest below, and we state it once. ZIP-227 normatively
cites ZIP-230 for all issuance encodings and ZIP-246 for all issuance
digests, yet upstream both are Withdrawn: ZIP-230 is the ``Withdrawn
Version 6 Transaction Format'', obsoleted by ZIP-229, its lead byte
replaced by the placeholder \texttt{\{\{LEADBYTE\}\}}, and ZIP-246 carries
``Digests for the Withdrawn Version 6 Transaction Format'', with v6
digests now assigned to ZIP-229 (ZIP-230 and ZIP-246, headers and
warnings @ zips \texttt{6961098}) --- and ZIP-229 carries no
issuance-bundle fields. The QEDit fork of the ZIPs repository retains
ZIP-230 and ZIP-246 at Status Draft with lead byte $\mathtt{0x03}$, and
its copy of ZIP-227 differs from upstream only cosmetically (zips-zsa @
\texttt{fd71419}). We therefore present the ZIP-230/246 issuance encoding
as what it is: the fork-specified design, deployed on QED-it's public
testnet (ZecHub Vol.\,61, 2025-12-17; \S\ref{sec:zsa-status}) and
implemented by the pinned fork libraries (the circuit layer audited by
Least Authority, 2025-01-30), with no merged upstream v6 specification
behind it.
The NU7 standing of the whole format stack is part of the deployment
picture of \S\ref{sec:zsa-scope}.

\paragraph{Implementation pins.} Claims about deployed behavior cite the
QEDit libraries at fixed commits: \texttt{orchard-zsa} branch
\texttt{zsa1} @ \texttt{cf801a5d} (2026-06-29), \texttt{librustzcash-zsa}
branch \texttt{zsa1} @ \texttt{3f9b53fc} (2026-04-17), and the patched
\texttt{zcash\_spec} (QED-it fork @ \texttt{d5e8426}) that carries the new
PRF domain byte; upstream \texttt{orchard}'s branch \texttt{ltf/zsa} is
context only, not ground truth (\S\ref{sec:zsa-sources},
\emph{Non-sources}). All \texttt{librustzcash-zsa}
transaction-format code sits behind
\texttt{cfg(zcash\_unstable = "nu7")}; file paths we cite within that
repository are relative to \texttt{zcash\_primitives/src/}. Except
where flagged inline,
everything in this section is \emph{specified}: ZIP text at rigor,
mirrored by these libraries. Three flags recur: fork-only specification
(the wire format above, and fees, \S\ref{sec:zsa-issuance-fees});
library-versus-ZIP divergence, noted where it occurs; and node-side
enforcement of global state, which no local repository exhibits
(\S\ref{sec:zsa-issuance-state}).

\subsection{Issuance keys and the issuer identifier}%
\label{sec:zsa-issuer-keys}

Issuance authority is a new capability, deliberately disjoint from spend
authority. The rationale in the ZIP: the issuance key tree is
``independent of but analogous to'' the spending key tree, so an
organization can delegate issuance without exposing funds, and a future
scheme can replace the issuance leaf with a multisignature without
touching spending keys (ZIP-227, \emph{Rationale}, and \emph{Issuance Key
Derivation}).

The key material itself is constructed once, in
\S\ref{sec:zsa-issuance-keys}: the BIP-340 signature scheme
(Definition~\ref{def:zsa-issueauthsig}), the hardened-only ZIP-32
derivation of $\mathsf{isk}$ in the $\mathsf{Issuance}$ context
(Construction~\ref{const:zsa-isk-derivation}), and the validating key with
the issuer identifier, $\mathsf{ik} := \mathit{PubKey}(\mathsf{isk})$ and
$\mathsf{issuer} = \mathsf{ik\_encoding} = [\mathtt{0x00}] \,\|\,
\mathsf{ik}$ (Definition~\ref{def:zsa-ik}). We restate only the parameters
the issuance layer consumes: seed $S$ of $32$ to $252$ bytes,
$\mathsf{Issuance.MKGDomain} = \texttt{ZcashSA\_Issue\_V1}$,
$\mathsf{Issuance.CKDDomain} = \mathtt{0x81}$, hardened path
$m_{\mathsf{Issuance}} \,/\, 227' \,/\, \mathit{coin\_type}' \,/\, 0'$,
and $\mathsf{isk} := \mathsf{sk}$ at the leaf (ZIP-227, \emph{Issuance Key
Derivation}). No rotation exists today: the on-chain issuer name is the
validating key itself, and the ZIP notes the equality might not hold in
future if key rotation is ever specified
(Remark~\ref{rem:zsa-rotation}; \S\ref{sec:zsa-issuance-final}).

\subsection{Asset identity: description, digest, and Asset Base}%
\label{sec:zsa-issuance-asset}

Every Custom Asset is named by its issuer plus a human-meaningful
description, and consumed by the protocol as a Pallas point. The full
derivation chain --- $\mathsf{assetDescHash}$, $\mathsf{AssetId}$,
$\mathsf{EncodeAssetId}$ with its version byte, $\mathsf{AssetDigest}$,
and $\mathsf{AssetBase} = \mathsf{ZSAValueBase}(\mathsf{AssetDigest})$ ---
is constructed once, in \S\ref{sec:zsa-asset-id} and
\S\ref{sec:zsa-asset-base} (Definition~\ref{def:zsa-asset-id}); the
hash-not-string rationale is Remark~\ref{rem:zsa-desc-hash}, and the UTF-8
SHOULD-versus-panic delta between ZIP and implementation is recorded in
\S\ref{sec:zsa-asset-id}. The native Asset bypasses the chain entirely:
its base is the Orchard value-commitment base $V^{\mathsf{Orchard}}$
(\S\ref{sec:zsa-native-asset}). An Asset Base equal to the Pallas identity
is invalid everywhere --- \texttt{AssetBase::custom} panics after hashing
to the curve and \texttt{AssetBase::from\_bytes} rejects the identity on
parse (Remark~\ref{rem:zsa-identity-point}). What the issuance layer adds
to the chain is consensus-side recomputation: validators re-derive every
Issue Note's $\mathsf{AssetBase}$ from the bundle's \texttt{issuer} and
the action's \texttt{assetDescHash}, never reading it off the wire
(\S\ref{sec:zsa-identity-wire}; consensus rules,
\S\ref{sec:zsa-issuance-state}).

\subsection{Issue notes and the derived $\rho$}%
\label{sec:zsa-issuance-notes}

\begin{definition}[Issue note]\label{def:zsa-issue-note}
An \emph{Issue Note} is a tuple
$(\mathsf{d}, \pkd, \mathsf{v}, \mathsf{AssetBase}, \rho,
\mathsf{rseed})$,
\[
\mathsf{Note}^{\mathsf{Issue}} :=
  \{0,1\}^{\ell_{\mathsf{d}}} \times
  \mathsf{KA}^{\mathsf{Orchard}}.\mathsf{Public} \times
  \{0\,..\,2^{\ell_{\mathsf{value}}}-1\} \times
  \Ep(\F_{p_{\Pallas}})^{\ast} \times
  \F_{p_{\Pallas}} \times
  \{0,1\}^{8 \cdot 32},
\]
a diversifier and diversified transmission key as in Orchard, a value
$\mathsf{v}$ with $\ell_{\mathsf{value}} = 64$, a non-identity Asset Base,
the nullifier-input field element $\rho$, and a $32$-byte
$\mathsf{rseed}$ that MUST be sampled uniformly at random by the issuer
(ZIP-227, \emph{Issue Note}; the ZIP writes the fourth and fifth
components $\mathbb{P}^{\ast}$ and $\F_{q_{\mathbb{P}}}$, our
$\Ep(\F_{p_{\Pallas}})^{\ast}$ and $\F_{p_{\Pallas}}$).
\end{definition}

Issue notes never appear on chain as commitments: the wire format carries
the note fields themselves (\S\ref{sec:zsa-issuance-wire}), and each
validator computes the note commitment --- by the OrchardZSA
$\NoteCommit$ of \S\ref{sec:zsa-transfer} --- and appends it to the
Orchard note commitment tree. Precisely because the nullifier key of the
recipient is unknown to validators, the note's future spend cannot be
linked to the issuance transaction (ZIP-227, \emph{Issue Note} and
\emph{Issuance Protocol}). Issuance is thus transparent at creation and
shielded from the first spend onward.

\paragraph{Computation of $\rho$.} An Orchard note takes
$\rho := \nfsym_{\mathrm{old}}$ from the Action that creates it (the
\emph{Orchard Guide}, \S``Transmission and detection of a note'');
an issue note has no such Action, so ZIP-227 derives $\rho$ from the
transaction's OrchardZSA transfer bundle instead:
\[
\mathsf{DeriveIssuedRho}(\nfsym, i_{\mathsf{A}}, i_{\mathsf{N}}) :=
\mathrm{ToBase}\bigl(\mathsf{PRFexpand}_{\nfsym}\bigl(
\mathtt{0x84} \,\|\, \mathrm{I2LEOSP}_{32}(i_{\mathsf{A}}) \,\|\,
\mathrm{I2LEOSP}_{32}(i_{\mathsf{N}})\bigr)\bigr),
\]
\[
\rho := \mathsf{DeriveIssuedRho}\bigl(\nfsym_{0,0},\
\mathsf{index}_{\mathrm{Action}},\ \mathsf{index}_{\mathrm{Note}}\bigr),
\]
with the PRF keyed by $\nfsym$ as $32$ little-endian bytes
($\mathrm{I2LEOSP}_{256}$),
where $\nfsym_{0,0}$ is the nullifier of the input note in the
\emph{first} Action of the \emph{first} Action Group of the OrchardZSA
bundle, and the two indices are the zero-based positions of the Issue
Action within the issuance bundle and of the note within its action
(ZIP-227, \emph{Computation of $\rho$}). The domain-separator byte
$\mathtt{0x84}$ extends the $\mathsf{PRFexpand}$ lead-byte registry (the
\emph{Wallet Guide}, \S``The $\mathsf{PRFexpand}$ lead-byte registry'');
the patched \texttt{zcash\_spec} carries it as
\texttt{ORCHARD\_DERIVED\_ISSUE\_RHO} (QED-it \texttt{zcash\_spec},
\texttt{src/prf\_expand.rs} line~107 @ \texttt{d5e8426}). The
implementation computes the identical expression --- $\mathsf{PRFexpand}$
keyed by the $32$-byte little-endian nullifier over
$\mathtt{0x84}$ followed by two $4$-byte little-endian indices, reduced by
$\mathrm{ToBase}$ (\texttt{orchard-zsa}, \texttt{src/note.rs}
lines~479--489 @ \texttt{cf801a5d}). This derivation is why an issuance
bundle cannot travel alone: requiring at least one Action Group both
supplies $\nfsym_{0,0}$ and blocks replay --- a transaction containing
only an issuance bundle would have a SIGHASH computed solely from that
bundle, so a duplicated bundle would share the SIGHASH (ZIP-227,
\emph{Rationale}).

In the builder, notes are created with $\rho$ unset, via
\texttt{new\_issue\_note}; once the transfer bundle's first nullifier
is known, \texttt{update\_rho\_for\_issuance\_note} sets $\rho$ and
resamples $\mathsf{rseed}$ in a loop until the note commitment is not
$\bot$\footnote{Strictly, this conditions $\mathsf{rseed}$ on a
negligible-probability event, a technical deviation from the ZIP's
``sampled uniformly at random'' (ZIP-227, \emph{Issue Note}).}
(\texttt{orchard-zsa}, \texttt{src/note.rs} lines~457--489 @
\texttt{cf801a5d}).

\paragraph{Reference notes.} The first time an Asset is issued, the
issuer MUST place a \emph{reference note} as the first note of that
Issue Action: a note of value $\mathsf{v} = 0$ whose recipient
$(\mathsf{d}, \pkd)$ is the default diversified payment address
(diversifier index $0$) derived from the \emph{all-zero} Orchard spending
key (ZIP-227, \emph{Reference Notes}). The implementation hard-codes the
resulting $43$-byte raw address as \texttt{RAW\_REFERENCE\_RECIPIENT},
byte-identical to the array in the ZIP, with
$\texttt{ReferenceKeys::sk()} =
\texttt{SpendingKey::from\_bytes([0;\,32])}$; the predicate
\texttt{is\_reference\_note} requires zero value and exactly this
recipient, \texttt{get\_reference\_note} accepts only the \emph{first}
note of an action, and the builder inserts the reference note itself when
told the issuance is a first issuance --- adding a recipient for an Asset
that already has notes fails with \texttt{CannotBeFirstIssuance}
(\texttt{orchard-zsa}, \texttt{src/constants/reference\_keys.rs}
lines~14--33 and \texttt{src/issuance.rs} lines~251--258, 459--507 @
\texttt{cf801a5d}). The reference note is retained in global state as
$\mathsf{note_{ref}}$ (\S\ref{sec:zsa-issuance-state}).

The ZIP mandates reference notes and stores them in consensus state, but
states no rationale for their existence. The evident function --- that
every issued Asset Base is guaranteed at least one note in the commitment
tree, available as a ZIP-226 split-input source for a first spend --- is
our reading, stated nowhere in the ZIP text or its rationale sections; we
classify the mechanism itself as \emph{specified} and this purpose as
\emph{designed but unspecified}.
%% TODO: if a QEDit design document or forum post states the reference-
%% note purpose, cite it and upgrade the bin; none is among the local
%% ground-truth sources.

\subsection{The issuance bundle and its encoding}%
\label{sec:zsa-issuance-wire}

An \emph{Issuance Bundle} carries the issuer identifier, the list of
Issue Actions, and one authorizing signature: $\mathsf{issuer}$, which
lets validators verify that each AssetId is properly associated with this
issuer; $\mathsf{vIssueActions}$, an array of Issue Actions; and
$\mathsf{issueAuthSig}$, the encoding of a signature over the transaction
SIGHASH by the issuance authorizing key $\mathsf{isk}$ (ZIP-227,
\emph{Issuance Bundle}). An \emph{Issuance Action} groups the notes
issued for one Asset: the $32$-byte $\mathsf{assetDescHash}$, the note
list, and a flags byte carrying $\mathsf{finalize}$ (ZIP-227,
\emph{Issuance Action}). Table~\ref{tab:zsa-issue-wire} gives the v6
serialization; the encoding rigor caveat of the section introduction
applies.

\begin{table}[htbp]
\centering
\small
\begin{tabular}{@{}lll@{}}
\toprule
Field & Type & Notes \\
\midrule
\texttt{issuerLength} & \texttt{compactSize} &
  MUST be $0$ if \texttt{nIssueActions} $= 0$ \\
\texttt{issuer} & \texttt{byte[issuerLength]} &
  $= \mathsf{ik\_encoding}$, $33$ bytes when present \\
\texttt{nIssueActions} & \texttt{compactSize} & \\
\texttt{vIssueActions} & \texttt{IssueAction[nIssueActions]} & \\
\texttt{issueAuthSig} & \texttt{IssueAuthSignature} &
  present iff \texttt{nIssueActions} $> 0$ \\
\midrule
\multicolumn{3}{@{}l@{}}{\emph{Per \texttt{IssueAction}:}} \\
\texttt{assetDescHash} & \texttt{byte[32]} & \\
\texttt{nNotes} & \texttt{compactSize} & may be $0$ (finalize-only) \\
\texttt{vNotes} & \texttt{IssueNoteDescription[nNotes]} &
  $115$ bytes each \\
\texttt{flagsIssuance} & \texttt{byte} &
  $\mathsf{finalize}$ at LSB; other bits MUST be $0$ \\
\midrule
\multicolumn{3}{@{}l@{}}{\emph{Per \texttt{IssueNoteDescription}
  ($43+8+32+32 = 115$ bytes):}} \\
\texttt{recipient} & \texttt{byte[43]} &
  $\mathrm{LEBS2OSP}_{88}(\mathsf{d}) \,\|\, \pkd^{\star}$ \\
\texttt{value} & \texttt{uint64} & little-endian \\
\texttt{rho} & \texttt{byte[32]} & as for Orchard notes \\
\texttt{rseed} & \texttt{byte[32]} & as for Orchard notes \\
\bottomrule
\end{tabular}
\caption{ZSA issuance fields of the v6 transaction format (fork ZIP-230,
\emph{Transaction Format}, \emph{Issuance Action Description}, and
\emph{Issue Note Description} @ zips-zsa \texttt{fd71419}; upstream text
identical apart from the Withdrawn status).}
\label{tab:zsa-issue-wire}
\end{table}

Three encoding facts deserve emphasis. First, the \texttt{recipient}
field is exactly an Orchard raw payment address (the \emph{Wallet Guide},
\S``Raw encodings''), and together with the bundle's $\mathsf{issuer}$
and the action's $\mathsf{assetDescHash}$, the four fields of an
\texttt{IssueNoteDescription} determine the Issue Note completely
(ZIP-230, \emph{Issue Note Description}; ZIP-227, \emph{Issue Note}) ---
the note commitment is deliberately absent, since validators recompute it
(\S\ref{sec:zsa-issuance-notes}). Second, an action with
\texttt{nNotes} $= 0$ is legal: it lets an issuer finalize an Asset
without issuing more of it (ZIP-230, \emph{Issuance Action Description};
finalization semantics in \S\ref{sec:zsa-issuance-final}). Third, the
flags byte is strict: \texttt{IssuanceFlags::from\_byte} returns
\texttt{None} whenever any bit other than the LSB is set, and
deserialization fails on such bytes, enforcing ZIP-230's ``remaining bits
are set to 0'' at parse time (\texttt{orchard-zsa},
\texttt{src/issuance.rs} lines~57--110; \texttt{librustzcash-zsa},
\texttt{transaction/components/issuance.rs} lines~89--91).

The \texttt{librustzcash-zsa} serializer pins down the absent-bundle
corner cases: \texttt{read\_bundle} reads \texttt{issuer} as a
length-prefixed vector and returns \texttt{None} only for an empty issuer
with zero actions, erroring on an empty issuer with non-empty actions and
on a non-empty issuer with zero actions; \texttt{write\_bundle} encodes
an absent bundle as two zero \texttt{compactSize}s. The authorization
serializes as a length-prefixed \texttt{sighashInfo} followed by the
length-prefixed signature encoding, and the parser accepts exactly
$[\mathtt{0x00}]$ as \texttt{sighashInfo}
(\texttt{transaction/components/issuance.rs} lines~20--191 and
\texttt{sighash\_versioning.rs} @ \texttt{3f9b53fc}; the meaning of
$[\mathtt{0x00}]$ is \S\ref{sec:zsa-issuance-auth}).

In the \texttt{orchard-zsa} object model, the type
\texttt{IssueBundle<T: IssueAuth>} holds the validating key
\texttt{ik} (an \texttt{IssueValidatingKey<ZSASchnorr>}), a non-empty
vector of actions, and an authorization typestate \texttt{T}; an
\texttt{IssueAction} holds the $32$-byte \texttt{asset\_desc\_hash},
its \texttt{Vec<Note>}, and its \texttt{IssuanceFlags}
(\texttt{src/issuance.rs} lines~47--124 @ \texttt{cf801a5d}).

\subsection{Authorization: BIP-340 over the transaction sighash}%
\label{sec:zsa-issuance-auth}

The issuance authorization signature scheme $\mathsf{IssueAuthSig}$ is
instantiated as a BIP-340 Schnorr signature over secp256k1 with the
standard parameters --- not RedPallas: issuance is a transparent
capability, and the choice buys direct compatibility with existing
secp256k1 tooling. Batch verification MAY be used, and precomputation MAY
be used if and only if it produces equivalent results (ZIP-227,
\emph{Issuance Authorization Signature Scheme}). The associated types
(same section):
\[
\mathsf{Message} = \mathbb{B}^{\mathbb{Y}^{[\mathbb{N}]}}, \quad
\mathsf{Signature} = \mathbb{B}^{\mathbb{Y}^{[64]}} \cup \{\bot\}, \quad
\mathsf{Public} = \mathbb{B}^{\mathbb{Y}^{[32]}} \cup \{\bot\}, \quad
\mathsf{Private} = \mathbb{B}^{\mathbb{Y}^{[32]}},
\]
that is, arbitrary byte-sequence messages, $64$-byte signatures,
$32$-byte x-only public keys, and $32$-byte secret keys, with $\bot$ as
the failure value. Signing fixes the BIP-340 auxiliary data to the
all-zero block:
\[
a := [\mathtt{0x00}]^{32}, \qquad
\sigma := \mathit{Sign}(\mathsf{isk}, M)\ \text{with auxiliary data}\ a,
\qquad
\mathtt{issueAuthSig} := [\mathtt{0x00}] \,\|\, \sigma,
\]
returning $\bot$ on failure; validation returns $0$ if
$\sigma = \bot$, else $1$ exactly when the BIP-340
$\mathit{Verify}(\mathsf{ik}, M, \sigma)$
succeeds\footnote{The ZIP's Validate definition writes
$\mathit{Verify}(\mathsf{key}, M, \sigma)$ for the parameter it names
$\mathsf{ik}$ --- an editorial slip we flag when citing (ZIP-227,
\emph{Issuance Authorization Signing and Validation}).} (ZIP-227,
\emph{Issuance Authorization Signing and Validation}). With the scheme
byte, the two wire encodings are $33$ bytes
($\mathsf{ik\_encoding}$) and $65$ bytes ($\mathtt{issueAuthSig}$).

\paragraph{What is signed.} The message is the transaction
$\mathsf{SigHash}$ as defined by protocol specification \S4.10 under the
ZIP-226 modifications, with sighash type $\mathsf{SIGHASH\_ALL}$, and the
consensus rule requires
$\mathsf{IssueAuthSig.Validate}(\mathsf{ik}, \mathsf{SigHash},
\mathtt{issueAuthSig}) = 1$ (ZIP-227, \emph{Specification: Consensus Rule
Changes}). Under ZIP-246 the signature digest includes branch S.5, the
$\mathsf{issuance\_digest}$, identical to the txid branch --- so the
issuance signature commits to \emph{all} effecting data of the
transaction, including the issuance bundle's own effecting data, though
never to any signature (ZIP-246, \emph{Signature Digest}). ZIP-246 also
introduces sighash algorithm versioning: each bundle carries
$\mathsf{sighashInfo} = [\mathsf{sighashVersion}] \,\|\,
\mathsf{associatedData}$, versions are numbered from $0$ per transaction
version, version $0$ is by convention the ``commit to all effecting
data'' algorithm, and for v6 only version $0$ is defined, with empty
associated data for the Transparent, Sapling, Orchard, and Issuance
bundles --- hence the constant $[\mathtt{0x00}]$ the parser demands
(ZIP-246, \emph{Sighash versioning}).

\paragraph{Implementation.} The \texttt{ZSASchnorr} scheme wraps the
\texttt{secp256k1} crate: signing builds a keypair from $\mathsf{isk}$
and calls \texttt{sign\_schnorr\_with\_aux\_rand(msg, keypair,
\&[0u8;~32])} --- the ZIP's $a = [\mathtt{0x00}]^{32}$ ---
$\mathsf{ik}$ is the x-only public key, and \texttt{encode()} of key and
signature each prepend \texttt{ALGORITHM\_BYTE} $= \mathtt{0x00}$, with
\texttt{decode()} rejecting any other lead byte (\texttt{orchard-zsa},
\texttt{src/issuance/auth.rs} lines~163--245 @ \texttt{cf801a5d}).
Bundle construction is a typestate chain --- \texttt{AwaitingNullifier}
$\rightarrow$ (\texttt{update\_rho}) $\rightarrow$
\texttt{AwaitingSighash} $\rightarrow$ (\texttt{prepare(sighash)})
$\rightarrow$ \texttt{Prepared} $\rightarrow$ (\texttt{sign(isk)})
$\rightarrow$ \texttt{Signed} --- so a bundle cannot be signed before its
notes' $\rho$ values and the sighash exist (\texttt{src/issuance.rs}
lines~261--304 and 404--599 @ \texttt{cf801a5d}). On the transaction
side, \texttt{v6\_signature\_hash} includes the txid
$\mathsf{issuance\_digest}$, and the builder signs the issue bundle with
\texttt{prepare(*shielded\_sig\_commitment)} where that commitment is
\texttt{signature\_hash(tx, SignableInput::Shielded, txid\_parts)} ---
the same $\mathsf{SIGHASH\_ALL}$ shielded digest the Sapling and Orchard
signatures sign (\texttt{librustzcash-zsa},
\texttt{transaction/sighash\_v6.rs} lines~10--51 and
\texttt{transaction/builder.rs} lines~1416--1417, 1453--1460 @
\texttt{3f9b53fc}). The builder also realizes the $\nfsym_{0,0}$ rule
constructively: \texttt{first\_nullifier} returns the first Action's
nullifier only for an OrchardZSA bundle, and building with an issuance
builder but no such bundle fails with
\texttt{BundleTypeNotSatisfiable}
(\texttt{transaction/builder.rs} lines~1356--1366 and 1594--1602 @
\texttt{3f9b53fc}).

\paragraph{A divergence to flag.} ZIP-246's txid and signature digests
include a sixth branch, T.6/S.6 $\mathsf{memo\_digest}$ (ZIP-246,
\emph{TxId Digest} and \emph{Signature Digest}), but
\texttt{librustzcash-zsa}'s \texttt{to\_hash} hashes only the header,
transparent, Sapling, Orchard, and issuance branches --- the fork does
not implement ZIP-231 memo bundles (\texttt{transaction/txid.rs}
lines~415--460 @ \texttt{3f9b53fc}). The digest the fork code signs
is therefore not
byte-identical to the digest ZIP-246 specifies; whichever is corrected,
the issuance signature's coverage claim above should be read against the
fork's five-branch digest today.

\subsection{Transaction digests}\label{sec:zsa-issuance-digests}

The v6 digest tree extends the ZIP-244 structure the \emph{Orchard
Guide} presents in \S``Conservation of value: value commitments and the
binding signature'': the txid digest is BLAKE2b-256 with personalization
$\texttt{ZcashTxHash\_} \,\|\, \mathrm{CONSENSUS\_BRANCH\_ID}$ over six
branches T.1--T.6 (header, transparent, Sapling, Orchard, issuance,
memo), and the signature digest has the same six-branch structure with
S.5 identical to T.5; a signature digest is produced for each transparent
input, Sapling input, OrchardZSA Action, and additionally for each
Issuance Action (ZIP-246, \emph{TxId Digest} and \emph{Signature
Digest}). The issuance branches are three nested hashes and one
authorizing hash, each over the wire encodings of
Table~\ref{tab:zsa-issue-wire}
(ZIP-246, \emph{txid\_digest} T.5 and \emph{A.4:
issuance\_auth\_digest}):%
\footnote{Two editorial defects in ZIP-246 to note when citing: the T.5
child list labels \texttt{issuerLength}/\texttt{issuer}/%
\texttt{issue\_actions\_digest} as T.5a/T.5b/T.5c, while the following
section headings call \texttt{issue\_actions\_digest} ``T.5a'' and
\texttt{issue\_notes\_digest} ``T.5a.i''; and the A.4 empty case reads
``In the case that the transaction has no Orchard Actions'' where ``no
Issuance Actions'' is evidently intended.}
\begin{align*}
\mathsf{issuance\_digest} &= \mathrm{BLAKE2b\text{-}256}\bigl(
  \texttt{ZTxIdSAIssueHash},\\
&\qquad \texttt{issuerLength} \,\|\, \texttt{issuer} \,\|\,
  \mathsf{issue\_actions\_digest}\bigr),\\
\mathsf{issue\_actions\_digest} &= \mathrm{BLAKE2b\text{-}256}\bigl(
  \texttt{ZTxIdIssuActHash},\\
&\qquad \textstyle\|_{\text{actions}}\;
  \texttt{assetDescHash} \,\|\,
  \mathsf{issue\_notes\_digest} \,\|\, \texttt{flagsIssuance}\bigr),\\
\mathsf{issue\_notes\_digest} &= \mathrm{BLAKE2b\text{-}256}\bigl(
  \texttt{ZTxIdIAcNoteHash},\\
&\qquad \textstyle\|_{\text{notes}}\;
  \texttt{recipient} \,\|\, \texttt{value} \,\|\,
  \texttt{rho} \,\|\, \texttt{rseed}\bigr),\\
\mathsf{issuance\_auth\_digest} &= \mathrm{BLAKE2b\text{-}256}\bigl(
  \texttt{ZTxAuthZSAOrHash},\
  \texttt{issueAuthSig}\ \text{field encoding}\bigr),
\end{align*}
where the A.4 preimage is the ZIP-230 \texttt{IssueAuthSignature}
composite
\[
\texttt{sizeSighashInfo} \,\|\, \texttt{sighashInfo} \,\|\,
\texttt{sizeSignature} \,\|\, \texttt{signature}
\]
(ZIP-230, \emph{Issuance Authorization Signature}), and a transaction
without issuance components hashes the empty string under the same
personalizations.

The code computes exactly these. On the txid side,
\texttt{hash\_issue\_bundle\_txid\_data} hashes
\[
\mathrm{compactSize}(|\mathsf{ik\_encoding}|) \,\|\,
\mathsf{ik\_encoding} \,\|\, \mathsf{issue\_actions\_digest}
\]
under \texttt{ZTxIdSAIssueHash}; the actions digest hashes, per action,
the $32$-byte hash, the notes digest, and the flags byte; the notes
digest hashes, per note, the $43$-byte recipient, the $8$-byte
little-endian value, and the two $32$-byte fields. On the authorizing
side, \texttt{hash\_issue\_bundle\_auth\_data} hashes
\[
\mathrm{compactSize}(|\mathsf{sighashInfo}|) \,\|\,
\mathsf{sighashInfo} \,\|\, \mathrm{compactSize}(65) \,\|\,
[\mathtt{0x00}] \,\|\, \sigma,
\]
asserting the $65$-byte length and the $\mathtt{0x00}$ lead. The
empty-bundle variants hash the empty string, and all four
personalizations match ZIP-246; both digests live in
\texttt{src/bundle/commitments/\allowbreak issuance.rs} lines~11--86
(\texttt{orchard-zsa} @ \texttt{cf801a5d}). For the wiring into
transaction identifiers, \texttt{TxIdDigester::digest\_issue} maps the
bundle to its \texttt{commitment()}; \texttt{to\_hash} appends the
digest only when the version \texttt{has\_orchard\_zsa()}; and
\texttt{BlockTxCommitmentDigester} takes the authorizing commitment
(\texttt{librustzcash-zsa}, \texttt{transaction/txid.rs}
lines~381--383, 445--452, 595--598 @ \texttt{3f9b53fc}).

Deterministic regression digests are pinned in the tests --- for a seeded
two-asset bundle,
\begin{center}\small
$\mathsf{issuance\_digest}$ =
\texttt{ee70e3b61674fd0428ac0020cc4fc5819386e39c4eb3c63357c84c998195bcdb},\\
$\mathsf{issuance\_auth\_digest}$ =
\texttt{6df77af7b5323d99376336b770e4c5b06ffc195de81ac7692d1b08b6eb19534d}
\end{center}
--- and cross-implementation test vectors ship in
\texttt{src/test\_vectors/} (\texttt{asset\_base.rs},
\texttt{issuance\_auth\_sig.rs}), exercised by the asset-base and
authorization tests (\texttt{orchard-zsa},
\texttt{src/bundle/commitments/issuance.rs} lines~156--190 and
\texttt{src/test\_vectors/} @ \texttt{cf801a5d}). The upstream ZIP-227
Test Vectors section, by contrast, still reads LINK TBD.

\subsection{Global issuance state and the consensus rules}%
\label{sec:zsa-issuance-state}

Supply integrity is the one property issuance cannot delegate to
cryptography: no transaction field simultaneously witnesses everything
issued and everything burnt for an Asset, so --- along the lines of the
ZIP-209 pool turnstile --- every node maintains a per-Asset circulation
record (ZIP-227, \emph{Rationale for Global Issuance State}).

\begin{definition}[Global issuance state]\label{def:zsa-issued-assets}
The map
\[
\begin{aligned}
\mathsf{issued\_assets} : \Ep(\F_{p_{\Pallas}})^{\ast} &\to
  \{0\,..\,\mathsf{MAX\_ISSUE}\} \times \{0,1\} \times
  \mathsf{Note}^{\mathsf{Issue}},\\
\mathsf{AssetBase} &\mapsto (\mathsf{balance}, \mathsf{final},
  \mathsf{note_{ref}}),
\end{aligned}
\]
records, for every issued Asset, the amount in circulation (issued minus
burnt), the finalization status --- which can never change from $1$ to
$0$ --- and the Asset's reference note. A missing entry
($\mathsf{issued\_assets}(\mathsf{AssetBase}) = \bot$) is read as
$\mathsf{balance} = 0$, $\mathsf{final} = 0$,
$\mathsf{note_{ref}} = \bot$; the constant $\mathsf{MAX\_ISSUE} =
2^{64} - 1$ caps the total supply of any Custom Asset (ZIP-227,
\emph{Global Issuance State}).
\end{definition}

The map threads through the chain deterministically: it MUST be updated
by a node during the processing of any transaction that contains burn
information or an issuance bundle; the input state for the OrchardZSA
activation block is the empty map; the input state of a block's first
transaction is the final state of the preceding block; each subsequent
transaction's input state is the preceding transaction's output state;
and the block's final state is its last transaction's output state
(ZIP-227, \emph{Management of the Global Issuance State}).

\paragraph{Per-transaction rules.} For a v6 transaction (ZIP-227,
\emph{Specification: Consensus Rule Changes}):
\begin{enumerate}[label=(T\arabic*)]
\item \texttt{nActionGroupsOrchard} MUST be $0$ or $1$, and
  \texttt{nAGExpiryHeight} MUST be $0$; the v6 parser of
  \texttt{librustzcash-zsa} enforces both at deserialization, in
  \texttt{transaction/components/\allowbreak orchard.rs}
  lines~107--132 @ \texttt{3f9b53fc};
\item the output state $\mathsf{issued\_assets_{OUT}}$ MUST be
  initialized to the input state $\mathsf{issued\_assets_{IN}}$;
\item the $\mathsf{assetBurn}$ set MUST satisfy the ZIP-226 rules
  (\S\ref{sec:zsa-transfer}: tuples $(\mathsf{AssetBase}, \mathsf{v})$
  with the Native Asset excluded, $0 < \mathsf{v} \le
  \mathsf{MAX\_BURN\_VALUE} = 2^{63}-1$, and no duplicate
  $\mathsf{AssetBase}$; ZIP-226, burn rules); and
\item for every $(\mathsf{AssetBase}, \mathsf{v}) \in
  \mathsf{assetBurn}$, the node MUST check $\mathsf{balance} \ge
  \mathsf{v}$ for the entry of $\mathsf{AssetBase}$ in
  $\mathsf{issued\_assets_{OUT}}$ and subtract $\mathsf{v}$ from that
  balance, \emph{prior to} processing the issuance bundle.
\end{enumerate}

\paragraph{Issuance-bundle rules.} For a transaction carrying an
issuance bundle (ZIP-227, \emph{Specification: Consensus Rule Changes}):
\begin{enumerate}[label=(I\arabic*)]
\item the transaction MUST also contain at least one OrchardZSA Action
  Group (the $\rho$/replay rationale of
  \S\ref{sec:zsa-issuance-notes});
\item the encoding $\mathsf{ik\_encoding}$ MUST be taken from the
  \texttt{issuer} field and MUST start with the byte $\mathtt{0x00}$
  (BIP-340);
\item the signature MUST satisfy
  $\mathsf{IssueAuthSig.Validate}(\mathsf{ik}, \mathsf{SigHash},
  \mathtt{issueAuthSig}) = 1$;
\item for each Issue Action: every \texttt{IssueNoteDescription} MUST be
  a valid ZIP-230 encoding; the $\mathsf{AssetBase}$ is derived from
  $\mathsf{assetDescHash}$ and $\mathsf{issuer}$
  (\S\ref{sec:zsa-issuance-asset}); and the Asset's $\mathsf{final}$
  bit in $\mathsf{issued\_assets_{OUT}}$ MUST NOT be $1$;
\item if $\mathsf{note_{ref}} = \bot$, the first Issue Note MUST be a
  reference note (\S\ref{sec:zsa-issuance-notes}), and the node MUST set
  $\mathsf{note_{ref}}$ to it;
\item for every Issue Note $\mathsf{note}_j$: its $\rho$ MUST equal the
  value computed per \S\ref{sec:zsa-issuance-notes}; the supply update
  \[
  \mathsf{issued\_assets_{OUT}}.\mathsf{balance} + \mathsf{v}_j
  \le \mathsf{MAX\_ISSUE}, \qquad
  \mathsf{issued\_assets_{OUT}}.\mathsf{balance} \mathrel{{+}{=}}
  \mathsf{v}_j
  \]
  MUST hold and be applied; and the node MUST compute the note
  commitment per ZIP-226;
\item if $\mathsf{finalize} = 1$ in \texttt{flagsIssuance}, the node
  MUST set $\mathsf{final} = 1$.
\end{enumerate}
The commitments enter the tree in a fixed order: for each Action Group
of the OrchardZSA bundle, the commitment of every new note per Action;
then, for each Issue Action of the issuance bundle, the commitment of
every Issue Note (ZIP-227, \emph{Addition to the Note Commitment Tree}).

The rationale ties the pieces together: balances of issued Assets must
never go negative, $\mathsf{MAX\_ISSUE}$ is a practical limit that also
lets an issuer create an Asset's complete supply in a single transaction,
and the same map records finalization so that further issuance of a
finalized Asset is rejected (ZIP-227, \emph{Rationale for Global
Issuance State}). Beyond consensus, the ZIP RECOMMENDS that full
validators keep a mutable \texttt{issuanceSupplyInfoMap} from AssetId to
(\texttt{totalSupply}, \texttt{finalize}) for wallets and applications
tracking total supply (ZIP-227, \emph{Implementing Zcash Nodes}).

\paragraph{The library realization.} Function
\texttt{verify\_issue\_bundle} takes the bundle, the sighash, a
caller-supplied view of global state (\texttt{get\_global\_records}),
and the first nullifier. It rejects any sighash kind other than
\texttt{AllEffecting}, verifies the BIP-340 signature over the $32$-byte
sighash with the bundle's $\mathsf{ik}$, folds the actions over a
\texttt{BTreeMap<AssetBase, AssetRecord>}, and returns the updated
records --- \texttt{amount}, \texttt{is\_finalized},
\texttt{reference\_note} --- for the caller to persist
(\texttt{orchard-zsa}, \texttt{src/issuance.rs} lines~647--802 @
\texttt{cf801a5d}). The fold realizes the per-action and per-note rules
as typed errors:
\begin{itemize}
\item \texttt{IncorrectRhoDerivation} --- a note's $\rho$ differs from
  the recomputed derivation of \S\ref{sec:zsa-issuance-notes};
\item \texttt{MissingReferenceNoteOnFirstIssuance} --- no record exists
  for the Asset, and the action's first note is not a reference note;
\item \texttt{IssueActionPreviouslyFinalizedAssetBase} --- reissuance
  against a finalized record;
\item \texttt{IssueBundleIkMismatchAssetBase} and
  \texttt{AssetBaseCannotBeIdentityPoint} --- a note's Asset Base
  differs from the one \texttt{IssueAction::verify} derives from
  $(\mathsf{ik}, \mathsf{assetDescHash})$, or the derived point is the
  Pallas identity;
\item \texttt{ValueOverflow} --- the checked \texttt{u64} value sum
  overflows, the code's realization of $\mathsf{MAX\_ISSUE} = 2^{64}-1$
  (\texttt{src/issuance.rs} lines~216--244; \texttt{src/value.rs}
  lines~143--145).
\end{itemize}
A signature-skipping variant,
\texttt{check\_issue\_bundle\_without\_sighash}, exists for verifying
historical blocks from a trusted checkpoint (\texttt{src/issuance.rs}
lines~647--705 @ \texttt{cf801a5d}).

Three boundary observations, each load-bearing for anyone re-deriving
consensus from these libraries:
\begin{itemize}
\item \emph{Empty-action divergence.} The implementation rejects an
  action with no notes unless it is finalized
  (\texttt{IssueActionWithoutNoteNotFinalized},
  \texttt{src/issuance.rs} lines~216--219), but ZIP-227's consensus-rule
  list contains no such MUST, and ZIP-230 merely notes that
  \texttt{nNotes} may be zero ``to only finalize''. A ZIP-conformant
  validator would accept an empty, non-finalizing action that the
  reference implementation rejects.
\item \emph{First-issuance detection.} The ZIP keys the reference-note
  requirement on $\mathsf{note_{ref}} = \bot$; the code keys it on the
  complete absence of an \texttt{AssetRecord}
  (\texttt{get\_global\_records} returning \texttt{None}). The two are
  equivalent under the invariant that $\mathsf{note_{ref}}$ is set
  exactly on first issuance --- rule (I5) --- and no rule ever clears
  it, so an entry exists if and only if its $\mathsf{note_{ref}}$ is
  set.
\item \emph{Node-side enforcement locus.} Neither
  \texttt{orchard-zsa} nor \texttt{librustzcash-zsa} contains the
  $\mathsf{issued\_assets}$ chain-state maintenance or the consensus
  caller of \texttt{verify\_issue\_bundle}: the libraries return updated
  records for a node to persist, and the QEDit node integration is not
  among our local ground-truth repositories. The global-state rules are
  therefore \emph{specified and library-supported}, with their node
  enforcement unverified locally.
%% TODO: verify the issued_assets chain-state fold in the QEDit node
%% fork (zebra/zcashd) once a local checkout is available, and cite it
%% here.
\end{itemize}

\subsection{Finalization}\label{sec:zsa-issuance-final}

The $\mathsf{finalize}$ boolean signals that no further issuance of the
specific Custom Asset is possible: every issuance action carries it in
\texttt{flagsIssuance} (a requirement of the ZIP's Requirements section),
setting it sets the Asset's $\mathsf{final}$ bit --- rule (I7) --- and
transactions attempting to issue further amounts of a previously
finalized Asset are rejected --- rule (I4) (ZIP-227,
\emph{Requirements} and \emph{Issuance Action}). The bit is monotone:
consensus state never changes $\mathsf{final}$ from $1$ to $0$
(Definition~\ref{def:zsa-issued-assets}).

Finalization is the ZIP's whole lifecycle policy, and its Concrete
Applications read as patterns over one bit (ZIP-227, \emph{Concrete
Applications}):
\begin{itemize}
\item \emph{One-time issuance}: set $\mathsf{finalize} = 1$ on the first
  action; the entire supply exists from one transaction ---
  $\mathsf{MAX\_ISSUE}$ is sized to permit this
  (\S\ref{sec:zsa-issuance-state}).
\item \emph{Stopping supply}: issue with $\mathsf{finalize} = 1$ in a
  final issuance transaction, or in one containing a single note of
  value $0$ (the wire format admits note-less finalizing actions as
  well, Table~\ref{tab:zsa-issue-wire}; the library insists on the
  finalizing variant, \S\ref{sec:zsa-issuance-state}).
\item \emph{NFTs}: issue multiple actions in one transaction, each with
  $\mathsf{value} = 1$ at the Asset's fundamental unit and
  $\mathsf{finalize} = 1$ --- each action's Asset is thereafter unique
  and non-replicable.
\end{itemize}

Finalization is also the entirety of the compromise story. No key
rotation exists for issuance keys; on compromise of $\mathsf{isk}$, the
ZIP's guidance is to set $\mathsf{finalize} = 1$ for all Assets issued
with that key, switch to a new $\mathsf{isk}$ --- hence a new issuer
identifier, since $\mathsf{issuer} = \mathsf{ik\_encoding}$ --- and
issue new Assets with new AssetIds (ZIP-227, \emph{Issuance Key
Compromise}). The old Assets keep circulating and burning; they simply
never grow again.

\subsection{Fees}\label{sec:zsa-issuance-fees}

The fee treatment of issuance is the sharpest fork/upstream split in the
ZSA stack, and we bin it explicitly. The QEDit fork of ZIP-317 adds two
contributions to the logical-action count of the conventional fee (the
\emph{Wallet Guide}, \S``Fees (ZIP-317)'', for the surrounding
machinery: $\mathit{marginal\_fee} = 5000$ zatoshis,
$\mathit{grace\_actions} = 2$):
\begin{align*}
\mathit{contribution}_{\mathsf{ZSAIssuance}} &=
  \mathit{nZSAIssueNotes},\\
\mathit{contribution}_{\mathsf{ZSACreation}} &=
  \mathit{creation\_cost} \cdot \mathit{nAssetCreations},
\end{align*}
where $\mathit{nZSAIssueNotes}$ counts \texttt{IssueNote} outputs,
$\mathit{nAssetCreations}$ counts Custom Assets newly added to the
global issuance state, and $\mathit{creation\_cost} = 100$ logical
actions (fork ZIP-317, \emph{Fee calculation} @ zips-zsa
\texttt{fd71419}). The stated rationale: every new Asset adds a row that
validators track in perpetuity
(\S\ref{sec:zsa-issuance-state}), and the two-orders-of-magnitude
creation cost disincentivizes junk assets.

Upstream, these contributions no longer exist: ZIP-317's change history
records, on 2026-06-26, ``Removed the ZSA issuance and asset-creation
fee contributions (ZIP 227)'' (ZIP-317, change history @ zips
\texttt{6961098}) --- while ZIP-227's own \emph{Changes to ZIP 317}
section still claims the conventional fee accounts for both. The
formulas above survive only in the QEDit fork and its libraries:
\emph{fork-specified}, implemented in \texttt{librustzcash-zsa} behind
\texttt{cfg(zcash\_unstable = "nu7")} (\texttt{fees/zip317.rs},
\texttt{CREATION\_COST} $= 100$ @ \texttt{3f9b53fc}), and awaiting an
upstream resolution that the NU7 process (\S\ref{sec:zsa-scope}) has yet
to provide.

\section{Transfer and burn (ZIP-226)}\label{sec:zsa-transfer}

ZIP-226 (\emph{Transfer and Burn of Zcash Shielded Assets}; Status: Draft;
owners Pablo Kogan and Vivek Arte of QEDit, Daira-Emma Hopwood, and Jack
Grigg) defines, jointly with the issuance protocol of ZIP-227
(\S\ref{sec:zsa-issuance}), the OrchardZSA transfer and burn protocol. We
develop it as four deltas against the Orchard baseline, each citing the
construction it modifies in the \emph{Orchard Guide} rather than re-deriving
it: the note gains one field, the \emph{Asset Base} (\S\ref{sec:zsa-note});
the note commitment becomes a case split that keeps ZEC commitments identical
to NU5 (\S\ref{sec:zsa-notecommit}); the value commitment trades its fixed
value base for the per-asset base, which turns the binding signature into a
per-asset balance check and gives burn its consensus meaning
(\S\ref{sec:zsa-value}, \S\ref{sec:zsa-burn}); and dummy spends give way to
\emph{Split Inputs} for Custom Assets (\S\ref{sec:zsa-split}). We then state
the extended Action statement (\S\ref{sec:zsa-circuit}) and the carrier
format with its digests and fees (\S\ref{sec:zsa-carrier}).

Ground truth for this section: upstream \texttt{zips} at commit
\texttt{69610984}, the QEDit fork \texttt{zips-zsa} at \texttt{fd71419a},
crate \texttt{orchard-zsa} at \texttt{cf801a5d} (branch \texttt{zsa1}), and
\texttt{librustzcash-zsa} at \texttt{3f9b53fc} (branch \texttt{zsa1}); we
verified every cited artifact at these commits. ZIP-226 names the QED-it
\texttt{zcash-test-vectors} and the QED-it forks of \texttt{zcashd},
\texttt{orchard}, \texttt{librustzcash}, and \texttt{halo2} as its test
vectors and reference implementations. Deployment standing is fixed once, in
\S\ref{sec:zsa-scope}, and not relitigated per object; two upstream facts
nonetheless shape the text. First, ZIP-226 states that ZSA ``is scheduled to
be deployed in Network Upgrade 7 (NU7)'' and ``currently assumes that this
will be the case'', while the transaction format it rides on has been
withdrawn upstream (\S\ref{sec:zsa-carrier}). Second, upstream ZIP-226 is now
defined relative to the protocol with ZIP-2005 (Ironwood Quantum
Recoverability) applied, which is expected to deploy before any ZSA
activation; the QEDit fork's ZIP-226 lacks this clause, and the pinned
implementation does not implement ZIP-2005 constructs. Where the difference
is material --- the split-note $\psi^{\nfsym}$
(\S\ref{sec:zsa-split}) and the lead byte (\S\ref{sec:zsa-carrier}) --- we
flag it in place.

\subsection{One new field: the OrchardZSA note}\label{sec:zsa-note}

\begin{definition}[OrchardZSA note]\label{def:zsa-note}
An OrchardZSA note extends the Orchard note --- the tuple
$(\mathsf{addr}, v, \rho, \psi, \rcm)$ of the \emph{Orchard Guide},
``Representation of a note: notes and note commitments'' --- by exactly one
field:
\[
\mathbf{n} :=
(\mathsf{addr},\, v,\, \mathsf{AssetBase},\, \rho,\, \psi,\, \rcm),
\]
where $\mathsf{AssetBase} \in \Ep^{*} := \Ep(\F_{p_{\Pallas}}) \setminus
\{\mathcal{O}\}$ is ``a valid group element that is not the identity and is
not bottom'' (ZIP-226; the ZIP writes $\mathbb{P}^*$ for $\Ep^{*}$). All
other components are as in protocol specification \S3.2. Formally,
\[
\begin{aligned}
\mathsf{Note}^{\mathsf{OrchardZSA}} :={}&
  \{0,1\}^{88} \times \mathsf{KA}^{\mathsf{Orchard}}.\mathsf{Public}
  \times \{0, \dots, 2^{64}-1\} \times \Ep^{*} \\
&\times \F_{p_{\Pallas}} \times \F_{p_{\Pallas}}
  \times \NoteCommit^{\mathsf{Orchard}}.\mathsf{Trapdoor},
\end{aligned}
\]
where the ZIP's $\F_{q_P}$ is this series' $\F_{p_{\Pallas}}$ and
$\ell_{\mathsf{value}} = 64$ (ZIP-226).
\end{definition}

The byte encoding of the new field is the star-encoding of the series'
conventions: ZIP-226 sets $\mathsf{asset\_base} :=
\mathrm{LEBS2OSP}_{\ell_P}(\mathrm{repr}_{\mathbb{P}}(\mathsf{AssetBase}))$,
which is $\mathsf{AssetBase}^\star \in \{0,1\}^{256}$ in $32$ bytes.

An Asset Base is never a free choice of the transactor. Every Custom Asset's
base is derived by the ZIP-227 issuance chain --- constructed in
\S\ref{sec:zsa-issuance}, reproduced here because the soundness argument of
\S\ref{sec:zsa-split} depends on its shape:
\begin{align*}
\mathsf{assetDescHash} &:=
  \mathrm{BLAKE2b\text{-}256}\bigl(\texttt{"ZSA-AssetDescCRH"},
  \ \mathsf{asset\_desc}\bigr), \\
\mathsf{AssetId} &:= (\mathsf{issuer},\ \mathsf{assetDescHash}), \\
\mathsf{EncodeAssetId}(\mathsf{AssetId}) &:= \mathtt{0x00} \,\|\,
  \mathsf{issuer}
  \,\|\, \mathsf{assetDescHash}, \\
\mathsf{AssetDigest} &:=
  \mathrm{BLAKE2b\text{-}512}\bigl(\texttt{"ZSA-Asset-Digest"},
  \ \mathsf{EncodeAssetId}(\mathsf{AssetId})\bigr), \\
\mathsf{AssetBase} &:= \GroupHash^{\Pallas}\bigl(\mathtt{z.cash:OrchardZSA},
  \ \mathsf{AssetDigest}\bigr)
\end{align*}
(ZIP-227, which names the last map $\mathsf{ZSAValueBase}$). For the native
Asset (ZEC/TAZ) the Asset Base is instead \emph{defined} to be
\[
V^{\mathsf{Orchard}} :=
\GroupHash^{\Pallas}(\mathtt{z.cash:Orchard\text{-}cv}, \texttt{"v"}),
\]
the fixed value-commitment base of the \emph{Orchard Guide}, ``Conservation
of value: value commitments and the binding signature'' (ZIP-226). This
single identification is what keeps every ZEC computation below identical to
NU5: ZIP-226 uses $\mathsf{ValueCommit}^{\mathsf{Orchard}}_{\rcv}(v)$ as
$\mathsf{ValueCommit}^{\mathsf{OrchardZSA}}_{\rcv}(V^{\mathsf{Orchard}}, v)$.
The implementation realizes it as \texttt{AssetBase::zatoshi()} $=$
\texttt{hash\_to\_curve("z.cash:Orchard-cv")(b"v")}
(\textsf{orchard-zsa}, \texttt{src/note/asset\_base.rs},
\texttt{src/constants/fixed\_bases.rs}).

\subsection{The note commitment: a case split}\label{sec:zsa-notecommit}

\begin{definition}[OrchardZSA note commitment]\label{def:zsa-notecommit}
For an OrchardZSA note,
\[
\begin{aligned}
&\NoteCommit^{\mathsf{OrchardZSA}}_{\rcm}(\gd^\star, \pkd^\star, v, \rho, \psi,
  \mathsf{AssetBase}) \\
&\qquad :=
\begin{cases}
\NoteCommit^{\mathsf{Orchard}}_{\rcm}(\gd^\star, \pkd^\star, v, \rho, \psi)
  & \text{if } \mathsf{AssetBase} = V^{\mathsf{Orchard}}, \\
\cmsym_{\mathsf{ZSA}} & \text{otherwise},
\end{cases}
\end{aligned}
\]
where
\[
\begin{aligned}
\cmsym_{\mathsf{ZSA}} :={}&
  \Sinsemilla_{\mathtt{z.cash:ZSA\text{-}NoteCommit\text{-}M}}
  \bigl(\gd^\star \,\|\, \pkd^\star \,\|\, \mathrm{LE}_{64}(v) \,\|\,
  \mathrm{LE}_{255}(\rho) \,\|\, \mathrm{LE}_{255}(\psi) \,\|\,
  \mathsf{AssetBase}^\star\bigr) \\
&\dotplus\; [\rcm]\,
  \GroupHash^{\Pallas}\bigl(\mathtt{z.cash:Orchard\text{-}NoteCommit\text{-}r},
  \ \texttt{""}\bigr)
\end{aligned}
\]
(ZIP-226; the ZIP writes $\mathrm{I2LEBSP}$ for the series'
$\mathrm{LE}_k$). The signature is
\[
\begin{aligned}
\NoteCommit^{\mathsf{OrchardZSA}} :{}&
  \NoteCommit^{\mathsf{Orchard}}.\mathsf{Trapdoor}
  \times \{0,1\}^{256} \times \{0,1\}^{256}
  \times \{0, \dots, 2^{64}-1\} \\
&\times \F_{p_{\Pallas}} \times \F_{p_{\Pallas}}
  \times \Ep^{*}
  \;\to\; \NoteCommit^{\mathsf{Orchard}}.\mathsf{Output},
\end{aligned}
\]
with trapdoor and output types those of $\NoteCommit^{\mathsf{Orchard}}$,
unchanged (ZIP-226).
\end{definition}

The two branches carry the design's two promises. In the native branch the
Asset Base is not hashed at all: a ZEC note commitment under OrchardZSA
\emph{is} the NU5 commitment, bit for bit. In the ZSA branch the message
appends $\mathsf{AssetBase}^\star$ ($256$ bits) after $\psi$, against the
baseline layout of the \emph{Orchard Guide}, ``Representation of a note'':
\[
\gd^\star \,\|\, \pkd^\star \,\|\, \mathrm{LE}_{64}(v) \,\|\,
\mathrm{LE}_{255}(\rho) \,\|\, \mathrm{LE}_{255}(\psi)
\;\longmapsto\;
\cdots \,\|\, \mathrm{LE}_{255}(\psi) \,\|\, \mathsf{AssetBase}^\star,
\]
under the fresh hash domain
\texttt{z.cash:ZSA-NoteCommit-M} but the \emph{same} randomness base
$\GroupHash^{\Pallas}(\mathtt{z.cash:Orchard\text{-}NoteCommit\text{-}r},
\texttt{""})$. ZIP-226's rationale: ``the nested structure of the
Sinsemilla-based commitment allows us to add the Asset Base as a final
recursive step'', and the output ``is still indistinguishable from the
original Orchard ZEC note commitments, by definition of the Sinsemilla hash
function''. Because the trapdoor and output types are unchanged, ZSA
commitments land in the same note commitment tree as ZEC commitments and are
indistinguishable from them there (ZIP-226). One caveat bounds the claim:
Orchard (v5) note commitments \emph{are} distinguishable from OrchardZSA
(v6) note commitments, because the two protocols use different transaction
versions (ZIP-226).

\begin{remark}[Implementation]\label{rem:zsa-notecommit-impl}
Function \texttt{NoteCommitment::derive} computes \emph{both} commitments.
For ZEC it calls \texttt{CommitDomain::new} on the prefix
\texttt{"z.cash:Orchard-NoteCommit"}; for ZSA it calls
\texttt{new\_with\_separate\_domains} on the $M$-prefix
\texttt{"z.cash:ZSA-NoteCommit"} and the shared $r$-prefix
\texttt{"z.cash:Orchard-NoteCommit"}. It then selects between
them in constant time on \texttt{asset.is\_zatoshi()} (\textsf{orchard-zsa},
\texttt{src/note/commitment.rs}). The \textsf{sinsemilla} crate appends the
\texttt{-M} and \texttt{-r} suffixes to these prefixes (\textsf{sinsemilla}
0.1.0, \texttt{src/lib.rs}).
\end{remark}

The nullifier of an ordinary (non-split) OrchardZSA note is generated
exactly as in Orchard (protocol specification \S4.16; \emph{Orchard Guide},
``Prevention of double-spending: nullifiers''); only Split Inputs use the
modified, randomized formula of \S\ref{sec:zsa-split} (ZIP-226).

\subsection{Per-asset value: the variable-base value commitment}
\label{sec:zsa-value}

\begin{definition}[OrchardZSA net value commitment]\label{def:zsa-cv}
Each Action commits to its net value change in its Asset:
\[
\begin{aligned}
\cvsym^{\mathsf{net}} &:=
\mathsf{ValueCommit}^{\mathsf{OrchardZSA}}_{\rcv}(\mathsf{AssetBase},
  v^{\mathsf{net}})
= [\,v^{\mathsf{net}}\,]\,\mathsf{AssetBase}
+ [\rcv]\,R^{\mathsf{Orchard}}, \\
v^{\mathsf{net}} &:= v_{\mathrm{old}} - v_{\mathrm{new}},
\qquad
R^{\mathsf{Orchard}} :=
\GroupHash^{\Pallas}(\mathtt{z.cash:Orchard\text{-}cv}, \texttt{"r"}),
\end{aligned}
\]
where $\mathsf{AssetBase}$ is the Action's Asset Base and the randomness
base $R^{\mathsf{Orchard}}$ is unchanged (ZIP-226).
\end{definition}

Exactly one thing changes against the baseline
$\cvsym^{\mathsf{net}} = [v_{\mathrm{old}} - v_{\mathrm{new}}]\,
V^{\mathsf{Orchard}} + [\rcv]\,R^{\mathsf{Orchard}}$ of the \emph{Orchard
Guide}, ``Conservation of value'': the fixed-base multiplication by
$V^{\mathsf{Orchard}}$ becomes a variable-base multiplication by the
Action's $\mathsf{AssetBase}$, while the randomness term keeps its fixed
base (ZIP-226). For ZEC, $\mathsf{AssetBase} = V^{\mathsf{Orchard}}$ and the
commitment reduces to the NU5 one (\S\ref{sec:zsa-note}).

The prover still signs the SIGHASH with
$\mathsf{bsk} = \sum_{i=1}^{n} \rcv_i$ over the $n$ Actions of a transfer
(ZIP-226); the verifier's key acquires burn terms --- one zero-trapdoor
commitment subtracted per $\mathsf{assetBurn}$ entry, each a bare multiple
of its base --- displayed once in this volume as \eqref{eq:zsa-bvk}
(\S\ref{sec:zsa-split-attack}). Why this checks balance \emph{per asset}:
``the committed values must sum to zero per Asset Base, as the Pedersen
commitments add up homomorphically only with respect to the same value
base point'' (ZIP-226). Partitioning the Actions into a ZEC set and a
Custom Asset set, the custom-asset commitments must cancel against the
declared burns up to pure randomness terms --- the cancellation equation
is displayed in \S\ref{sec:zsa-split-input} --- so
$\mathsf{bvk} = [\mathsf{bsk}]\,R^{\mathsf{Orchard}}$ holds if and only
if, per Custom Asset, the sum of net Action values equals its
$\mathsf{assetBurn}$ entry (or $0$ if absent), and the ZEC component
equals $v_{\mathrm{balance}}^{\mathsf{Orchard}}$ (ZIP-226). The
binding-signature mechanics --- $\mathsf{bsk}$/$\mathsf{bvk}$ cancellation
and RedPallas on the base $R^{\mathsf{Orchard}}$ --- are those of the
\emph{Orchard Guide}, ``Conservation of value'', and we do not restate them.

\begin{remark}[Implementation: $\mathsf{derive\_bvk}$]\label{rem:zsa-bvk-impl}
Function \texttt{derive\_bvk\_raw} computes the ZIP equation verbatim:
$\sum \cvsym^{\mathsf{net}}$ minus the zero-trapdoor commitment of
\texttt{value\_balance} at \texttt{AssetBase::zatoshi()} minus, over the
burn set, the zero-trapdoor commitment of each \texttt{value} at its
\texttt{asset} (\textsf{orchard-zsa}, \texttt{src/bundle.rs}). Function
\texttt{ValueCommitment::derive} computes $V \cdot \mathsf{value} + R \cdot
\rcv$ with $V$ the point \texttt{asset.cv\_base()} and $R$ the point
\texttt{hash\_to\_curve("z.cash:Orchard-cv")} applied to \texttt{b"r"},
the value signed per the sign of the value sum (\textsf{orchard-zsa},
\texttt{src/value.rs}).
\end{remark}

\subsection{Burn}\label{sec:zsa-burn}

Burning is the value-balance mechanism read as an exit door. The mechanism
``does NOT send Assets to a non-spendable address, it simply reduces the
total number of units of a given Custom Asset in circulation''; it is
enforced at consensus level as an extension of the value-balance mechanism,
and it ``makes it globally provable that a given amount of a Custom Asset
has been destroyed''. The Native Asset (ZEC/TAZ) cannot be burnt this way
(ZIP-226).

ZIP-226 attaches three consensus rules to the burn set, whose cardinality it
denotes $L$:
\begin{enumerate}[leftmargin=*]
\item for every $(\mathsf{AssetBase}, v) \in \mathsf{assetBurn}$:
  $\mathsf{AssetBase} \neq V^{\mathsf{Orchard}}$;
\item $v > 0$ and $v \leq \mathsf{MAX\_BURN\_VALUE} := 2^{63} - 1$;\footnote{%
ZIP-230's $\mathsf{AssetBurn}$ table instead says $\mathsf{valueBurn}$ is
``checked by consensus to be non-zero, and less than
$\mathsf{MAX\_BURN\_VALUE}$''. The implementation rejects only amounts
strictly greater than $\mathsf{MAX\_BURN\_VALUE}$, i.e.\ permits equality,
agreeing with ZIP-226 (\textsf{orchard-zsa},
\texttt{src/bundle/burn\_validation.rs}, which defines
$\mathsf{MAX\_BURN\_VALUE} = (1 \ll 63) - 1$). A specification-text defect,
not a consensus divergence.}
\item every $\mathsf{AssetBase}$ has at most one entry in
  $\mathsf{assetBurn}$.
\end{enumerate}

Burn data lives \emph{per Action Group}: the
$\mathsf{ActionGroupDescription}$ carries $\mathsf{nAssetBurn}$
(compactSize) and $\mathsf{vAssetBurn}$, an array of $40$-byte entries ---
$\mathsf{asset\_base}$ (byte[32]) followed by $\mathsf{valueBurn}$ (uint64,
little-endian) --- so that parties assembling separate Action Groups can
each burn their own Assets (withdrawn ZIP-230; \S\ref{sec:zsa-carrier}).

The transaction-level field $\mathsf{valueBalanceOrchard}$ (int64) remains
the net \emph{ZEC-only} flow, ``the net value of Orchard spends minus
outputs'' (withdrawn ZIP-230), and the transparent protocol is unchanged, so
a Custom Asset cannot be unshielded by transfer: ``All Custom Assets are
contained within the shielded pool'' (ZIP-226). Burning is the only exit.
What a burn reveals is exactly its $\mathsf{assetBurn}$ entry: ``the amount
and identifier of the Asset being burnt''. What transfers hide: ``the
overall balance is checked without revealing which (or how many distinct)
Assets are being transferred''; OrchardZSA ZEC commitments are
indistinguishable from non-native OrchardZSA commitments, and the amounts
and identifiers of transferred Assets stay private (ZIP-226).

A burn also updates ZIP-227's global issuance state:
$\mathsf{issued\_assets}(\mathsf{AssetBase}).\mathsf{balance} \in \{0,
\dots, \mathsf{MAX\_ISSUE}\}$ is ``computed as the amount of the Asset that has
been issued less the amount of the Asset that has been burnt'', and the map
MUST be updated during processing of any transaction containing burn
information (ZIP-227). Supply sufficiency --- that a burn not exceed the
circulating supply --- is therefore a state-level rule of the issuance
protocol, not one of ZIP-226's three per-transaction rules above; we
construct the state machine in \S\ref{sec:zsa-issuance}. The implementation
enforces both layers side by side: \texttt{validate\_bundle\_burn} returns
the per-transaction errors \texttt{ZatoshiAsset}, \texttt{ZeroAmount},
\texttt{InvalidAmount}, and \texttt{DuplicateAsset}, and the state-level
errors \texttt{AssetNotFoundInState} and \texttt{InsufficientSupply}
(\textsf{orchard-zsa}, \texttt{src/bundle/burn\_validation.rs}).

\subsection{Split Inputs}\label{sec:zsa-split-inputs}

Padding is where the asset dimension bites. An Orchard transfer with more
outputs than inputs pads with \emph{dummy} spends: zero-value notes whose
Merkle membership check is gated off by $v_{\mathrm{old}}$ (constraint A3 of
the \emph{Orchard Guide}, ``The Action statement, formally''). For Custom
Assets that padding is unsound, because the circuit forces the output note's
Asset Base to equal the input note's (\S\ref{sec:zsa-circuit}) --- and a
dummy input's Asset Base is chosen freely by the prover. ZIP-226 states the
attack that must be excluded: without enforcing that every Custom Asset
input exists in the note commitment tree, ``the prover could input a
multiple (or linear combination) of an existing Asset Base, and thereby
attack the network by overflowing the ZEC value balance and hence
counterfeiting ZEC funds''. The fix: dummy notes remain in use for ZEC notes
only; Custom Assets always prove Merkle membership, padding instead with
Split Inputs, which force the output's Asset Base to equal a genuinely
issued, $\GroupHash$-derived base (ZIP-226).

The construction is developed once, in \S\ref{sec:zsa-split}: the
\emph{Split Input} --- a copy of a previously issued, tree-resident note
of the target Asset, spent with $\mathsf{split\_flag} = 1$ and its value
zeroed in the value-commitment computation only
(Definition~\ref{def:zsa-split-input}); its randomized nullifier with the
offset point $L^{\mathsf{Orchard}}$, \eqref{eq:zsa-split-nf}; the
three-way divergence over the derivation of $\psi^{\nfsym}$
(\S\ref{sec:zsa-split-nf}); and the builder mechanics of padding per
asset, dummy notes for ZEC and split spends for Custom Assets
(\S\ref{sec:zsa-split-builder}). The three wallet strategies ZIP-226
offers for sourcing the copied note are enumerated in
\S\ref{sec:zsa-split-input} (\texttt{zip-0226.rst}:283--287). Burn entry
points reject duplicate Assets and amounts that do not fit in $63$ bits
(\texttt{add\_burn}, \textsf{orchard-zsa}, \texttt{src/builder.rs}).

\subsection{The Action statement, extended}\label{sec:zsa-circuit}

A bin statement first. ZIP-226's ``Circuit Statement'' section describes the
constraint changes informally and defers ``all modifications in the
Circuit'' to an external document; no statement at the rigor of the
\emph{Orchard Guide}'s Definition ``Orchard Action statement'' exists in any
specification. The implemented circuit is the de facto specification, and
every gate-level claim below cites it (\textsf{orchard-zsa},
\texttt{src/circuit/circuit\_zsa.rs}). Bin: \emph{designed and implemented,
but unspecified}.

ZIP-226 states the deltas against A1--A9 as follows:
\begin{enumerate}[label=(\alph*), leftmargin=*]
\item the Asset Base is witnessed once as an auxiliary input, and both note
commitment integrity constraints (A1, A2) switch to
$\NoteCommit^{\mathsf{OrchardZSA}}$ with that same witness --- one variable
simultaneously enforces input/output Asset equality and commitment
correctness;
\item an auxiliary boolean $\mathsf{is\_native\_asset}$ is constrained to
equal $1$ exactly when $\mathsf{AssetBase} = V^{\mathsf{Orchard}}$;
\item $\mathsf{enableZSA} = 0$ forces $\mathsf{is\_native\_asset} = 1$;
\item the value commitment's fixed-base multiplication (A4) becomes a
variable-base multiplication by the witnessed Asset Base;
\item a boolean $\mathsf{split\_flag}$ replaces the spent value inside
the value commitment: the committed net value becomes
$v' - v_{\mathrm{new}}$, with $v' = 0$ if $\mathsf{split\_flag} = 1$ and
$v' = v_{\mathrm{old}}$ otherwise; $\mathsf{split\_flag} = 1$ forces
$\mathsf{is\_native\_asset} = 0$;
\item the Merkle path constraint (A3) becomes
$(v_{\mathrm{old}} = 0 \wedge \mathsf{is\_native\_asset} = 1) \vee
\mathrm{ValidMerklePath}$ --- the dummy skip survives only for the native
Asset;
\item nullifier integrity is randomized for split notes
(\S\ref{sec:zsa-split}).
\end{enumerate}

The circuit realizes these deltas in a single gate, \texttt{q\_orchard},
whose twelve named constraints are stated once, in
\S\ref{sec:zsa-split-gate}, in the code's names (the code names the ZIP's
$\mathsf{is\_native\_asset}$ as $\mathsf{is\_zatoshi\_asset}$); we do not
restate them here.

Two gadgets carry the arithmetic weight. The in-circuit value commitment
computes $[\mathrm{magnitude}]\,\mathsf{asset\_base}$ by variable-base long
scalar multiplication on the witnessed non-identity point, applies
$\mathsf{mul\_sign}$ for the sign, and adds the fixed-base
$[\rcv]\,R^{\mathsf{Orchard}}$ (\texttt{src/circuit/value\_commit\_orchard.rs}).
The in-circuit note commitment shares one Sinsemilla evaluation between the
two branches of Definition~\ref{def:zsa-notecommit}: it muxes the initial
point $Q$ between $Q_{\mathrm{ZEC}}$ (domain
\texttt{z.cash:Orchard-NoteCommit-M}) and $Q_{\mathrm{ZSA}}$ (domain
\texttt{z.cash:ZSA-NoteCommit-M}; the hard-coded point is test-asserted in
\texttt{src/constants/sinsemilla.rs}) on $\mathsf{is\_zatoshi\_asset}$,
hashes the common message prefix once, hashes the two suffixes (ZEC: the
zero-padded piece $h_2$; ZSA: the Asset Base bits), muxes the resulting hash
point, and adds the shared blinding term $[\rcm]\,R$ --- the shared
randomness base of Definition~\ref{def:zsa-notecommit} is what makes this
sharing possible (\texttt{src/circuit/note\_commit.rs}). The Asset Base
message pieces: $h_{\mathrm{ZSA}} = (\text{bits } 249\text{--}253 \text{ of }
\psi) \,\|\, (\text{bit } 254 \text{ of } \psi) \,\|\, (\text{bits }
0\text{--}3 \text{ of } x(\mathsf{asset}))$; $i = \text{bits }
4\text{--}253$ of $x(\mathsf{asset})$;
$j = (\text{bit } 254 \text{ of } x(\mathsf{asset})) \,\|\,
(\tilde{y}(\mathsf{asset})) \,\|\, 8 \text{ zero bits}$ (ibid.).

The public instance grows from $9$ to $10$ field elements:
\[
\begin{aligned}
\bigl(&\mathsf{anchor},\ \cvsym^{\mathsf{net}}\!.x,\ \cvsym^{\mathsf{net}}\!.y,
\ \nfsym_{\mathrm{old}},\ \rk.x,\ \rk.y,\ \cmx, \\
&\mathsf{enable\_spends},\ \mathsf{enable\_outputs},
  \ \mathsf{enable\_zsa}\bigr),
\end{aligned}
\]
with the new $\mathsf{enable\_zsa}$ last. In the vanilla flavor
$\mathsf{enable\_zsa}$ is always false, and since \textsf{halo2} pads
instances with zeros, old nine-element proofs and new proofs behave
identically; the crate accordingly carries two circuit flavors,
\texttt{OrchardVanilla} and \texttt{OrchardZSA} (\textsf{orchard-zsa},
\texttt{src/circuit.rs}, \texttt{src/flavor.rs}). At the bundle level the
$\mathsf{enableZSAs}$ flag is bit $2$ of $\mathsf{flagsOrchard}$ (LSB order,
after $\mathsf{enableSpendsOrchard}$ and $\mathsf{enableOutputsOrchard}$;
the implementation constant is \texttt{FLAG\_ZSA\_ENABLED} $=$
\texttt{0b0000\_0100}), and coinbase transactions MUST set
$\mathsf{enableSpendsOrchard}$ and $\mathsf{enableZSAs}$ to $0$ (withdrawn
ZIP-230; \textsf{orchard-zsa}, \texttt{src/bundle.rs}).

\subsection{The carrier: transaction format, digests, and fees}
\label{sec:zsa-carrier}

The byte-level carrier has no merged deployment path upstream. Upstream
ZIP-230 (the OrchardZSA v6 transaction format) is Status: Withdrawn, as is
upstream ZIP-246 (its digests), and the new ZIP-229 (Version 6 Transaction
Format; Draft, created 2026-06-13) states as a non-requirement that ``the
v6 transaction format need not support ZSAs''. The QEDit fork retains
ZIP-230 and ZIP-246 as Draft. We therefore document the carrier from the
withdrawn upstream texts, the fork, and the implementation, and bin it
accordingly: the transfer/burn cryptography above is Draft-specified, while
the encoding it rides on is \emph{designed and implemented but not
specified upstream} (\S\ref{sec:zsa-scope}).

\paragraph{Note plaintext.} The OrchardZSA note plaintext adds
$\mathsf{asset\_base}$ ($32$ bytes) after $\mathsf{rseed}$. Per withdrawn
ZIP-230 the v6 Orchard note plaintext is
\[
\mathsf{np} := \bigl[\mathsf{leadByte}\bigr] \,\|\,
\mathrm{LEBS2OSP}_{88}(d) \,\|\, \mathrm{I2LEOSP}_{64}(v) \,\|\,
\mathsf{rseed} \,\|\, \mathsf{asset\_base} \,\|\, K^{\mathsf{memo}},
\]
with a $32$-byte memo key $K^{\mathsf{memo}}$ replacing the $512$-byte memo
field, giving $\mathsf{encCiphertext}$ of $132$ bytes and an
$\mathsf{OrchardZSAAction}$ of $372$ bytes ($\cvsym$ $32$ + nullifier $32$ +
$\rk$ $32$ + $\cmx$ $32$ + $\mathsf{ephemeralKey}$ $32$ +
$\mathsf{encCiphertext}$ $132$ + $\mathsf{outCiphertext}$ $80$). ZIP-230
notes that an Orchard note ``is essentially equivalent to an OrchardZSA
note with $\mathsf{AssetBase} = V^{\mathsf{Orchard}}$'' and that the two
need not be distinguished in plaintexts.

\begin{remark}[The lead byte is unassigned upstream]\label{rem:zsa-leadbyte}
Upstream ZIP-226 requires $\mathsf{leadByte}$ to be the unresolved
placeholder \texttt{\{\{ZSALEADBYTE\}\}} when the protocol's note-creation
sections (\S4.7.3 ``Sending Notes'', \S4.8.3 ``Dummy Notes'') are invoked
in the computation of $\rho$ and $\psi$ for OrchardZSA notes --- a
sentence with no counterpart
in the fork's ZIP-226 and no implementation counterpart. Upstream ZIP-230
replaced its original $\mathtt{0x03}$ with the placeholder
\texttt{\{\{LEADBYTE\}\}} because ZIP-2005 separately claims lead byte
$\mathtt{0x03}$. The QEDit fork ZIPs still say $\mathtt{0x03}$, and the
implementation hard-codes $\mathsf{NOTE\_VERSION\_BYTE\_V3} = \mathtt{0x03}$
with $\mathsf{MEMO\_SIZE} = 512$ and no $K^{\mathsf{memo}}$ --- the
pre-memo-bundle layout, giving $\mathsf{COMPACT\_NOTE\_SIZE\_ZSA} = 52 + 32
= 84$ and a $612$-byte $\mathsf{encCiphertext}$, diverging from ZIP-230's
$132$ bytes (\textsf{orchard-zsa},
\texttt{src/primitives/zcash\_note\_encryption\_domain.rs},
\texttt{src/primitives/orchard\_primitives.rs}). No lead byte is currently
reserved for ZSA notes upstream. Bin: \emph{open}.
\end{remark}

\paragraph{Digests.} Withdrawn ZIP-246 extends the ZIP-244 digest tree ---
the baseline is the \emph{Orchard Guide}, ``Conservation of value'', its
ZIP-244 paragraph --- with a per-Action-Group
$\mathsf{orchard\_burn\_digest}$ (node T.4a.vi) under
$\mathsf{orchard\_action\_groups\_digest}$ (personalization
\texttt{ZTxIdOrcActGHash}), beside the compact and non-compact action
digests, $\mathsf{flagsOrchard}$, $\mathsf{anchorOrchard}$, and
$\mathsf{nAGExpiryHeight}$. The compact digest covers nullifier, $\cmx$,
$\mathsf{ephemeralKey}$, and $\mathsf{encCiphertext}$ (personalization
\texttt{ZTxId6OActC\_Hash}); the non-compact digest covers $\cvsym$, $\rk$,
and $\mathsf{outCiphertext}$ (\texttt{ZTxId6OActN\_Hash}). The burn digest
is BLAKE2b-256 with personalization \texttt{ZTxIdOrcBurnHash} over, per
$\mathsf{assetBurn}$ tuple,
\[
\mathsf{assetBase} \,\|\, \mathsf{valueBurn} \quad
(\text{$64$-bit unsigned little-endian}),
\]
the empty burn set hashing to
$\mathrm{BLAKE2b\text{-}256}(\texttt{ZTxIdOrcBurnHash}, [\,])$ (ZIP-246).
One disambiguation: the header-digest field T.1g $\mathsf{burn\_amount}$ is
ZIP-233's network-sustainability burn, unrelated to asset burn (ZIP-246).

\paragraph{Parsing and bundle types.} Function \texttt{read\_v6\_bundle}
(gated by \texttt{zcash\_unstable = "nu7"}) reads
$\mathsf{nActionGroups}$, which must be $0$ or exactly $1$ (``A V6
transaction must contain at most one action group''), then the actions,
flags, anchor, $\mathsf{nAGExpiryHeight}$ (must be $0$), the burn vector
($32$-byte Asset Base plus note value), the proof, per-action versioned
signatures, $\mathsf{valueBalance}$, and a versioned binding signature; the
write path is symmetric via \texttt{write\_burn}
(\textsf{librustzcash-zsa},
\texttt{zcash\_primitives/src/transaction/components/orchard.rs}). The
bundle type carries the burn set as the field \texttt{burn:}
\texttt{Vec<(AssetBase, NoteValue)>} beside \texttt{value\_balance}, and
\texttt{binding\_validating\_key()} returns \texttt{derive\_bvk(actions,}
\texttt{value\_balance, burn)}; the builder collects burns in a
\texttt{BTreeMap<AssetBase, NoteValue>}, canonically ordered and
duplicate-free by construction (\textsf{orchard-zsa},
\texttt{src/bundle.rs}, \texttt{src/builder.rs}). One open seam: ZIP-230
permits $\mathsf{nActionGroupsOrchard} > 1$ with per-group burn fields (the
multi-party rationale of \S\ref{sec:zsa-burn}), the implementation rejects
more than one group, and ZIP-226's $\mathsf{bvk}$ equation is written over
a single flat action/burn set; how multi-group aggregation would interact
with the binding signature is unexercised. Bin: \emph{open}.

\paragraph{Tree ordering.} Per Action Group, every new OrchardZSA note
commitment is appended to the note commitment tree per Action; afterwards,
Issue Note commitments are appended per Issue Action (ZIP-227, ``Addition
to the Note Commitment Tree''; \S\ref{sec:zsa-issuance}).

\paragraph{Fees.} The QEDit fork of ZIP-317 adds a
$\mathsf{contribution\_ZSAIssuance}$ term --- issue notes and asset
creations --- to the conventional-fee formula; it adds no burn term, so a
burn contributes to fees only through the Actions that fund it. Upstream
ZIP-317 contains no ZSA changes. The conventional-fee baseline is the
\emph{Wallet Guide}, ``Fees (ZIP-317)''.

%% ZSA Guide --- draft section: Split notes and the counterfeiting boundary
%% Ground truth verified at: zcash/zips main @ 69610984 (2026-07-09); QEDit
%% forks: orchard-zsa branch zsa1 @ cf801a5d ("Merge upstream 20260609
%% (#251)", 2026-06-29), zips-zsa zsa1 @ fd71419a, librustzcash-zsa zsa1 @
%% 3f9b53fc. Deployment status (Draft ZIPs, testnet, NU7-contested) is
%% stated once in this volume's scope section and not restated here.

\section{Split notes and the counterfeiting boundary}\label{sec:zsa-split}

An Orchard bundle pads to a uniform shape with \emph{dummy notes}: an Action
with no real input spends a zero-valued note under a freshly sampled key,
and the membership constraint waves it through because the check is gated by
the spent value --- constraint A3 of the Action statement,
$v_{\mathrm{old}} \cdot (\mathsf{root} - \mathsf{anchor}) = 0$, ``how an
Action pads with no real input'' (the \emph{Orchard Guide}, ``The Action
statement, formally''; its node-verification section, ``What a node
verifies, without ever observing the note'', records that zero-value dummy
spends are accepted by design). The construction is protocol \S4.8.3,
implemented in \texttt{orchard-zsa}, \texttt{src/builder.rs},
\texttt{SpendInfo::dummy}: a random spending key, a zero-valued note, and a
random Merkle path that is never checked. The exemption is sound there
because every Orchard value commitment rides the \emph{fixed} base
$V^{\mathsf{Orchard}}$ (\emph{Orchard Guide}, ``Conservation of value:
value commitments and the binding signature''): a dummy spend can move no
value, whatever path it fails to prove. ZIP-226 replaces the fixed base
with a witnessed one, and the same exemption becomes a mint. This section
develops the attack that forces the change, the \emph{split note} that
ZIP-226 substitutes for the dummy note on custom assets, and the exact
constraint set --- specified and implemented --- that draws the
counterfeiting boundary.

Ground truth for this section: ZIP-226 and ZIP-227 at \texttt{zcash/zips}
commit \texttt{69610984} (2026-07-09); the pinned implementation, crate
\texttt{orchard} of the QED-it fork (\texttt{orchard-zsa}, branch
\texttt{zsa1}) at commit \texttt{cf801a5d} (2026-06-29); the QED-it
\texttt{zips} fork at \texttt{fd71419a} where it diverges from upstream.
Split-note logic lives entirely in the \texttt{orchard} crate; the
\texttt{librustzcash} fork consumes its bundle API and contains no
split-note code of its own (verified by search at \texttt{3f9b53fc}).

\subsection{The counterfeit construction the boundary excludes}
\label{sec:zsa-split-attack}

ZIP-226 generalises the net value commitment from the fixed base
$V^{\mathsf{Orchard}}$ to a per-asset base (ZIP-226, ``Value Commitment''):
\begin{equation}\label{eq:zsa-cvnet}
\begin{split}
\cvsym^{\mathsf{net}} &:=
\mathsf{ValueCommit}^{\mathsf{OrchardZSA}}_{\rcv}\!\bigl(
\mathsf{AssetBase}_{\mathsf{AssetId}},\,
v^{\mathsf{net}}_{\mathsf{AssetId}}\bigr) \\
&\;= [\,v^{\mathsf{net}}_{\mathsf{AssetId}}\,]\,
\mathsf{AssetBase}_{\mathsf{AssetId}} + [\rcv]\,R^{\mathsf{Orchard}},
\end{split}
\end{equation}
with $v^{\mathsf{net}} := v_{\mathrm{old}} - v_{\mathrm{new}}$ and
$R^{\mathsf{Orchard}}$ the randomness base of the \emph{Orchard Guide}'s
net value commitment (``Conservation of value: value commitments and the
binding signature''), unchanged. The native asset keeps the old base:
``the Asset Base of the ZEC Asset will be kept as the original value base
point $\mathcal{V}^{\mathsf{Orchard}}$'' (ZIP-226, ``Asset Identifiers'';
in code \texttt{AssetBase::zatoshi()} is
$\GroupHash^{\Pallas}(\mathtt{z.cash:Orchard\text{-}cv}, \texttt{"v"})$,
\texttt{orchard-zsa}, \texttt{src/note/asset\_base.rs},
\texttt{src/constants/fixed\_bases.rs}). Verification subtracts the
declared native balance and the declared burns (ZIP-226, ``Value Balance
Verification''; implemented in \texttt{src/bundle.rs},
\texttt{derive\_bvk\_raw}):
\begin{equation}\label{eq:zsa-bvk}
\begin{split}
\mathsf{bvk} = \Bigl(\sum_{i=1}^{n} \cvsym^{\mathsf{net}}_{i}\Bigr)
&- \mathsf{ValueCommit}^{\mathsf{OrchardZSA}}_{0}\!\bigl(
V^{\mathsf{Orchard}},\, v^{\mathsf{balanceOrchard}}\bigr) \\
&- \sum_{(\mathsf{AssetBase},\, v)\, \in\, \mathsf{assetBurn}}
\mathsf{ValueCommit}^{\mathsf{OrchardZSA}}_{0}(\mathsf{AssetBase},\, v),
\end{split}
\end{equation}
and the binding signature under $\mathsf{bvk}$ proves knowledge of
$\mathsf{bsk} = \sum_i \rcv_i$ exactly as in Orchard.

Suppose now that the dummy exemption survived unchanged: a padding Action
proves no membership, and --- since the base of \eqref{eq:zsa-cvnet} is a
private witness --- may witness \emph{any} point $A$ as its asset base.
Choose $A := [-1]\,V^{\mathsf{Orchard}}$, witness $v_{\mathrm{old}} = 0$
and $v_{\mathrm{new}} = v$:
\[
\cvsym^{\mathsf{net}}_{1} = [\,0 - v\,]\,A + [\rcv_1]\,R^{\mathsf{Orchard}}
= [\,v\,]\,V^{\mathsf{Orchard}} + [\rcv_1]\,R^{\mathsf{Orchard}}.
\]
Pair it with a second Action that outputs a genuine $v$-valued ZEC note,
\[
\cvsym^{\mathsf{net}}_{2} = [\,0 - v\,]\,V^{\mathsf{Orchard}}
+ [\rcv_2]\,R^{\mathsf{Orchard}},
\]
and declare $v^{\mathsf{balanceOrchard}} = 0$ with no burns. The
$V^{\mathsf{Orchard}}$-components cancel in \eqref{eq:zsa-bvk},
\[
\mathsf{bvk} = \cvsym^{\mathsf{net}}_{1} + \cvsym^{\mathsf{net}}_{2}
= [\rcv_1 + \rcv_2]\,R^{\mathsf{Orchard}},
\]
the attacker knows $\mathsf{bsk} = \rcv_1 + \rcv_2$, the binding signature
verifies --- and a $v$-valued ZEC note now exists that no one paid for. Any
known multiple or linear combination of existing bases serves equally.
(\emph{Classification:} the construction is assembled here from the
specified formulas \eqref{eq:zsa-cvnet} and \eqref{eq:zsa-bvk}; ZIP-226
states the attack only in words.) The ZIP's own statement, verbatim: the
output note of split Actions ``cannot contain just any Asset Base. We must
enforce it to be an actual output of a GroupHash computation\dots\ Without
this enforcement the prover could input a multiple (or linear combination)
of an existing Asset Base, and thereby attack the network by overflowing
the ZEC value balance and hence counterfeiting ZEC funds'' (ZIP-226,
``Split Notes'', rationale).

The boundary to draw is therefore precise: every asset base that enters a
value commitment must be proven to be an output of $\GroupHash$ --- a point
of unknown discrete-log relation to $V^{\mathsf{Orchard}}$ and
$R^{\mathsf{Orchard}}$ --- and the one place an unproven base could enter
is an Action with no real input.

\subsection{The split input}\label{sec:zsa-split-input}

\begin{definition}[Split Input, Split Action]\label{def:zsa-split-input}
ZIP-226 defines a \emph{Split Input} as ``an Action input used to ensure
that the output note of that Action is of a validly issued AssetBase when
there is no corresponding real input note, in situations where the number
of outputs are larger than the number of inputs'', and a \emph{Split
Action} as an Action containing a Split Input (ZIP-226, ``Terminology'').
\end{definition}

Mechanically, a Split Input is a \emph{copy of a previously issued note} ---
one whose commitment is already in the note commitment tree --- with
exactly two changes: a boolean witness $\mathsf{split\_flag}$ is set to
$1$, and the note's value is replaced by $0$ \emph{in the value-commitment
computation only} (ZIP-226, ``Split Notes''). The copied note's commitment
$\cmsym_{\mathrm{old}}$ still opens with the true value: in the circuit the
value is zeroed solely in the $\cvsym^{\mathsf{net}}$ equation, while the
old-note commitment check consumes $v_{\mathrm{old}}$ unchanged
(\texttt{orchard-zsa}, \texttt{src/circuit/circuit\_zsa.rs}: the gate
zeroes $v_{\mathrm{old}}$ via the $(1-\mathsf{split\_flag})$ factor in the
value equation, lines 127--131, whereas the $\NoteCommit$ call at lines
637--655 takes $v_{\mathrm{old}}$ as witnessed). The copy is therefore not
a spend --- it moves no value and, as \S\ref{sec:zsa-split-nf} shows,
reveals no usable nullifier --- but it drags the original note's asset base
into the proof, where an equality constraint passes it to the output note.

The split Action still balances, per asset. Committed values must sum to
zero for each asset base separately, ``as the Pedersen commitments add up
homomorphically only with respect to the same value base point'' (ZIP-226,
``Value Balance Verification''): for a correct transaction the
custom-asset commitments satisfy
\[
\sum_{j \in S_{\mathsf{CA}}} \cvsym^{\mathsf{net}}_{j}
- \sum_{(\mathsf{AssetBase},\,v)\,\in\,\mathsf{assetBurn}}
[\,v\,]\,\mathsf{AssetBase}
= \sum_{j \in S_{\mathsf{CA}}} [\rcv_j]\,R^{\mathsf{Orchard}},
\]
so the split Action's $\cvsym^{\mathsf{net}} = [\,0 -
v_{\mathrm{new}}\,]\,\mathsf{AssetBase} + [\rcv]\,R^{\mathsf{Orchard}}$
cancels against the real spending Actions of the same asset. The builder
implements exactly this: \texttt{ActionInfo::value\_sum} uses a spent value
of $0$ when $\mathsf{split\_flag}$ is set (\texttt{orchard-zsa},
\texttt{src/builder.rs}, lines 495--506), and \texttt{derive\_bvk\_raw}
computes \eqref{eq:zsa-bvk} over the resulting commitments
(\texttt{src/bundle.rs}, lines 482--517).

For the note to copy, ZIP-226 offers the wallet three options: (1) another
note of the same Asset Base being spent in the same transaction; (2) a
different unspent note of that Asset Base, which is \emph{not} thereby
spent; (3) an already spent note of that Asset Base --- the zeroed value
prevents double spending (ZIP-226, ``Split Notes''). The ZIP adds that
``we do not care about whether the previously issued note copied to create
a Split Input is owned by the sender, or whether it was nullified
before''; \S\ref{sec:zsa-split-ownership} examines how far that sentence
actually holds against the constraint system.

\subsection{Membership as well-formedness: the issuance anchor}
\label{sec:zsa-split-membership}

The fix ZIP-226 specifies, verbatim: ``for Custom Assets we enforce that
every input note to an OrchardZSA Action must be proven to exist in the
set of note commitments in the note commitment tree. We then enforce this
real note to be unspendable in the sense that its value will be zeroed in
split Actions and the nullifier will be randomized\dots\ the proof itself
ensures that the output note is of the same Asset Base as the input note.
In the circuit, the split note functionality will be activated by a
boolean private input (aka the split\_flag boolean)'' (ZIP-226, ``Split
Notes'', rationale).

Why does tree membership prove that an asset base is a $\GroupHash$
output? Because for custom assets the tree is fed exclusively from two
sources that both enforce it. Transfer outputs inherit their base from a
proven input by the in-circuit equality constraint
(\S\ref{sec:zsa-split-circuit}). Issuance outputs are checked publicly: a
consensus rule obliges validators to recompute each Issue Note's
$\mathsf{AssetBase}$ from the issuance bundle's issuer key and the
IssueAction's $\mathsf{assetDescHash}$ (ZIP-227, ``Specification:
Consensus Rule Changes''), along the derivation chain constructed in
\S\ref{sec:zsa-asset-base} (ZIP-227, ``Asset Bases''; recomputed in
\texttt{orchard-zsa}, \texttt{src/note/asset\_base.rs}),
and then appends the issue-note commitments to the \emph{same} Orchard
note commitment tree --- transfer-Action commitments first, then issuance
notes (ZIP-227, ``Addition to the Note Commitment Tree''). Hence a
custom-asset note commitment resident in the tree commits, by induction
along the two rules, to a genuine $\GroupHash$ output. ZIP-226 states the
conclusion: ``This ensures that the value base points of all output notes
of a transfer are actual outputs of a GroupHash, as they originate in the
Issuance protocol which is publicly verified'' (``Split Notes'',
rationale); and, on the identifier side, ``We prevent a potential
malleability attack on the Asset Identifier by ensuring the output notes
receive an Asset Base that exists on the global state'' (``Rationale for
Asset Identifiers'').

One clarification the ZIP's sentence invites: ``every input note to an
OrchardZSA Action'' is scoped to \emph{custom-asset} inputs. Native-asset
(ZEC) Actions in a v6 bundle retain the dummy-note exemption unchanged ---
``The ZEC-based Actions will still include dummy input notes, whereas the
OrchardZSA Actions will include split input notes and will not include
dummy input notes'' (ZIP-226, ``Backward Compatibility''; the ``Split
Notes'' rationale states the same). The constraint-level form of this
scoping is the disjunction in \S\ref{sec:zsa-split-circuit}.

\subsection{The specified circuit delta}\label{sec:zsa-split-circuit}

We state ZIP-226's circuit changes (``Circuit Statement'') as a delta
against the Action statement A1--A9 of the \emph{Orchard Guide} (``The
Action statement, formally''). Two new boolean witnesses appear:
$\mathsf{split\_flag}$, and $\mathsf{is\_native\_asset}$, constrained to
equal $1$ exactly when the witnessed base is the native one,
$\mathsf{AssetBase} = V^{\mathsf{Orchard}}$. One public input appears: the
bundle-level $\mathsf{enableZSA}$ flag (a bit of the v6
$\mathsf{flagsOrchardZSA}$ field), with the constraint
$\mathsf{enableZSA} = 0 \implies \mathsf{is\_native\_asset} = 1$ ---
switching the flag off collapses the circuit to native-asset behaviour
and, transitively via the split constraint below, disables split notes
(ZIP-226, ``Overview'' and ``Circuit Statement'', The enableZSA Flag; in
code the flag is the tenth public input, taken from
\texttt{Flags::zsa\_enabled()}, \texttt{orchard-zsa},
\texttt{src/circuit.rs}, lines 476--494, 560).

\begin{itemize}[leftmargin=*]
\item \textbf{A1, A2 (note commitment integrity).} Both commitment checks
replace $\NoteCommit^{\mathsf{Orchard}}$ with
$\NoteCommit^{\mathsf{OrchardZSA}}$, which takes the witnessed
$\mathsf{AssetBase}$ as an additional final argument; the \emph{same}
once-witnessed base feeds the old-note check, the new-note check, and the
value commitment (ZIP-226, ``Circuit Statement'', Asset Base Equality).
For the native asset $\NoteCommit^{\mathsf{OrchardZSA}}$ coincides with
$\NoteCommit^{\mathsf{Orchard}}$; otherwise it is the Sinsemilla
commitment over the extended message ending in the asset-base encoding, in
the hash domain \texttt{"z.cash:ZSA-NoteCommit-M"} with the blinding base
shared with Orchard (\texttt{"z.cash:Orchard-NoteCommit-r"}) (ZIP-226,
``Note Structure and Commitment''). The full
$\NoteCommit^{\mathsf{OrchardZSA}}$ construction is
Definition~\ref{def:zsa-notecommit}; its in-circuit mux is described in
\S\ref{sec:zsa-circuit}.

\item \textbf{A3 (membership).} The value gate becomes a disjunction
(ZIP-226, ``Circuit Statement'', Asset Identifier Consistency for Split
Actions --- Merkle Path Validity):
\[
\bigl(v_{\mathrm{old}} = 0 \,\land\, \mathsf{is\_native\_asset} = 1\bigr)
\;\lor\; \bigl(\text{Valid Merkle Path}\bigr),
\]
so the dummy skip survives \emph{only} for the native asset. Every
custom-asset input --- split or not, zero-valued or not --- must prove
membership.

\item \textbf{A4 (value commitment).} Two changes. First, ``the
fixed-base multiplication constraint between the value and the value base
point of the value commitment, cv, is replaced with a variable-base
multiplication between the two'' (ZIP-226, ``Circuit Statement'', Value
Commitment Correctness), the base being the witnessed
$\mathsf{AssetBase}$. Second, the input value is replaced by a generic
$v'$ (ZIP-226, ``Circuit Statement'', Asset Identifier Consistency for
Split Actions):
\[
\cvsym^{\mathsf{net}} =
\mathsf{ValueCommit}^{\mathsf{OrchardZSA}}_{\rcv}\!\bigl(
\mathsf{AssetBase},\, v' - v_{\mathrm{new}}\bigr), \qquad
v' = \begin{cases} 0 & \text{if } \mathsf{split\_flag} = 1,\\
v_{\mathrm{old}} & \text{otherwise,}\end{cases}
\]
together with
$\mathsf{split\_flag} = 1 \implies \mathsf{is\_native\_asset} = 0$: split
notes exist only for custom assets, and a zero-valued native dummy cannot
masquerade as one.

\item \textbf{A5 (nullifier).} The derivation changes for split notes
``to prevent the identification of notes'' (ZIP-226, ``Circuit
Statement'', Asset Identifier Consistency for Split Actions);
\S\ref{sec:zsa-split-nf} develops it.

\item \textbf{A6, A7 (spend authority, address integrity).} Untouched by
the ZIP's change list --- and, decisively for
\S\ref{sec:zsa-split-ownership}, they remain \emph{unconditional}: no
$\mathsf{split\_flag}$ factor relaxes them.
\end{itemize}

\subsection{The implemented gate}\label{sec:zsa-split-gate}

A naming note, once: the ZIP's $\mathsf{is\_native\_asset}$ is
\texttt{is\_zatoshi\_asset} throughout the implementation, assigned $1$
exactly when the witnessed asset equals \texttt{AssetBase::zatoshi()}
(\texttt{orchard-zsa}, \texttt{src/circuit/gadget.rs},
\texttt{circuit\_zsa.rs}). We keep the ZIP name in prose and the code name
inside code-derived formulas.

The implemented circuit concentrates the delta in a single gate,
\texttt{q\_orchard}, which carries twelve named constraints
where the vanilla circuit's corresponding gate has four
(\texttt{src/circuit/circuit\_zsa.rs}, lines 68--185, versus
\texttt{src/circuit/circuit\_vanilla.rs}, lines 59--100). In the code's
names:
\begin{align*}
&\mathsf{bool\_check}(\mathsf{split\_flag}), \qquad
\mathsf{bool\_check}(\mathsf{is\_zatoshi\_asset}),\\
&v_{\mathrm{old}} \cdot (1 - \mathsf{split\_flag}) - v_{\mathrm{new}}
= \mathrm{magnitude} \cdot \mathrm{sign},\\
&(v_{\mathrm{old}} + 1 - \mathsf{is\_zatoshi\_asset}) \cdot
(\mathsf{root} - \mathsf{anchor}) = 0,\\
&v_{\mathrm{old}} = 0 \,\lor\, \mathsf{enable\_spends} = 1, \qquad
v_{\mathrm{new}} = 0 \,\lor\, \mathsf{enable\_outputs} = 1,\\
&\mathsf{is\_zatoshi\_asset} \cdot (\mathsf{asset}_x -
\mathsf{zatoshi}_x) = 0, \qquad
\mathsf{is\_zatoshi\_asset} \cdot (\mathsf{asset}_y -
\mathsf{zatoshi}_y) = 0,\\
&(1 - \mathsf{is\_zatoshi\_asset}) \cdot (\Delta_x \cdot
\Delta_x^{\mathsf{inv}} - 1) \cdot (\Delta_y \cdot
\Delta_y^{\mathsf{inv}} - 1) = 0,\\
&(1 - \mathsf{split\_flag}) \cdot (\psi_{\mathrm{old}} -
\psi^{\nfsym}) = 0, \qquad
\mathsf{split\_flag} \cdot \mathsf{is\_zatoshi\_asset} = 0,\\
&(1 - \mathsf{enable\_zsa}) \cdot (1 - \mathsf{is\_zatoshi\_asset}) = 0,
\end{align*}
where $\Delta_x, \Delta_y$ abbreviate the code's
\texttt{diff\_asset\_x}, \texttt{diff\_asset\_y} --- the witnessed asset's
coordinates minus the native base's --- and $\psi^{\nfsym}$ is the code's
\texttt{psi\_nf} (\S\ref{sec:zsa-split-nf}). Four points merit unpacking.

\paragraph{The folded value equation.} The single constraint
$v_{\mathrm{old}} \cdot (1 - \mathsf{split\_flag}) - v_{\mathrm{new}} =
\mathrm{magnitude} \cdot \mathrm{sign}$ folds the ZIP's case-split $v'$
into one expression; the witness generator supplies the matching net
value, $-v_{\mathrm{new}}$ for split spends and $v_{\mathrm{old}} -
v_{\mathrm{new}}$ otherwise (\texttt{circuit\_zsa.rs}, lines 462--501).

\paragraph{The Merkle factor.} The disjunction of A3 is encoded as the
single product $(v_{\mathrm{old}} + 1 - \mathsf{is\_zatoshi\_asset})
\cdot (\mathsf{root} - \mathsf{anchor}) = 0$. The code comment argues the
encoding safe from wrap-around: $\mathsf{is\_zatoshi\_asset}$ is
boolean-checked and $v_{\mathrm{old}}$ is range-checked to $64$ bits in
the note-commitment evaluation, so over $\F_{p_{\Pallas}}$ the first
factor vanishes exactly when $v_{\mathrm{old}} = 0$ and
$\mathsf{is\_zatoshi\_asset} = 1$ (\texttt{circuit\_zsa.rs}, lines
132--139). The consequence bears repeating: for every custom-asset input
the factor is nonzero, so $\mathsf{root} = \mathsf{anchor}$ is mandatory
--- a strictly stronger membership requirement than vanilla Orchard's
$v_{\mathrm{old}}$-gate, applying even to zero-valued non-split
custom-asset inputs.

\paragraph{Proving inequality.} The assignment
$\mathsf{is\_zatoshi\_asset} = 0$ must itself be sound: the circuit must
verify that the witnessed base \emph{differs} from $V^{\mathsf{Orchard}}$.
The prover witnesses field inverses $\Delta_x^{\mathsf{inv}},
\Delta_y^{\mathsf{inv}}$ (assigned $0$ when the corresponding difference
vanishes), and the triple product above forces at least one difference to
be exhibited invertible (\texttt{circuit\_zsa.rs}, lines 160--170 for the
gate, 792--841 for the assignment).

\paragraph{The variable-base value commitment.} With the base witnessed,
the vanilla circuit's fixed-base \emph{short} multiplication
$[\,v^{\mathsf{net}}\,]\,V^{\mathsf{Orchard}}$ --- whose $64$-bit range
check is implicit in the short-multiplication decomposition --- is no
longer available. The ZSA gadget makes the range check explicit: the
magnitude is decomposed against the Sinsemilla lookup table as six
$10$-bit windows plus one $4$-bit short check ($64$ bits), then promoted
to a full-width scalar (\texttt{ScalarVar::from\_base}) and fed to the
variable-base long multiplication --- the code notes it is ``optimized
for $\sim$255-bit scalar'' and leaves a TODO to specialise it for $64$
bits --- after which the sign is applied and the blind
$[\rcv]\,R^{\mathsf{Orchard}}$ added
(\texttt{src/circuit/value\_commit\_orchard.rs}, lines 38--127). Two
supporting details: the ZSA circuit uses the extended lookup
configuration \texttt{PallasLookupRangeCheck4\_5BConfig} (an additional
lookup-table column for $4$- and $5$-bit short checks), and it refuses to
synthesise under \texttt{CircuitVersion::InsecureUnanchoredBase}, the
pre-NU6.2 \texttt{halo2\_gadgets} variable-base scalar-multiplication
version whose base point was not anchored --- directly load-bearing here,
because the entire counterfeiting boundary rests on the multiplication
being by the \emph{witnessed} $\mathsf{AssetBase}$ and nothing else
(\texttt{circuit\_zsa.rs}, lines 22, 51, 190--197, 334--337;
\texttt{circuit.rs}, lines 151--185).

The flag interactions are tested asymmetrically: the three legal
combinations of $(\mathsf{is\_zatoshi\_asset}, \mathsf{split\_flag})$ ---
$(1,0)$, $(0,1)$, and $(0,0)$ --- are proven passing, while the $(1,1)$
combination is excluded at witness-generation time by the test helper's
assertion (``We cannot create a split note with a zatoshi asset''), so
the constraint $\mathsf{split\_flag} \cdot \mathsf{is\_zatoshi\_asset} =
0$ is not itself exercised by a failing-proof test
(\texttt{circuit\_zsa.rs}, lines 1174--1312). Documented cost: the ZSA
proof measures $5120$ bytes for one Action and $7392$ for two, against
the vanilla $4992$ and $7264$, at the shared circuit size $K = 11$; the
$\mathsf{enableZSA}$ bit is the tenth public input, at index $9$ (the
proof-size tests of \texttt{circuit\_zsa.rs} and
\texttt{circuit\_vanilla.rs}; \texttt{circuit.rs}, lines 68--80,
536--563).
%% TODO (open): no row/region-utilisation delta is documented anywhere ---
%% ZIP-226 delegates circuit detail to an external, unversioned Google Doc
%% ([#circuit-modifications]). Candidate worked example per CONVENTIONS
%% (script over CircuitCost::measure), to be decided at volume level.

\subsection{The split nullifier}\label{sec:zsa-split-nf}

A split input copies a note that may sit unspent in someone's wallet; its
Action must nevertheless publish \emph{some} nullifier. Publishing the
note's true nullifier would be harmful in both directions: it would mark
an unspent note spent (bricking it), and it would collide with the honest
nullifier if the same note were spent for real in the same transaction.
ZIP-226 therefore randomises the split note's nullifier (``Split
Notes''):
\begin{equation}\label{eq:zsa-split-nf}
\nfsym_{\mathrm{old}} = \Extract\bigl(\,[\,F_{\nksym}(\rho_{\mathrm{old}})
+ \psi^{\nfsym}\,]\; G^{\mathsf{nf}} \;+\; \cmsym_{\mathrm{old}} \;+\;
L^{\mathsf{Orchard}}\,\bigr),
\end{equation}
in the notation of the \emph{Orchard Guide}'s A5 (``The Action statement,
formally''; ZIP-226 writes $\mathcal{K}^{\mathsf{Orchard}}$ for
$G^{\mathsf{nf}}$ and computes the bracketed sum in the Pallas base
field). Two deltas against A5: the note's $\psi_{\mathrm{old}}$ is
replaced by a fresh $\psi^{\nfsym}$, and a fixed offset point
\[
L^{\mathsf{Orchard}} :=
\GroupHash^{\Pallas}(\mathtt{z.cash:Orchard}, \texttt{"L"})
\]
is added (ZIP-226, ``Split Notes''; the constant is pinned and
test-asserted against the $\GroupHash$ evaluation in
\texttt{orchard-zsa}, \texttt{src/constants/nullifier\_l.rs}). The gate
constraint $(1 - \mathsf{split\_flag}) \cdot (\psi_{\mathrm{old}} -
\psi^{\nfsym}) = 0$ of \S\ref{sec:zsa-split-gate} makes the substitution
split-only: non-split OrchardZSA nullifiers coincide with the Orchard
derivation, as the ZIP states (``Note Structure and Commitment'': the
non-split nullifier ``is generated in the same manner as in the Orchard
protocol'').

In circuit, the nullifier gadget computes the standard point
$\cmsym_{\mathrm{old}} + [\,F_{\nksym}(\rho_{\mathrm{old}}) +
\psi^{\nfsym}\,]\,G^{\mathsf{nf}}$ from the witnessed $\psi^{\nfsym}$,
forms the offset variant by adding the constant $L^{\mathsf{Orchard}}$,
and muxes the two on $\mathsf{split\_flag}$ before extraction
(\texttt{src/circuit/derive\_nullifier.rs}, lines 54--111). Out of
circuit, \texttt{Nullifier::derive} mirrors the computation and selects
the offset variant in constant time exactly when the note carries a
\texttt{rseed\_split\_note} (\texttt{src/note/nullifier.rs}, lines
73--90; \texttt{src/note.rs}, lines 415--426).

\paragraph{Deriving $\psi^{\nfsym}$: a three-way divergence.} The value
of $\psi^{\nfsym}$ is a private witness, constrained by the circuit only
in the \emph{non}-split case; how a wallet derives it for split notes is
therefore wallet-level, not consensus-level. The three extant texts
disagree; per the volume's policy (\S\ref{sec:zsa-sources}), the upstream
ZIP formula comes first. Upstream ZIP-226 --- written against the
ZIP-2005 baseline, ``expected to deploy before any ZSA activation'' ---
specifies
\[
\psi^{\nfsym} :=
\mathrm{ToBase}\bigl(\mathsf{PRFexpand}_{\mathsf{rseed\_nf}}
(\mathtt{0x0A} \,\|\, \rho_{\mathrm{old}})\bigr)
\]
with $\mathsf{rseed\_nf}$ uniform on $32$ bytes and the domain byte
$\mathtt{0x0A}$ ``reserved by ZIP 2005 for this use''
(\texttt{zip-0226.rst}:297 @ \texttt{69610984}). The QED-it \texttt{zips}
fork specifies neither derivation: there $\psi^{\nfsym}$ ``is sampled
uniformly at random'' on the Pallas base field (fork
\texttt{zip-0226.rst}:293 @ \texttt{fd71419a}). The implementation
derives
\[
\psi^{\nfsym} :=
\mathrm{ToBase}\bigl(\mathsf{PRFexpand}_{\mathsf{rseed\_split\_note}}
(\mathtt{0x09} \,\|\, \rho_{\mathrm{old}})\bigr),
\]
the ordinary $\psi$ domain byte keyed by a fresh random $32$-byte
\texttt{rseed\_split\_note} (\texttt{orchard-zsa},
\texttt{src/circuit.rs}:316--319, \texttt{src/note.rs}:114--116 and
417--425; domain byte \texttt{PSI} $= \mathtt{0x09}$ per the QED-it
\texttt{zcash\_spec} fork, \texttt{src/prf\_expand.rs}:89 @
\texttt{d5e84264}). All three yield a uniformly distributed witness; the
implementation matches neither ZIP text exactly. (\emph{Classification:}
each derivation is specified in its own text; the divergence is
consensus-invisible.)

\paragraph{Why the offset $L^{\mathsf{Orchard}}$?} The randomised
$\psi^{\nfsym}$ already makes the split nullifier unlinkable to the
note's true nullifier. ZIP-226 states no rationale for additionally
adding $L^{\mathsf{Orchard}}$. A plausible purpose is robustness against
a prover who \emph{chooses} $\psi^{\nfsym} = \psi_{\mathrm{old}}$ --- the
circuit does not forbid it in the split case --- and would thereby
publish the note's true nullifier from a zero-valued split Action,
bricking the original note or colliding with an honest spend of it
elsewhere in the transaction; the constant offset makes even that choice
yield a different group element. (\emph{Classification:}
designed-but-unspecified; the preceding sentence is our inference, stated
in neither ZIP nor code.)
%% TODO: hunt for a citable rationale (QED-it security analysis,
%% zcash/zips PR discussion) before this survives adversarial review.

\subsection{Ownership, in fact and in prose}\label{sec:zsa-split-ownership}

Recall the ZIP's claim that ``we do not care about whether the previously
issued note copied to create a Split Input is owned by the sender''
(\S\ref{sec:zsa-split-input}). The constraint system says otherwise.
Spend authority (A6) and address integrity (A7) are unchanged \emph{and
unconditional} in the ZSA circuit: the proof establishes
$\rk = [\alpha]\,G^{\mathsf{SpendAuth}} + \aksym$ and
$\ivk = \CommitIvk_{\rivk}(\Extract(\aksym), \nksym)$ with
$\pkd_{\mathrm{old}} = [\ivk]\,\gd_{\mathrm{old}}$ --- for split inputs
too (\texttt{orchard-zsa}, \texttt{src/circuit/circuit\_zsa.rs}, lines
554--623; the old-note commitment equality at lines 625--659). Because
the copied note's $\cmsym_{\mathrm{old}}$ must open to the actual tree
leaf (Sinsemilla binding plus the mandatory Merkle check of
\S\ref{sec:zsa-split-gate}), the witnessed $\pkd_{\mathrm{old}}$ is the
original recipient's, and A7 then forces the prover to know $(\aksym,
\nksym, \rivk)$ consistent with it --- that is, to hold the copied note's
full viewing key. The ZIP's ownership sentence holds for \emph{spentness}
--- a nullified note is a perfectly good split input, precisely because
the value is zeroed --- but not for ownership in the key-material sense.
The implementation never exercises the freedom the prose suggests: the
builder only ever copies one of the sender's own spends in the same
bundle, option (1) of \S\ref{sec:zsa-split-input}
(\texttt{src/builder.rs}, \texttt{pad\_spend}). (\emph{Classification:}
the constraint reading is specified --- the ZIP's constraint list and the
implemented circuit agree; the rationale prose overstates.)

One genuine exception exists, and neither ZIP states it. ZIP-227 requires
the first note of any asset's first issuance to be a \emph{reference
note}: value $0$, recipient the default diversified address of the
all-zero Orchard spending key --- a key anyone can derive (ZIP-227,
``Reference Notes''; \texttt{orchard-zsa},
\texttt{src/constants/reference\_keys.rs} pins the all-zero spending key
and the $43$-byte raw address). Because that spending key is public,
anyone can satisfy A6/A7 for the reference note, which makes it a
universally available split input for its asset --- the one case in which
``ownership does not matter'' holds literally. (\emph{Classification:}
inference; the connection between reference notes and split inputs is
stated in neither ZIP nor the implementation.)

\subsection{Builder mechanics: padding, per asset}
\label{sec:zsa-split-builder}

The two padding regimes coexist per ZIP-226's rule quoted in
\S\ref{sec:zsa-split-membership}: dummy inputs for ZEC Actions, split
inputs for custom-asset Actions. The builder partitions spends and
outputs by asset, pads each side of each partition to equal length,
shuffles within each asset, then shuffles the assembled Actions across
assets (\texttt{orchard-zsa}, \texttt{src/builder.rs}, lines 1081--1164).
The padding spend is asset-dependent (\texttt{pad\_spend}, lines
928--944): for the native asset, \texttt{SpendInfo::dummy} exactly as in
Orchard; for a custom asset, \texttt{create\_split\_spend} on the first
real spend of that asset in the bundle --- and if the bundle spends none,
\texttt{BuildError::NoSplitNoteAvailable} (``No split note has been
provided for this asset''): unlike a dummy, a split input cannot be
conjured from nothing. Output-side padding uses dummy outputs of the
padding asset in both regimes.

The copy itself is minimal. Function \texttt{Note::create\_split\_note}
duplicates the note and sets exactly one new field,
$\mathsf{rseed\_split\_note}$ --- a fresh random $32$-byte seed, drawn by
\texttt{RandomSeed::random} and resampled until the derived
$\mathsf{esk}$ is valid --- asserting the asset is not the native one;
\texttt{SpendInfo::create\_split\_spend} keeps the original note's full
viewing key and Merkle path and raises $\mathsf{split\_flag}$
(\texttt{src/note.rs}, lines 436--442, 95--103; \texttt{src/builder.rs},
lines 302--318). Recipient, value, asset, $\rho$, $\mathsf{rseed}$ ---
hence $\cmsym_{\mathrm{old}}$ --- are identical to the tree-resident
original; the only witness-level differences downstream are
$\mathsf{split\_flag} = 1$ and the $\psi^{\nfsym}$ drawn from the new
seed. Witness assembly routes exactly that:
\texttt{Witnesses::from\_action\_context\_unchecked} takes
$\psi^{\nfsym}$ from $\mathsf{rseed\_split\_note}$ when present and from
$\mathsf{rseed}$ otherwise --- so $\psi^{\nfsym} = \psi_{\mathrm{old}}$
for every non-split spend --- and packs $(\psi^{\nfsym}, \mathsf{asset},
\mathsf{split\_flag})$ into the ZSA witness extension that the vanilla
flavour leaves unknown (\texttt{src/circuit.rs}, lines 199--244,
300--345).

The counterfeiting boundary, restated once in full: a custom-asset output
note's base equals its Action's input-note base (the A1/A2 equality on
the shared witness); that input note is a tree leaf (the mandatory Merkle
factor); every custom-asset tree leaf descends from a publicly recomputed
$\GroupHash$ output (the issuance anchor of
\S\ref{sec:zsa-split-membership}); and the copy that makes padding
possible moves no value (the zeroed $v'$) and emits an unlinkable,
non-colliding nullifier \eqref{eq:zsa-split-nf}. Each link is enforced by
a constraint quoted above; breaking any one of them re-opens the mint of
\S\ref{sec:zsa-split-attack}.

\section{Transaction integration and outlook}\label{sec:zsa-integration}

The preceding sections construct the OrchardZSA objects --- asset bases
(\S\ref{sec:zsa-asset-identity}), notes and Actions
(\S\ref{sec:zsa-transfer}, \S\ref{sec:zsa-split}), issuance
(\S\ref{sec:zsa-issuance}), and burn (\S\ref{sec:zsa-burn}).
This section places them in a transaction: the v6 wire format (ZIP-230), the
digest tree that yields transaction identifiers, sighashes, and the
authorizing-data commitment (ZIP-246), the transaction-level consensus rules
of ZIP-226/227, the fee treatment (ZIP-317), the state of the QED-it
implementation stack at commit-pinned sources, the obligations the format
imposes on wallets, and the deployment outlook. Deployment status --- Draft
ZIPs, externally audited at circuit scope, running on QED-it's public
testnet, contested for NU7 --- is stated once in the scope section
(\S\ref{sec:zsa-status}) and not relitigated per object;
what this section adds is the classification of each format artifact, because
the ZSA transaction layer is the part of the protocol where specification and
implementation have diverged furthest. Table~\ref{tab:zsa-sources-integration}
pins the sources; all repository facts below were verified firsthand at the
stated commits, on 2026-07-14.

\begin{table}[htbp]
\centering\small
\begin{tabular}{@{}llll@{}}
\toprule
Source & Commit & Date & Role \\
\midrule
\texttt{zcash/zips} & \texttt{6961098} & 2026-07-09
  & upstream ZIP texts \\
QED-it \texttt{zips} fork & \texttt{fd71419} & 2026-02-11
  & fork ZIP texts (pre-withdrawal) \\
QED-it \texttt{orchard} fork, \texttt{zsa1} & \texttt{cf801a5d} & 2026-06-29
  & OrchardZSA bundle, circuit, issuance \\
QED-it \texttt{librustzcash} fork, \texttt{zsa1} & \texttt{3f9b53fc}
  & 2026-04-17 & v6 assembly, digests, fees \\
\texttt{zcash/orchard} @ \texttt{ltf/zsa} & \texttt{a799082e} & 2025-11-25
  & upstream tracking branch \\
\bottomrule
\end{tabular}
\caption{Pinned sources for this section. The \texttt{ltf/zsa} branch merges
ZSA-compatible Halo~2 plus PCZT external spend-authorization signatures; it
is the beginning of an upstreaming path, not a deployable stack.}
\label{tab:zsa-sources-integration}
\end{table}

\subsection{Format status: version 6 no longer means OrchardZSA}
\label{sec:zsa-format-status}

The carrier format for everything in this volume is ZIP-230, the OrchardZSA
v6 transaction format, with its digests in ZIP-246. Both are \emph{withdrawn}
upstream. ZIP-230 is retitled ``Withdrawn Version 6 Transaction Format''; its
warning states that it is obsoleted by ZIP-229 and ``will not be deployed.
Transaction version number 6 is now defined by ZIP 229.''
(\texttt{zip-0230.rst}, header and warning). ZIP-246 is likewise withdrawn
(``Digests for the Withdrawn Version 6 Transaction Format''); v6 transaction
identifiers, authorizing-data commitments, and signature digests are now
specified by ZIP-229 (\texttt{zip-0246.rst}, header). The QED-it fork copies
of both remain Status: Draft under their original titles --- the fork
(@~\texttt{fd71419}) predates the withdrawal.

ZIP-229 (Status: Draft, created 2026-06-13) redefines v6 for NU6.3,
introducing the Ironwood pool in the aftermath of the Orchard soundness
vulnerability whose mitigation deployed as NU6.2 (ZIP-257, Status: Final).
Its Non-requirements state explicitly that ``the v6 transaction format need
not support ZSAs, the \texttt{zip233Amount} field, the explicit fee field, or
per-transparent-input sighash information that appeared in the withdrawn
ZIP 230,'' and that version number~6 need not avoid collision with withdrawn
ZIP-230 or draft ZIP-248 (ZIP-229, Non-requirements). The same transaction
version number therefore now denotes two incompatible byte formats, and the
ZSA one is the orphan.

Classification, per the three-bin discipline of this series (the
\emph{Tachyon Guide}, \S``The three-bin method''): the byte-level format
below is fixed at ZIP rigor only by withdrawn upstream text and pinned
fork code --- \emph{designed but unspecified} (\S\ref{sec:zsa-method}) ---
with no deployment path under its current ZIP numbers. We document it because
it is the format the pinned implementation speaks and the reference against
which any respecification will be diffed; \S\ref{sec:zsa-outlook} states what
would have to happen for it to become normative again.

\subsection{The v6 wire format (ZIP-230)}\label{sec:zsa-v6}

\paragraph{Common and transparent fields.} The common header carries
\texttt{header} (\texttt{fOverwintered} and the version),
\texttt{nVersionGroupId}, \texttt{nConsensusBranchId}, \texttt{lock\_time},
\texttt{nExpiryHeight}, an
explicit 8-byte \texttt{fee} (\texttt{uint64}, per ZIP-2002), and an 8-byte
\texttt{zip233Amount} (\texttt{uint64}, per ZIP-233) (ZIP-230, Transaction
Format). Making the fee an explicit field --- rather than the implicit
difference between transparent inputs and outputs --- is one of the format's
two departures from v5 visible already in the header; the second is that
signatures become versioned, below. The transparent fields add
\texttt{vSighashInfo}: one \texttt{TransparentSighashInfo} (a
\texttt{compactSize} length followed by \texttt{sighashInfo} bytes) per
transparent input (ZIP-230, Transaction Format).

\paragraph{Orchard fields: Action Groups.} The Orchard component becomes a
vector of \emph{Action Groups}: \texttt{nActionGroupsOrchard},
\texttt{vActionGroupsOrchard}, then a transaction-level
\texttt{valueBalanceOrchard} (\texttt{int64}) and
\texttt{bindingSigOrchard}. The balance and binding signature are present iff
\texttt{nActionGroupsOrchard} $> 0$; an absent \texttt{valueBalanceOrchard}
means $v_{\mathrm{balance}}^{\mathsf{Orchard}} = 0$ (ZIP-230, Transaction
Format). Each \texttt{ActionGroupDescription} contains, in order
(ZIP-230, OrchardZSA Action Group Description):
\begin{itemize}[itemsep=1pt]
\item \texttt{nActionsOrchard} (\texttt{compactSize}, MUST be $> 0$) and
  \texttt{vActionsOrchard}, at $372$ bytes per \texttt{OrchardZSAAction};
\item \texttt{flagsOrchard}, LSB-first: \texttt{enableSpendsOrchard},
  \texttt{enableOutputsOrchard}, \texttt{enableZSAs}, remaining bits zero;
\item \texttt{anchorOrchard} ($32$ bytes) --- each group carries its own
  anchor;
\item \texttt{nAGExpiryHeight} (\texttt{uint32});
\item \texttt{nAssetBurn} and \texttt{vAssetBurn}, at $40$ bytes per
  \texttt{AssetBurn} entry;
\item \texttt{sizeProofsOrchard} and \texttt{proofsOrchard}, one aggregated
  proof for the group;
\item \texttt{vSpendAuthSigsOrchard}, one \texttt{OrchardSignature} per
  Action.
\end{itemize}
Field \texttt{nAGExpiryHeight} exists for forward compatibility with ZSA
atomic swaps (ZIP-228, currently a Reserved stub); for OrchardZSA it is zero
by consensus (ZIP-230, Rationale for \texttt{nAGExpiryHeight};
\S\ref{sec:zsa-tx-rules}). The same rationale section states ``In NU7,
\texttt{nExpiryHeight} MUST be set to 0'' --- read literally, a constraint on
the transaction-level ZIP-203 field that would contradict ordinary expiry
semantics; context, the ZIP-227 consensus rule, and the implementation (which
constrains only \texttt{nAGExpiryHeight}) all indicate
\texttt{nAGExpiryHeight} is meant. We flag this as an apparent editorial
defect in the withdrawn text. Burn fields sit inside each Action Group
rather than at transaction level so that, in a future multi-party Action
Group world, each party can burn assets within its own group (ZIP-230,
Rationale for including Burn fields inside OrchardZSA Action Groups).

\paragraph{Action and burn encodings.} An \texttt{OrchardZSAAction} occupies
$372$ bytes: \texttt{cv} ($32$), \texttt{nullifier} ($32$), \texttt{rk}
($32$), \texttt{cmx} ($32$), \texttt{ephemeralKey} ($32$),
\texttt{encCiphertext} ($132$), and \texttt{outCiphertext} ($80$) (ZIP-230,
OrchardZSA Action Group Description). The $132$-byte ciphertext is the
memo-bundle-era layout: the note plaintext gains the $32$-byte asset base and
replaces the in-band $512$-byte memo with a $32$-byte memo key
$K^{\mathsf{memo}}$ (ZIP-231), the memo chunks themselves moving to a
transaction-level memo bundle (fields \texttt{nMemoChunks},
\texttt{pruned}, \texttt{vMemoChunks}, present iff $\texttt{fAllPruned} =
0$) whose pruned entries are replaced by their digests (ZIP-230,
Transaction Format; \S\ref{sec:zsa-digests}). The v6 note
plaintext is
\[
\mathtt{leadByte} \,\|\, \mathrm{LEBS2OSP}_{88}(d) \,\|\,
\mathrm{I2LEOSP}_{64}(v) \,\|\, \mathsf{rseed} \,\|\,
\mathsf{asset\_base} \,\|\, K^{\mathsf{memo}},
\]
with $\mathsf{asset\_base} = \mathrm{LEBS2OSP}_{\ell_P}(\mathrm{repr}_P
(\mathsf{AssetBase}))$, and the lead byte mandatory for all v6 note
plaintexts (ZIP-230, note plaintext encoding). Upstream, the lead byte is the
unassigned placeholder \texttt{\{\{LEADBYTE\}\}}: the original assignment
$\mathtt{0x03}$ was taken by ZIP-2005, and no replacement is assigned
(\S\ref{sec:zsa-impl}). For the same $K^{\mathsf{memo}}$ reason the Sapling
\texttt{OutputDescriptionV6} shrinks \texttt{encCiphertext} to $100$ bytes
(ZIP-230, Sapling Output Description). An \texttt{AssetBurn} entry is $40$
bytes: \texttt{AssetBase} ($32$, the encoding of
$\mathsf{AssetBase}^{\mathsf{Orchard}}$) and \texttt{valueBurn}
(\texttt{uint64}), consensus-checked non-zero and less than
$\mathsf{MAX\_BURN\_VALUE}$ (ZIP-230, OrchardZSA Action Group Description;
\S\ref{sec:zsa-tx-rules}).

\paragraph{Issuance bundle.} The issuance component carries five fields:
\texttt{issuerLength} (\texttt{compactSize}), \texttt{issuer},
\texttt{nIssueActions}, \texttt{vIssueActions}, and
\texttt{issueAuthSig}. The signature is present
iff $\texttt{nIssueActions} > 0$; the first four fields are always present,
and if $\texttt{nIssueActions} = 0$ then \texttt{issuerLength} MUST be $0$
--- an empty issuance bundle is exactly two zero \texttt{compactSize}
values (ZIP-230, Transaction Format). An \texttt{IssueAction} carries
\texttt{assetDescHash} ($32$ bytes), \texttt{nNotes} and \texttt{vNotes} at
$115$ bytes per \texttt{IssueNoteDescription}, and \texttt{flagsIssuance}
with \texttt{finalize} in the LSB; \texttt{nNotes} may be zero, so a
finalize-only action is expressible. An \texttt{IssueNoteDescription} is
\texttt{recipient} ($43$ bytes, the raw Orchard payment-address encoding),
\texttt{value} (\texttt{uint64}), \texttt{rho} ($32$), \texttt{rseed} ($32$).
The \texttt{IssueAuthSignature} is \texttt{sizeSighashInfo} and
\texttt{sighashInfo}, then \texttt{sizeSignature} and a signature ``preceded
by a \texttt{0x00} byte indicating a BIP 340 signature'' (ZIP-230, issuance
field encodings) --- the same $\mathtt{0x00}$ scheme byte that prefixes the
issuer key encoding (ZIP-227, Issuance Authorization Signature Scheme).

\paragraph{Versioned signatures.} The \texttt{SaplingSignature} and
\texttt{OrchardSignature} types prepend a \texttt{compactSize}-prefixed
\texttt{sighashInfo} to the $64$-byte RedJubjub or RedPallas signature
(ZIP-230, signature encodings). ZIP-246 defines the versioning:
\[
\mathsf{sighashInfo} := [\,\mathsf{sighashVersion}\,] \,\|\,
\mathsf{associatedData},
\]
where version $0$ is the ``commit to all effecting data'' algorithm; for v6,
only version $0$ is defined, with empty \texttt{associatedData}, for all four
of Transparent, Sapling, Orchard, and Issuance (ZIP-246, Sighash versioning).
The machinery buys the ability to introduce --- and, critically, to retire
--- sighash algorithms per transaction version without a new format.

\paragraph{Coinbase.} For coinbase transactions the
\texttt{enableSpendsOrchard} and \texttt{enableZSAs} bits MUST be zero
(ZIP-230, consensus rules), and ZIP-235 requires all coinbase transactions to
be v6 from its activation, scheduled for NU7 (ZIP-230, Implications for
Mining).

\subsection{Digests: txid, sighash, and authorizing data (ZIP-246)}
\label{sec:zsa-digests}

ZIP-246 extends the ZIP-244 digest tree --- constructed for vanilla Orchard
in the \emph{Orchard Guide}, \S``Conservation of value: value commitments and
the binding signature'' --- from four top-level branches to six:
\[
\mathsf{txid} := \mathrm{BLAKE2b\text{-}256}\bigl(\texttt{ZcashTxHash\_}
\,\|\, \mathsf{branchid},\;
d_{\mathrm{hdr}} \,\|\, d_{\mathrm{transp}} \,\|\, d_{\mathrm{Sap}} \,\|\,
d_{\mathrm{Orch}} \,\|\, d_{\mathrm{issue}} \,\|\, d_{\mathrm{memo}}\bigr),
\]
the nodes T.1--T.6 of the ZIP, each itself a personalized
$\mathrm{BLAKE2b\text{-}256}$ (ZIP-246, TxId Digest). New relative to ZIP-244
are the fee and burn fields in the header digest, the Sapling output digests,
the entire Orchard Action-Group subtree, the issuance subtree, and the memo
subtree. The header digest T.1 is
\begin{align*}
d_{\mathrm{hdr}} := \mathrm{BLAKE2b\text{-}256}\bigl(
&\texttt{ZTxIdHeadersHash},\;
\mathsf{version} \,\|\, \mathsf{versionGroupId} \,\|\, \mathsf{branchid}
\,\|\, \mathsf{lockTime} \\
&\,\|\, \mathsf{expiryHeight} \,\|\,
\mathrm{LE}_{64}(\mathsf{fee}) \,\|\, \mathrm{LE}_{64}(\mathsf{burnAmount})
\bigr)
\end{align*}
(ZIP-246, T.1). The Orchard branch T.4 hashes the concatenated per-group
digests and the balance,
$d_{\mathrm{Orch}} = \mathrm{BLAKE2b\text{-}256}(\texttt{ZTxIdOrchardHash},\,
d_{\mathrm{AGs}} \,\|\,
\mathrm{LE}_{64}(v_{\mathrm{balance}}^{\mathsf{Orchard}}))$,
with a defined empty case; each group digest (T.4a) covers a compact digest
($\nfsym$, $\cmx$, \texttt{ephemeralKey}, \texttt{encCiphertext}), a
non-compact digest ($\cvsym$, $\rk$, \texttt{outCiphertext}),
\texttt{flagsOrchard}, \texttt{anchorOrchard}, \texttt{nAGExpiryHeight}, and
a burn digest over the $(\mathsf{AssetBase}, \mathsf{valueBurn})$ tuples
(ZIP-246, T.4a). The issuance branch T.5 hashes
$\texttt{issuerLength} \,\|\, \texttt{issuer}$ and a per-action digest
(per action: \texttt{assetDescHash}, a per-note digest over
$(\texttt{recipient}, \texttt{value}, \rho, \mathsf{rseed})$,
\texttt{flagsIssuance}), each level with a defined empty case (ZIP-246,
T.5). The memo branch T.6 hashes a nonce and the concatenation of per-chunk
digests; the v6 pruned memo bundle stores exactly these
\texttt{memo\_chunk\_digest} / \texttt{memo\_digest} values in place of
pruned chunks, so pruning preserves the txid (ZIP-246, T.6; ZIP-230,
Transaction Format). Table~\ref{tab:zsa-personalizations} collects the
personalization strings; the implementation defines the same constants
(\S\ref{sec:zsa-impl}).

\begin{table}[htbp]
\centering\small
\begin{tabular}{@{}lll@{}}
\toprule
Node & Personalization & Content \\
\midrule
T.1 & \texttt{ZTxIdHeadersHash} & header, incl.\ fee and burn amount \\
T.3b & \texttt{ZTxId6SOutC\_Hash}, \texttt{ZTxId6SOutN\_Hash}
  & Sapling v6 outputs, compact/non-compact \\
T.4 & \texttt{ZTxIdOrchardHash} & Action-Group digests, balance \\
T.4a & \texttt{ZTxIdOrcActGHash} & one Action Group \\
T.4a.i & \texttt{ZTxId6OActC\_Hash}
  & $\nfsym$, $\cmx$, \texttt{ephemeralKey}, \texttt{encCiphertext} \\
T.4a.ii & \texttt{ZTxId6OActN\_Hash}
  & $\cvsym$, $\rk$, \texttt{outCiphertext} \\
T.4a.vi & \texttt{ZTxIdOrcBurnHash} & burn tuples \\
T.5 & \texttt{ZTxIdSAIssueHash} & issuer, issue-actions digest \\
--- & \texttt{ZTxIdIssuActHash} & per action: desc.\ hash, notes, flags \\
--- & \texttt{ZTxIdIAcNoteHash}
  & per note: recipient, value, $\rho$, $\mathsf{rseed}$ \\
T.6 & \texttt{ZTxIdMemo\_\_\_Hash} & nonce, chunks digest \\
--- & \texttt{ZTxIdMemoCksHash}, \texttt{ZTxIdMemoCk\_Hash}
  & memo chunk digests \\
\midrule
A.1 & \texttt{ZTxAuthTransHash}
  & scripts, now incl.\ per-input \texttt{sighashInfo} \\
A.3 & \texttt{ZTxAuthOrchaHash}
  & group auth digests, \texttt{bindingSigOrchard} \\
A.3a & \texttt{ZTxAuthOrcAGHash} & \texttt{proofsOrchard}, spend-auth digest \\
A.3a.ii & \texttt{ZTxAuthOrSASHash} & \texttt{spendAuthSig} encodings \\
A.4 & \texttt{ZTxAuthZSAOrHash} & \texttt{issueAuthSig} encoding \\
\bottomrule
\end{tabular}
\caption{ZIP-246 digest personalizations (ZIP-246, TxId Digest and
Authorizing Data Commitment). Unnumbered rows are children whose labels the
ZIP text assigns inconsistently; see the closing remark of this
subsection.}
\label{tab:zsa-personalizations}
\end{table}

\paragraph{Signature digest.} The sighash reuses the tree with transparent
and Sapling branches replaced by input-specific variants:
\[
\mathsf{sighash} := \mathrm{BLAKE2b\text{-}256}\bigl(
\texttt{ZcashTxHash\_} \,\|\, \mathsf{branchid},\;
d_{\mathrm{hdr}} \,\|\, d^{S}_{\mathrm{transp}} \,\|\,
d^{S}_{\mathrm{Sap}} \,\|\, d_{\mathrm{Orch}} \,\|\, d_{\mathrm{issue}}
\,\|\, d_{\mathrm{memo}}\bigr),
\]
produced per transparent input, per Sapling input, per OrchardZSA Action,
and --- new in v6 --- per Issue Action; the Orchard, issuance, and memo
branches are identical to their txid-side counterparts, so the issuance
bundle and the memo bundle enter every sighash (ZIP-246, Signature Digest).
On this digest ZIP-227 hangs its authorization rule: the SigHash is the
\S4.10 SIGHASH transaction hash with the ZIP-226 modifications, computed
with \texttt{SIGHASH\_ALL}, and the issuance authorization signature MUST
satisfy $\mathsf{IssueAuthSig.Validate}(\mathsf{ik}, \mathsf{SigHash},
\mathtt{issueAuthSig}) = 1$, where $\mathsf{ik}$ is decoded from the
\texttt{issuer} field and its encoding is required to begin with
$\mathtt{0x00}$, the BIP-340 scheme byte (ZIP-227, Specification: Consensus
Rule Changes).

\paragraph{Authorizing-data commitment.} The commitment to authorizing data
follows the same pattern: A.1 covers transparent scripts together with the
new per-input \texttt{TransparentSighashInfo}; A.2 keeps the ZIP-244
structure with field encodings altered by \texttt{sighashInfo}; A.3 commits
to the per-group proofs and spend-authorization signatures and to
\texttt{bindingSigOrchard}; and a new A.4 commits to the
\texttt{issueAuthSig} field encoding, with a defined empty case (ZIP-246,
Authorizing Data Commitment).

\begin{remark}[The reference text is itself unfinished]
\label{rem:zsa-zip246-defects}
ZIP-246 is internally incomplete even as withdrawn text: the
authorizing-data section carries ``TODO: Update the Authorizing Data
Commitment below for the \texttt{sighashInfo} changes to the tx format,''
Rationale and Reference implementation are TBD, and the subtree labels are
inconsistent --- \texttt{issue\_actions\_digest} is listed as T.5c but headed
T.5a with children T.5a.i, and the burn-digest children are labelled
T.4b.i/T.4b.ii under node T.4a.vi (\texttt{zip-0246.rst}). Anyone
respecifying the format inherits these defects as open editorial work; the
implementation, not the ZIP, is currently the more precise artifact.
\end{remark}

\subsection{Transaction-level consensus rules}\label{sec:zsa-tx-rules}

\paragraph{One Action Group.} Although the format is a vector of Action
Groups, ZIP-227 collapses it for OrchardZSA: per transaction,
\texttt{nActionGroupsOrchard} MUST be $0$ or $1$, and
\texttt{nAGExpiryHeight} MUST be $0$ (ZIP-227, Specification: Consensus Rule
Changes). The vector form is forward provisioning for ZSA swaps (ZIP-228),
nothing more.

\paragraph{Issuance requires an Action Group.} A transaction with an
issuance bundle MUST also contain at least one OrchardZSA Action Group
(ZIP-227, Specification: Consensus Rule Changes). The rule is load-bearing
twice over. First, issue-note $\rho$ derivation consumes the first Action's
nullifier: each issued note's $\rho$ is computed by the function ZIP-227
names $\mathsf{DeriveIssuedRho}$ --- constructed in
\S\ref{sec:zsa-issuance-notes} --- keyed by $\nfsym_{0,0}$, the nullifier
of the input note of the first Action in the first Action Group, over the
Issue Action and note indices (ZIP-227, Computation of $\rho$), which
gives every issued note a globally unique $\rho$ because nullifiers are
unique.
Second, without a spend in the transaction the issuance-only SIGHASH would
be replayable; the mandatory Action Group ties the issuance signature to a
nullifier that consensus consumes (ZIP-227, Specification: Consensus Rule
Changes).

\paragraph{Global issuance state.} Nodes maintain a map
$\mathsf{issued\_assets}$ from asset bases to triples
$(\mathsf{balance}, \mathsf{final}, \mathsf{note\_ref})$, chained empty map
$\rightarrow$ per-block $\rightarrow$ per-transaction. Burns decrement
$\mathsf{balance}$, requiring $\mathsf{balance} \geq v$, \emph{before} the
transaction's issuance bundle is processed; issuance requires
$\mathsf{final} \neq 1$, increments $\mathsf{balance}$ subject to the
overflow bound $\mathsf{MAX\_ISSUE} = 2^{64} - 1$, and sets
$\mathsf{final}$ when \texttt{finalize} $= 1$. The first issuance of an
asset MUST include a \emph{reference note}: value $0$, paid to the default
address of the all-zero Orchard spending key (ZIP-227, Global Issuance State
and Reference Notes). The 43-byte raw encoding of that address is fixed in
the ZIP and reproduced as a constant in the implementation
(\S\ref{sec:zsa-impl}).

\paragraph{Burn rules and the extended balance equation.} For every
$(\mathsf{AssetBase}, v)$ in \texttt{assetBurn}: $\mathsf{AssetBase} \neq
V^{\mathsf{Orchard}}$ (ZEC cannot be burned this way), $0 < v \leq
\mathsf{MAX\_BURN\_VALUE}$, and no two entries share an
$\mathsf{AssetBase}$; the bound $\mathsf{MAX\_BURN\_VALUE} := 2^{63} - 1$ is
chosen for compatibility with signed 64-bit value-balance fields (ZIP-226,
Burn Mechanism). Burns enter value balance through the binding key. Where
vanilla Orchard computes $\mathsf{bvk}$ by subtracting the committed
declared balance from the sum of net value commitments (the \emph{Orchard
Guide}, \S``Conservation of value: value commitments and the binding
signature''), ZIP-226 subtracts one zero-randomness commitment per burn
entry as well:
\begin{align*}
\mathsf{bvk} :=\ &\Bigl(\sum_{i=1}^{n} \cvsym^{\mathsf{net},i}\Bigr)
- \mathsf{ValueCommit}^{\mathsf{OrchardZSA}}_{0}\bigl(V^{\mathsf{Orchard}},
v_{\mathrm{balance}}^{\mathsf{Orchard}}\bigr) \\
&- \sum_{(\mathsf{AssetBase},\, v) \,\in\, \mathtt{assetBurn}}
\mathsf{ValueCommit}^{\mathsf{OrchardZSA}}_{0}(\mathsf{AssetBase}, v),
\end{align*}
and the prover signs the SIGHASH with $\mathsf{bsk} = \sum_{i=1}^{n}
\rcv_i$ (ZIP-226, Burn Mechanism). The signature verifies only if every
non-ZEC asset nets to exactly its declared burns --- the same
Pedersen-homomorphism argument as for ZEC, run per asset base.

\paragraph{Commitment-tree order.} Note commitments enter the global tree
in a fixed order: first, per Action Group, the commitments of the new
OrchardZSA notes of each Action; then, per Issue Action, the issue-note
commitments (ZIP-227, Specification: Consensus Rule Changes). Wallets and
nodes must agree on this order to agree on note positions.

\subsection{Fees (ZIP-317): specified in the fork, removed upstream}
\label{sec:zsa-fees}

The fee treatment chains through three documents: ZIP-226 (Transaction Fees)
defers to ZIP-317 ``as described in ZIP 227''; ZIP-227 (Changes to ZIP 317)
states that the conventional fee is altered for the presence of issuance
actions and the creation of new Custom Assets, pointing at ZIP-317's Fee
calculation section. Fees remain payable in ZEC only (ZIP-227, rationale).
All variants share the conventional-fee form documented in the
\emph{Wallet Guide}, \S``Fees (ZIP-317)'':
\[
\mathsf{conventionalFee} := \mathsf{marginalFee} \cdot
\max(\mathsf{graceActions},\ \mathsf{logicalActions}),
\]
with $\mathsf{marginalFee} = 5000$~zatoshis and $\mathsf{graceActions} = 2$
(ZIP-317, Fee calculation).

The pointer now dangles upstream: ZIP-317 (Last-Updated 2026-06-26) removed
the ZSA contributions --- its Change History records ``Removed the ZSA
issuance and asset-creation fee contributions (ZIP 227)'' --- leaving
Revision~0 (active base), Revision~1 (Ironwood-pool Actions, NU6.3), and
Revision~2 (memo bundles, draft) (ZIP-317, Revisions and Change History).
The ZSA treatment survives in the QED-it fork copy
(@~\texttt{fd71419}): a requirement that ``users should be discouraged from
issuing new `garbage' Custom Assets,'' the parameter
$\mathsf{creationCost} = 100$ logical actions, and two contributions summed
into $\mathsf{logicalActions}$ alongside the base terms of the deployed
rule:
\begin{align*}
\mathsf{contribution}_{\mathsf{ZSAIssuance}} &:= \mathsf{nZSAIssueNotes},\\
\mathsf{contribution}_{\mathsf{ZSACreation}} &:= \mathsf{creationCost}
\cdot \mathsf{nAssetCreations},
\end{align*}
where $\mathsf{nZSAIssueNotes}$ counts issue-note outputs and
$\mathsf{nAssetCreations}$ counts Custom Assets newly added to the global
issuance state (fork ZIP-317, Fee calculation). The rationale is state rent
paid once: each new asset adds an $\mathsf{issued\_assets}$ row tracked in
perpetuity, while subsequent issuance, finalization, and burn only mutate
the row. The implementation matches the fork text:
\texttt{zcash\_primitives}' \texttt{fees/zip317.rs} defines
\texttt{CREATION\_COST} $= 100$ behind
\texttt{cfg(zcash\_unstable = "nu7")} and adds
$\mathsf{nZSAIssueNotes} + 100 \cdot \mathsf{nAssetCreations}$ to the
standard logical-action count. As with the deployed rule (the \emph{Wallet
Guide}, \S``The deployed fee rule''), paying the conventional fee is not a
consensus requirement in any ZIP-317 variant. Classification: the ZSA fee
rule is \emph{fork-specified}; whether it returns upstream as a ZIP-317
revision, as Ironwood's did, is undecided.

\subsection{Implementation status at the pinned commits}\label{sec:zsa-impl}

\paragraph{Constants and serialization.} All NU7 code in the librustzcash
fork sits behind the configuration gate
\texttt{cfg(zcash\_unstable = "nu7")}. The constants are
$\texttt{V6\_TX\_VERSION} = 6$, $\texttt{V6\_VERSION\_GROUP\_ID} =
\mathtt{0x77777777}$, $\texttt{BranchId::Nu7} = \mathtt{0x77190AD8}$, with
NU7 activation \texttt{None} on Mainnet, $4{,}000{,}000$ on Testnet, $1$ on
Regtest (\texttt{zcash\_protocol}, \texttt{constants.rs} and
\texttt{consensus.rs}). The v6 reader parses the header fragment, one
\texttt{vSighashInfo} per transparent input (required to equal the version-0
encoding), the Sapling v6 bundle, the Orchard v6 bundle, and the issuance
bundle (\texttt{zcash\_primitives}, \texttt{transaction/mod.rs}). The
Orchard reader rejects more than one Action Group (``A V6 transaction must
contain at most one action group''), requires at least one Action, and
rejects a nonzero \texttt{nAGExpiryHeight}; the writer emits exactly one
group with \texttt{nAGExpiryHeight} $= 0$
(\texttt{transaction/components/orchard.rs}) --- the ZIP-227 consensus rules
hard-coded at the serialization layer. Issuance-bundle serialization matches
ZIP-230, including the two-zero-\texttt{compactSize} empty bundle and
rejection of the inconsistent issuer/action combinations
(\texttt{transaction/components/issuance.rs}). The sighash implementation
extends the v5 hash with the issuance digest for ZSA-capable versions, with
sighash-info constants $\texttt{ORCHARD\_SIGHASH\_INFO\_V0} =
\texttt{ISSUE\_SIGHASH\_INFO\_V0} = [\mathtt{0x00}]$
(\texttt{sighash\_v6.rs}, \texttt{sighash\_versioning.rs}). The orchard fork
defines the ZIP-246 personalizations of
Table~\ref{tab:zsa-personalizations} (\texttt{bundle/commitments.rs},
\texttt{bundle/commitments/issuance.rs}), asserting the $\mathtt{0x00}$
BIP-340 byte when hashing the 65-byte issuance signature encoding.

\paragraph{Verification code.} Function \texttt{derive\_bvk\_raw} implements
the extended balance equation of \S\ref{sec:zsa-tx-rules} verbatim
(\texttt{orchard} fork, \texttt{bundle.rs}); \texttt{validate\_bundle\_burn}
enforces the ZIP-226 burn rules --- non-ZEC asset, non-zero amount, amount
$\leq 2^{63} - 1$, no duplicates --- plus presence in and sufficient supply
under the global issuance state (\texttt{bundle/burn\_validation.rs}); and
\texttt{verify\_issue\_bundle} checks the BIP-340 signature over the 32-byte
sighash, the per-note $\rho$ against
$\mathsf{DeriveIssuedRho}$, non-finalization, \texttt{u64}-checked supply
addition, and the reference note on first issuance, whose 43-byte
\texttt{RAW\_REFERENCE\_RECIPIENT} constant equals the array in ZIP-227
(\texttt{issuance.rs}, \texttt{constants/reference\_keys.rs}). The
$\mathtt{0x84}$ domain separator for issued-note $\rho$ lives in the QED-it
\texttt{zcash\_spec} fork (rev \texttt{d5e8426}). Proof sizes record the
circuit cost: the ZSA aggregated proof has base size $2848$ bytes and
$2272$ bytes per Action, against a vanilla base of $2720$
(\texttt{primitives/orchard\_primitives\_zsa.rs},
\texttt{orchard\_primitives\_vanilla.rs}).

\paragraph{Divergences from the withdrawn ZIP text.} Four are material;
each is an explicit interim measure or a pending reassignment, not a bug,
but each blocks byte-compatibility with ZIP-230/246 as written.
\begin{enumerate}[itemsep=2pt]
\item \emph{No explicit fee field.} The v6 reader serializes no fee, and
  \texttt{zip233Amount} only behind the \texttt{zip-233} cargo feature;
  the header txid digest correspondingly omits ZIP-246's T.1f fee field,
  with the comment ``TODO: Factor this out \dots\ when implementing ZIP 246
  in full'' (\texttt{transaction/mod.rs}, \texttt{txid.rs}).
\item \emph{No memo bundles.} The forks implement the pre-ZIP-231 layout:
  OrchardZSA notes carry a full $512$-byte memo, giving a $612$-byte
  \texttt{encCiphertext} with an $84$-byte compact prefix ($52$ vanilla
  bytes plus the $32$-byte asset base), and Sapling v6 outputs are read in
  the v5 $580$-byte format
  (\texttt{primitives/zcash\_note\_encryption\_domain.rs},
  \texttt{transaction/components/sapling.rs}). No memo fields and no T.6
  are implemented; instead the action-group txid digest inserts a
  non-spec memos digest (personalization \texttt{ZTxIdOrcActMHash})
  between the compact and non-compact digests and hashes only the
  $84$-byte compact slice into the compact digest, with TODOs ``Remove
  once new Memo Bundles are implemented (ZIP-231)''
  (\texttt{primitives/orchard\_primitives\_zsa.rs}).
\item \emph{Lead byte.} The implementation hard-codes
  $\texttt{NOTE\_VERSION\_BYTE\_V3} = \mathtt{0x03}$; upstream ZIP-230
  withdrew that claim after ZIP-2005 took $\mathtt{0x03}$, leaving the ZSA
  lead byte an unassigned placeholder. Reassignment is a hard prerequisite
  for any deployment: the lead byte selects the note-plaintext parser.
\item \emph{Split-note $\psi^{\mathsf{nf}}$.} Upstream ZIP-226 now derives
  the split-note nullifier randomizer as
  $\psi^{\mathsf{nf}} = \mathrm{ToBase}\bigl(
  \mathsf{PRFexpand}_{\mathsf{rseed}_{\mathsf{nf}}}(\mathtt{0x0A} \,\|\,
  \rho^{\mathrm{old}})\bigr)$ with $\mathsf{rseed}_{\mathsf{nf}}$ sampled
  uniformly at random, the first-byte domain $\mathtt{0x0A}$ ``reserved by
  ZIP 2005 for this use'' (ZIP-226, split-note nullifier); the fork at
  \texttt{cf801a5d} realizes the older fork text (uniformly random
  $\psi^{\mathsf{nf}}$) via a separate \texttt{rseed\_split\_note} under
  the ordinary $\psi$ domain $\mathtt{0x09}$ (\texttt{note.rs},
  \texttt{circuit.rs}).
\end{enumerate}

\paragraph{Audit scope.} Public reporting supports both standing claims,
at scope: Least Authority's \emph{OrchardZSA Protocol Security Audit
Report} (ZCG-commissioned, Updated Final 30 January 2025) reviewed the
QED-it \texttt{orchard} \texttt{zsa1}-vs-\texttt{main} diff (PR
QED-it/orchard\#7, the ZSA circuit extensions) and the QED-it
\texttt{halo2} changes (PR QED-it/halo2\#18) at revisions
\texttt{a7c02d22}/\texttt{01e85a5f} and identified no issues; ZecHub
\emph{Shielded News} Vol.\,61 (2025-12-17) reports a feature-complete ZSA
build on QED-it's public Zebra testnet; both sources are pinned in
\S\ref{sec:zsa-status}. Neither covers \texttt{librustzcash-zsa} or the
v6 carrier layer, and the audited revisions predate \texttt{cf801a5d}.
Two facts bound what the past audit covers: the audited revisions predate
both the NU6.2 Orchard circuit remediation (ZIP-257) --- and the ZSA
circuit is a fork of the Orchard circuit --- and the ZIP-2005-relative
rebasing of upstream ZIP-226. A mainnet path implies re-audit against
both.

\subsection{Wallet integration}\label{sec:zsa-wallets}

ZIP-230 and ZIP-226 impose two obligations. First, all wallets SHOULD
switch to sending only v6 transactions once permitted: a v5 transaction
reveals that all its Orchard notes are ZEC, and only v6 OrchardZSA outputs
are quantum-recoverable (ZIP-230, Implications for Wallets; ZIP-226,
Backwards
Compatibility with ZEC Notes). Second, wallets MUST support parsing and
trial-decrypting v6 outputs --- the new lead byte and the memo-bundle
changes --- by activation, or rescan once support lands: funds sent to an
un-upgraded wallet are temporarily inaccessible, because the new lead byte
selects different $\rcm$/$\psi$ derivations and a v5-decryption wallet
rejects the note (ZIP-230, Implications for Wallets). Payment requests have
a
specified extension path: ZIP-321's \texttt{req-asset} parameter carries an
Asset Identifier and amount (the \emph{Wallet Guide}, \S``Custom assets:
\texttt{req-asset} (specified, not deployed)'').

Multi-party and hardware signing is the open end. No PCZT structure exists
for OrchardZSA bundles or issuance bundles: the librustzcash fork's
transaction extractor panics with ``PCZT support for ZSA is not
implemented'' when handed an OrchardZSA bundle, the orchard fork's PCZT
module is written against the vanilla flavor only, and the draft PCZT v2
targets ZIP-229's Ironwood v6, not ZSA (the \emph{Wallet Guide},
\S``Partially created transactions (PCZT)'' and \S``PCZT v2: deferred
anchors and the Ironwood bundle (specified, not deployed)'').
Classification: \emph{open} --- hardware-wallet and multi-party signing for
ZSAs is unspecified in any draft.

\subsection{Outlook}\label{sec:zsa-outlook}

ZIP-226 and ZIP-227 remain Status: Draft. ZIP-226's Deployment section reads
``The Zcash Shielded Assets protocol is scheduled to be deployed in Network
Upgrade 7 (NU7). This ZIP currently assumes that this will be the case'';
ZIP-227's Deployment is TBD and its test-vector and reference-implementation
entries are ``LINK TBD,'' with ZIP-226 naming the QED-it forks and
\texttt{QED-it/zcash-test-vectors} as reference artifacts (ZIP-226/227,
Deployment). The deployment vehicle is thinner still: the NU7 deployment
draft is unnumbered (\texttt{draft-arya-deploy-nu7}, Status: Draft,
replacing the withdrawn ZIP-254), fixes
$\texttt{CONSENSUS\_BRANCH\_ID} = \mathtt{0x77190AD8}$ and minimum protocol
versions ($170150$ Testnet, $170160$ Mainnet) with activation heights TBD,
and lists \emph{no} ZIPs. The repository's NU7 candidate list (ZIPs 218,
230, 231, 233, 234, 235, 2002, 2003, plus a possible ZIP-317 update) does
not include ZIP-226/227, still names the withdrawn ZIP-230, and states that
``no decision has been made on whether to include each of these ZIPs''
(\texttt{zcash/zips} README, NU7 Candidate ZIPs). Governance sentiment is
split: in the February 2026 NU7 sentiment polls, ZCAP supported ZSAs at
70.4\% while the coinholder poll ran 98.6\% against; the Zcash Foundation
did not recommend ZSAs for NU7, writing that ``before any decision is
made, we need to understand why coinholders oppose ZSAs so strongly''
(zfnd.org, ``NU7 Polling Results: What We Heard and Where We Go From
Here'', 2026-02-23,
\url{https://zfnd.org/nu7-polling-results-what-we-heard-and-where-we-go-from-here/}).

Between the pinned implementation and any deployment stand three
respecification problems, none of which is per-object protocol work.
\begin{enumerate}[itemsep=2pt]
\item \emph{A carrier.} Transaction version 6 now belongs to NU6.3's
  Ironwood format (ZIP-229), whose non-requirements exclude the ZSA
  machinery (\S\ref{sec:zsa-format-status}). The designated successor for
  what ZIP-230/246 deferred --- an extensible transaction format --- is
  ZIP-248, which exists only as open PR \texttt{zcash/zips\#1156}. The
  entire ZSA and memo-bundle format stack therefore has no merged
  specification; if ZSAs are selected for NU7, the byte format of
  \S\ref{sec:zsa-v6} must be respecified under new numbers, and
  \S\ref{sec:zsa-digests} re-anchored to whatever digest ZIP results.
\item \emph{A base protocol.} Upstream ZIP-226 is now ``defined relative to
  the Zcash protocol with the changes specified in ZIP 2005 applied,'' with
  ZIP-2005 (Ironwood Quantum Recoverability) ``expected to deploy before any
  ZSA activation'' (ZIP-226, header notes). NU6.3 (ZIP-258, testnet
  activation $4{,}134{,}000$, Mainnet TBD) introduces the Ironwood pool and
  restricts transfers into the Orchard pool (ZIP-2006, currently a Reserved
  stub). No ZIP specifies whether OrchardZSA then targets the restricted
  Orchard pool or is rebased onto Ironwood; the audited revisions
  predate the entire NU6.2/NU6.3 pivot. Classification: \emph{open}.
\item \emph{Reconvergence of spec and code.} The divergence list of
  \S\ref{sec:zsa-impl} --- fee field, memo bundles, lead byte,
  $\psi^{\mathsf{nf}}$ --- plus the orphaned fee treatment of
  \S\ref{sec:zsa-fees} must land on one side or the other: either the
  respecified ZIPs adopt what the code does, or the code moves. Every item
  is known and mechanical; none is done.
\end{enumerate}

The summary judgment this volume can defend: the ZSA transaction layer is
\emph{designed and implemented, but currently unspecified} in
the sense this series requires --- there is no merged normative text a
mainnet node could be held to. The byte-level facts of this section are
fixed by withdrawn ZIP text and pinned fork code, and they will remain the
diff base for the respecification; the open items are enumerated above and
in the scope section, and every one of them is legible, which is more than
can be said for most protocols at this stage.

\end{document}
